
%% ============================================================
%%  main.tex — الملف الرئيسي للتقرير
%%  تقرير حول الشغل المؤقت والمناولة في مدونة الشغل المغربية
%% ============================================================

\documentclass[12pt, a4paper]{report}

%% ─── الحزم الأساسية ───────────────────────────────────────
\usepackage{fontspec}           % دعم الخطوط OpenType (XeLaTeX/LuaLaTeX)
\usepackage{polyglossia}        % دعم تعدد اللغات بما فيها العربية
\usepackage{geometry}           % هوامش الصفحة
\usepackage{graphicx}           % إدراج الصور
\usepackage{xcolor}             % الألوان
\usepackage{titlesec}           % تخصيص عناوين الفصول
\usepackage{titletoc}           % تخصيص فهرس المحتويات
\usepackage{fancyhdr}           % ترويسات وتذييلات مخصصة
\usepackage{hyperref}           % روابط تشعبية داخلية
\usepackage{microtype}          % جودة الطباعة
\usepackage{setspace}           % تباعد الأسطر
\usepackage{enumitem}           % تخصيص القوائم
\usepackage{booktabs}           % جداول احترافية
\usepackage{array}
\usepackage{multirow}
\usepackage{colortbl}           % تلوين خلايا الجداول (\rowcolor)
\usepackage{amsmath}
\usepackage{amssymb}
\usepackage{tcolorbox}          % صناديق ملونة
\usepackage{emptypage}          % صفحات فارغة بدون ترويسة

%% إصلاح توافق bidi مع جداول LaTeX (حل خطأ UseMathForPositioningText في bidi 2026)
\providecommand{\UseMathForPositioningText}{}
\providecommand{\UseTaggingSocket}[1]{}

%% ─── إعداد اللغة العربية ──────────────────────────────────
\setdefaultlanguage[numerals=maghrib]{arabic}
\setotherlanguage{french}
\setotherlanguage{english}

\newfontfamily\arabicfont[
    Script=Arabic,
    Scale=1.1,
    Ligatures=TeX
]{Amiri}                        % خط أميري (يمكن تغييره إلى Scheherazade New)

%% خط أحادي المسافة
\setmonofont{Courier New}

%% ─── هوامش الصفحة ────────────────────────────────────────
\geometry{
    top=2.5cm,
    bottom=2.5cm,
    left=2.5cm,
    right=3cm,
    headheight=30pt
}

%% ─── الألوان المخصصة ──────────────────────────────────────
\definecolor{maincolor}{RGB}{0, 51, 102}       % أزرق داكن رسمي
\definecolor{accentcolor}{RGB}{180, 0, 0}      % أحمر داكن
\definecolor{lightgray}{RGB}{245, 245, 245}
\definecolor{sectioncolor}{RGB}{0, 80, 140}

%% ─── تنسيق العناوين ──────────────────────────────────────
\titleformat{\chapter}[display]
  {\normalfont\huge\bfseries\color{maincolor}}
  {\chaptertitlename\ \thechapter}{20pt}{\Huge}

\titleformat{\section}
  {\normalfont\Large\bfseries\color{sectioncolor}}
  {\thesection}{1em}{}

\titleformat{\subsection}
  {\normalfont\large\bfseries\color{accentcolor}}
  {\thesubsection}{1em}{}

%% ─── ترويسات وتذييلات ─────────────────────────────────────
\pagestyle{fancy}
\fancyhf{}
%% نستخدم \protect\RLE لتفادي تحذيرات endL/endR من bidi+fancyhdr
\fancyhead[R]{\small\color{maincolor}\protect\RLE{\rightmark}}
\fancyhead[L]{\small\color{maincolor}\protect\RLE{الشغل المؤقت والمناولة}}
\fancyfoot[C]{\small\thepage}
\renewcommand{\headrulewidth}{0.5pt}
\renewcommand{\footrulewidth}{0.3pt}

%% ─── تباعد الأسطر ────────────────────────────────────────
\setstretch{1.5}

%% ─── إعداد الروابط ───────────────────────────────────────
\hypersetup{
    colorlinks=true,
    linkcolor=maincolor,
    urlcolor=accentcolor,
    citecolor=sectioncolor,
    pdfauthor={فريق الدراسة - المجموعة 8},
    pdftitle={الشغل المؤقت والمناولة في مدونة الشغل المغربية},
    pdfsubject={دراسة قانونية وتطبيقية},
    pdflang={ar}
}

%% ============================================================
%%  بداية الوثيقة
%% ============================================================
\begin{document}

%% ─── صفحة الغلاف ─────────────────────────────────────────
%% ============================================================
%%  cover.tex — صفحة الغلاف الاحترافية
%% ============================================================
\begin{titlepage}
\thispagestyle{empty}
\pagecolor{maincolor}
\color{white}

\vspace*{0.8cm}
{\centering

%% ─── الجهة المؤسسية ───────────────────────────────────────
{\large\bfseries المملكة المغربية}\\[0.3cm]
{\large وزارة التعليم العالي والبحث العلمي والابتكار}\\[0.2cm]
{\normalsize وحدة البحث والتكوين في القانون الاجتماعي والعلاقات المهنية}\\[0.4cm]
\rule{0.9\textwidth}{0.6pt}\\[1cm]

%% ─── صندوق العنوان الرئيسي ────────────────────────────────
\setlength{\fboxsep}{18pt}
\fcolorbox{white}{white}{\color{maincolor}%
\parbox{0.82\textwidth}{\centering
    {\LARGE\bfseries الشغل المؤقت والمناولة}\\[0.5cm]
    {\large\bfseries في ضوء مدونة الشغل المغربية (القانون رقم 65.99)}\\[0.4cm]
    \rule{0.55\textwidth}{0.8pt}\\[0.4cm]
    {\normalsize دراسة قانونية واجتماعية وتطبيقية — دراسة الحالة: المجموعة \textenglish{8}}
}}

\vspace{1.5cm}

%% ─── بيانات الإعداد ──────────────────────────────────────
{\large\bfseries إعداد الطلبة (المجموعة الثامنة --- \textenglish{Groupe 8}):}\\[0.4cm]
{\large \textenglish{Haron Jadid}}\\[0.15cm]
{\large \textenglish{Said Nichan}}\\[0.15cm]
{\large \textenglish{Khouka Hajar}}\\[0.15cm]
{\large \textenglish{Jaouad El Hail}}\\[0.15cm]
{\large \textenglish{Nour Elhouda Ouhba}}\\[0.7cm]

\rule{0.9\textwidth}{0.5pt}\\[0.5cm]

%% ─── المشرف ──────────────────────────────────────────────
{\large\bfseries تحت إشراف:}\\[0.3cm]
{\large الأستاذ(ة) المشرف(ة) على وحدة قانون الشغل}\\[0.8cm]

\rule{0.9\textwidth}{0.5pt}\\[0.5cm]

%% ─── الوحدة والسنة الدراسية ──────────────────────────────
{\normalsize\bfseries الوحدة الدراسية:}
{\normalsize قانون الشغل والعلاقات المهنية}\\[0.3cm]
{\normalsize\bfseries السنة الدراسية:}
{\normalsize \textenglish{2025 -- 2026}}

\vfill

\rule{0.9\textwidth}{0.5pt}\\[0.3cm]
{\small \textenglish{2025-2026} \quad|\quad جميع الحقوق محفوظة}

} %% end \centering
\end{titlepage}

%% إعادة اللون الأبيض للخلفية بعد الغلاف
\pagecolor{white}
\color{black}


%% ─── صفحة فارغة بعد الغلاف ───────────────────────────────
\cleardoublepage

%% ─── الفهرس ──────────────────────────────────────────────
\tableofcontents
\cleardoublepage

%% ─── المقدمة العامة ───────────────────────────────────────
% =====================================================================
% المقدمة العامة وإشكالية الدراسة
% =====================================================================

\chapter*{مقدمة عامة وإشكالية الدراسة}
\addcontentsline{toc}{chapter}{مقدمة عامة وإشكالية الدراسة}
\markboth{مقدمة عامة وإشكالية الدراسة}{مقدمة عامة وإشكالية الدراسة}

%% ---------------------------------------------------------------------
\section*{تمهيد وتأصيل تاريخي واقتصادي}
%% ---------------------------------------------------------------------

شهدت المنظومة الاقتصادية والإنتاجية العالمية خلال العقود الأخيرة تحولات جذرية فرضتها عولمة الأسواق، واحتدام المنافسة، وثورة التكنولوجيا وسرعة تدفق رؤوس الأموال. وقد ألقت هذه التحولات بظلالها المباشرة على العلاقات المهنية وأنماط إدارة الموارد البشرية، حيث تراجعت المقاولات الحديثة تدريجياً عن النمط التقليدي للتشغيل القائم على الاستقرار الوظيفي وعقد الشغل الفردي غير المحدد المدة (CDI) كقاعدة عامة وحيدة، لصالح البحث المحموم عما بات يعرف في الأدبيات الاقتصادية والقانونية بـ \textbf{«مرونة الشغل»} (\textit{Flexibilité de l'emploi}).

وتتجلى هذه المرونة في قدرة المنشأة على ملاءمة حجم وقدرات مواردها البشرية مع تقلبات السوق ودورات الإنتاج وفترات الذروة المفاجئة دون تحمل الأعباء المالية والتنظيمية الثقيلة المترتبة على إنهاء عقود الشغل القارة. وفي هذا السياق، برز نمطان أساسيان ومترابطان لإعادة هيكلة قوة العمل:
\begin{enumerate}
    \item \textbf{الشغل المؤقت (\textit{Travail temporaire / Intérim}):} وهو علاقة ثلاثية الأبعاد تقوم بمقتضاها مقاولة متخصصة مرخص لها بوضع أجراء ترتبط معهم بروابط تعاقدية رهن إشارة منشأة أخرى مستعملة، لأداء مهام غير دائمة وبشروط محددة حصراً.
    \item \textbf{المناولة أو المقاولة من الباطن (\textit{Sous-traitance}):} وهي تقنية تعاقدية تجارية يعهد بموجبها مقاول أصلي إلى مقاول آخر مستقل بإنجاز عمل أو تقديم خدمة متخصصة، محتفظاً بالتبعية المهنية والفنية لأجرائه.
\end{enumerate}

وعلى صعيد النسيج الاقتصادي المغربي، تسارع اللجوء إلى هاتين الآليتين منذ أواخر تسعينيات القرن الماضي مع انفتاح الاقتصاد الوطني وتوقيع اتفاقيات التبادل الحر وتنامي قطاعات استراتيجية كالسيارات والطيران واللوجستيك والصناعات الغذائية. ورغم أن هذا التوجه وفّر متنفساً حقيقياً للمقاولات لتخفيض تكاليف التسيير وضمان التنافسية، وفتح أمام فئات واسعة من الشباب والباحثين عن شغل بوابة للاندماج المهني واكتساب المهارات، إلا أنه أفرز في المقابل إشكالات بنيوية بالغة الخطورة؛ حيث أدى تفكيك الرابطة العقدية الثنائية المباشرة بين الأجير ورب العمل الفعلي إلى إضعاف الحماية الاجتماعية، وتشتيت عناصر التبعية القانونية، وتنامي مظاهر الهشاشة المهنية وتدني الأجور.

%% ---------------------------------------------------------------------
\section*{الإطار الدستوري والمرجعية الدولية}
%% ---------------------------------------------------------------------

تتقاطع إشكالية الشغل المؤقت والمناولة مع ثنائية دستورية بالغة الدقة؛ فمن جهة، يكرس \textbf{دستور المملكة المغربية لسنة 2011} في فصله الخامس والثلاثين (35) حرية المبادرة والمقاولة، والتنافس الحر، وضمان مناخ ملائم للاستثمار. ومن جهة ثانية، ينص الفصل الحادي والثلاثون (31) على التزام الدولة والمؤسسات العمومية والجماعات الترابية بتعبئة كل الوسائل المتاحة لتيسير استفادة المواطنات والمواطنين من الحق في الشغل، وضمان شروط عمل لائقة وعادلة، والاستفادة من الحماية الاجتماعية والتغطية الصحية والتكوين المهني.

ودولياً، تكتسي هذه المسألة عناية خاصة لدى \textbf{منظمة العمل الدولية (OIT)}، التي وضعت معايير متقدمة للتوفيق بين مرونة السوق وحماية العمال:
\begin{itemize}
    \item \textbf{الاتفاقية رقم 181 بشأن وكالات الاستخدام الخاصة (1997):} التي كرست الاعتراف بالدور الإيجابي لشركات التشغيل المؤقت في تسيير سوق الشغل، شريطة إخضاعها لترخيص صارم، ومنع اقتطاع أي مصاريف أو أتعاب من أجور العمال، وإقرار مسؤولية مشتركة واضحة تضمن تمتع الأجراء بكافة حقوقهم الأساسية والتفاوض الجماعي.
    \item \textbf{التوصية رقم 198 بشأن علاقة العمل (2006):} التي شددت على المعيار الجوهري المتمثل في \textbf{«مبدأ أولوية الواقع على المظاهر والشكل»} (\textit{Principe de la primauté des faits})؛ حيث تقضي بوجوب تكييف علاقة الشغل بالنظر إلى الوقائع الميدانية وممارسة السلطة الفعلية، بصرف النظر عن التسميات أو العقود الصورية التي يضفيها الأطراف للتنصل من موجبات القانون.
\end{itemize}

%% ---------------------------------------------------------------------
\section*{سياق الدراسة وأهميتها}
%% ---------------------------------------------------------------------

تنطلق هذه الدراسة من واقعة عملية ذات راهنية قصوى تتكرر بصورة شبه نمطية في العديد من المنشآت الصناعية والإنتاجية الكبرى بالمغرب:
\begin{quote}
\textit{«شركة صناعية كبرى تعقد اتفاقاً مع مقاولة متخصصة تحت غطاء "المناولة"، غير أن المحتوى الفعلي للعقد ينحصر في توفير اليد العاملة والعمالة اليدوية. ويشتغل هؤلاء العمال داخل أسوار المصنع ذاته لسنوات متتالية، جنباً إلى جنب مع أجراء المقاولة الأصلية، مؤدين نفس المهام والوظائف المندرجة في النشاط العادي والمستمر للمصنع، وتحت الإشراف والتوجيه المباشر لأطر المقاولة المستفيدة. ورغم ذلك، يتقاضى هؤلاء العمال أجوراً أدنى، ويحرمون من التعويضات والترقيات والامتيازات الاجتماعية الممنوحة لنظرائهم القارين، وسط غموض كامل يكتنف الجهة المسؤولة قانونياً عن سلامتهم وحمايتهم وضمان حقوقهم».}
\end{quote}

وتتفرع أهمية تفكيك هذه النازلة إلى أربعة أبعاد متداخلة:
\begin{itemize}
    \item \textbf{البعد القانوني والتنظيمي:} ضبط الفوارق الماهوية الفاصلة بين الشغل المؤقت، والمقاولة من الباطن، وعقد الشغل المحدد المدة (CDD)، والتمييز بين الإسناد المشروع للخدمات وبين جريمة "إعارة اليد العاملة غير المشروعة بمقابل" والتحايل على أحكام مدونة الشغل (القانون رقم 65.99).
    \item \textbf{البعد الاجتماعي والإنساني:} صيانة كرامة الأجراء ودرء الهشاشة المزدوجة الناتجة عن تدني الأجر وعدم استقرار الرابطة العقدية، وتكريس المبدأ الحقوقي الكوني: «أجر متساوٍ عن عمل متساوي القيمة»، وضمان المساواة التامة في شروط الوقاية والسلامة وحوادث الشغل.
    \item \textbf{البعد الاقتصادي والتنافسي:} تحديد التخوم الفاصلة بين المرونة الاقتصادية المشروعة التي تعزز تنافسية المقاولة وتساعدها على امتصاص الأزمات الظرفية، وبين "الإغراق الاجتماعي" (\textit{Dumping social}) القائم على ضغط التكاليف عبر استغلال الثغرات التشريعية، بما يضر بالتنافسية الشريفة بين المقاولات.
    \item \textbf{البعد العملي والمؤسساتي:} رسم معالم المسؤولية القانونية المدنية والجنائية للأطراف (المشغل الأصلي والوسيط)، وتحديد نطاق التدخل الرقابي المتاح لمفتشية الشغل والقضاء الاجتماعي لإنفاذ مقتضيات النظام العام الاجتماعي وحماية الطرف الضعيف في العلاقة.
\end{itemize}

%% ---------------------------------------------------------------------
\section*{الإشكالية المركزية}
%% ---------------------------------------------------------------------

في ضوء التجاذب الدائم بين ضرورات المرونة التدبيرية للمقاولات من جهة، ومقتضيات الاستقرار المهني والأمن التعاقدي للأجراء من جهة ثانية، تتمحور الإشكالية المركزية لهذه الدراسة حول التساؤل الجوهري الآتي:

\begin{tcolorbox}[colback=lightgray, colframe=maincolor, arc=5pt, title=\textbf{الإشكالية المحورية}, fonttitle=\bfseries]
\centering
\large
إلى أي حد يوفق الإطار القانوني المغربي في تحقيق التوازن الحرج بين تلبية متطلبات المقاولات في المرونة وتقليص التكاليف عبر الشغل المؤقت والمناولة، وبين صيانة الحقوق الجوهرية للأجراء وحمايتهم من مظاهر الاستغلال والتحايل التعاقدي على أرض الواقع؟
\end{tcolorbox}

وتتفرع عن هذه الإشكالية الأم جملة من التساؤلات الفرعية النسقية:
\begin{enumerate}
    \item ما هي الخصائص المفاهيمية المميزة لكل من الشغل المؤقت وعقد المقاولة من الباطن، وما هي المعايير الضابطة للفصل بينهما وبين عقود الشغل محددة المدة (CDD)؟
    \item ما مدى دقة الحالات الحصرية والقيود الزمنية والموضوعية التي رسمها المشرع المغربي في المواد 495 إلى 506 من مدونة الشغل لتأطير الاستعانة بمقاولات التشغيل المؤقت؟
    \item ما هي الالتزامات والضمانات المقررة لصالح الأجير تجاه كل من المشغل القانوني (الوسيط) والمستعمل الفعلي، لا سيما في مجالات الأجر، والتصريح بالضمان الاجتماعي، والتأمين ضد حوادث الشغل والأمراض المهنية؟
    \item ما هي حدود المسؤولية القانونية الملقاة على المقاول الأصلي في عقود المناولة وفقاً للمادتين 89 و90 من المدونة، وما مدى كفاية شروط هذه المسؤولية لضمان حقوق العمال في حالة إعسار المقاول من الباطن؟
    \item هل يملك القضاء الاجتماعي وأجهزة تفتيش الشغل الأدوات والآليات الكفيلة بزجر الممارسات الصورية، وإعادة تكييف عقود المناولة التي تستعمل كغطاء لتوفير يد عاملة دائمة إلى عقود شغل مباشرة وغير محددة المدة؟
\end{enumerate}

%% ---------------------------------------------------------------------
\section*{فرضيات البحث}
%% ---------------------------------------------------------------------

سعياً للإحاطة بمختلف جوانب الموضوع واختبار صحة الطروحات الفقهية والقضائية، تنطلق دراستنا من الفرضيات العلمية الأربع التالية:
\begin{enumerate}
    \item \textbf{الفرضية الأولى (التباين الماهوي في محل الالتزام):} يختلف عقد الشغل المؤقت جوهرياً عن عقد المناولة من حيث موضوع الالتزام؛ فالشغل المؤقت ينصب حصراً على وضع طاقة عمل بشرية رهن إشارة المستعمل مع انتقال جزء من التبعية القانونية، في حين تنصب المناولة على التزام المقاول من الباطن بتحقيق نتيجة مستقلة (شغل أو خدمة) عبر وسائله وإشرافه التقني المباشر دون تفويت التبعية.
    \item \textbf{الفرضية الثانية (عدم مشروعية تغطية الحاجيات الدائمة):} إن استخدام الأجراء المؤقتين أو عمال المناولة لشغل مناصب قارة ولأداء مهام تدخل في صميم النشاط الدائم والعادي للمقاولة المستعملة لسنوات متواصلة، يمثل خرقاً جوهرياً للنظام العام الاجتماعي، ويبرر قانوناً وقضائياً إسقاط الحجاب التعاقدي وإعادة تكييف العلاقة إلى عقد شغل غير محدد المدة (CDI) مباشر مع المقاولة المستعملة.
    \item \textbf{الفرضية الثالثة (قصور نظام المسؤولية الاحتياطية بالمناولة):} يعاني نظام المسؤولية التضامنية أو الاحتياطية المنصوص عليه في المادة 89 من مدونة الشغل من ثغرات هيكلية حادة؛ لكونه يعلق مساءلة المقاول الأصلي على شروط عسيرة في مقدمتها ثبوت إعسار المقاول من الباطن وعدم تقييده في السجل التجاري وعدم توفره على أصل تجاري، ويحصر الأداء في حدود المبالغ المترتبة بذمته، مما يفرغ الحماية الاجتماعية للأجراء من محتواها العملي.
    \item \textbf{الفرضية الرابعة (رهان الفعالية الرقابية):} إن صيانة التوازن الاجتماعي وحماية العمال في هذه الأنماط المرنة تظل رهينة بتعزيز الصلاحيات التقديرية والضبطية لمفتشي الشغل، وتكريس شجاعة القضاء الاجتماعي في تطبيق مبدأ "أولوية الواقع على المظهر التعاقدي"، فضلاً عن مراجعة العقوبات الزجرية الهزيلة المنصوص عليها في المادة 505 من المدونة.
\end{enumerate}

%% ---------------------------------------------------------------------
\section*{المنهجية المعتمدة وتقسيم التقرير}
%% ---------------------------------------------------------------------

ولمعالجة الإشكالية وفحص الفرضيات، تعتمد هذه الدراسة تكاملاً منهجياً متعدد الروافد:
\begin{itemize}
    \item \textbf{المنهج التحليلي التأصيلي:} لتشريح وتفسير النصوص القانونية المؤطرة في مدونة الشغل (القانون 65.99) والقوانين المكملة (قانون الضمان الاجتماعي وقانون التعويض عن حوادث الشغل)، ورصد مكامن قوتها وثغراتها.
    \item \textbf{المنهج المقارن:} لمقارنة المنظومة القانونية المغربية مع نظيرتها في التشريعات المقارنة (خاصة التشريع الفرنسي) ومعايير منظمة العمل الدولية، للوقوف على التطورات الحديثة في تكييف المسؤوليات.
    \item \textbf{منهج دراسة الحالة والمنهج العملي:} لتنزيل القواعد المجردة على واقع النازلة المعروضة، وتشخيص أدوار الفاعلين الحقيقيين (المشغل، الوسيط، الأجراء، مفتش الشغل، القضاء).
\end{itemize}

وبناءً على هذه المقاربة، جرى تقسيم هذا التقرير البحثي إلى عشرة فصول متناسقة ومبنية وفق هندسة موحدة:
\begin{itemize}
    \item \textbf{الفصل الأول:} الإطار المفاهيمي (الشغل المؤقت، المناولة، التمييز بينهما).
    \item \textbf{الفصل الثاني:} الإطار القانوني في مدونة الشغل المغربية.
    \item \textbf{الفصل الثالث:} الحالات التي يسمح فيها القانون باللجوء إلى الشغل المؤقت.
    \item \textbf{الفصل الرابع:} التزامات شركة التشغيل المؤقت والمقاولة المستعملة.
    \item \textbf{الفصل الخامس:} المسؤولية القانونية لكل من الشركة الأصلية وشركة المناولة.
    \item \textbf{الفصل السادس:} دور مفتش الشغل ومراقبة احترام القانون.
    \item \textbf{الفصل السابع:} الآثار الاجتماعية والاقتصادية (الأجراء، المقاولة، سوق الشغل).
    \item \textbf{الفصل الثامن:} دراسة الحالة وتحليل أدوار الفاعلين.
    \item \textbf{الفصل التاسع:} الاجتهاد القضائي والتجارب المقارنة.
    \item \textbf{الفصل العاشر:} الاختلالات والتوصيات.
\end{itemize}

وينتهي التقرير بخاتمة عامة تركيبية تجيب عن الإشكالية المحورية وتلخص الحلول القانونية لأسئلة النازلة، متبوعة بملحقين (نصوص قانونية مختارة من مدونة الشغل ونماذج عقود استرشادية) ولائحة المراجع المعتمدة.
\cleardoublepage

%% ─── الفصل الأول: الإطار المفاهيمي ───────────────────────
% =====================================================================
% الفصل الأول: الإطار المفاهيمي: الشغل المؤقت والمناولة والتمييز بينهما
% =====================================================================

\chapter{الإطار المفاهيمي: الشغل المؤقت والمناولة والتمييز بينهما}
\label{chap:chapitre1}

\section*{مقدمة الفصل}

يعد الضبط المفاهيمي منطلقاً أساسياً وحجر الزاوية في أي دراسة قانونية واجتماعية معمقة؛ إذ إن الخلط بين المفاهيم في حقل العلاقات المهنية لا يقف عند حدود التوصيف النظري، بل ينعكس مباشرة على تحديد القواعد الآمرة الواجبة التطبيق، وتعيين أطراف المسؤولية، ورسم معالم الحماية المقررة للطرف الضعيف في الرابطة العقدية وهو الأجير[cite: 1, 2]. 

ويبرز هذا التداخل جلياً في الممارسة العملية بين تقنيتي \textbf{الشغل المؤقت} (\textit{Travail temporaire}) و\textbf{المناولة أو المقاولة من الباطن} (\textit{Sous-traitance})؛ فهما يشتركان ظاهرياً في إقحام طرف ثالث وسيط بين العامل والمنشأة المستفيدة من جهده العضلي أو الفكري[cite: 2]. غير أن هذا التشابه الشكلي يخفي تبايناً جوهرياً في الفلسفة القانونية ومحل الالتزام وطبيعة التبعية[cite: 1, 2].

يهدف هذا الفصل إلى تأصيل مفهوم الشغل المؤقت وتفكيك هندسته التعاقدية الثلاثية (المبحث الأول)، ثم تحليل مفهوم المقاولة من الباطن وشروط مشروعيتها (المبحث الثاني)، قبل استعراض معايير التمييز النسقية بينهما وبين عقد الشغل المحدد المدة (CDD)، وتنزيل ذلك تحليلياً على واقع النازلة الميدانية المعروضة في سياق هذه الدراسة (المبحث الثالث)[cite: 1, 2, 3].

% ---------------------------------------------------------------------
\section{مفهوم الشغل المؤقت وهندسته التعاقدية}

\subsection{التعريف القانوني لمقاولة التشغيل المؤقت}
عرف المشرع المغربي مقاولة التشغيل المؤقت في المادة 495 من مدونة الشغل بأنها: \textit{«كل شخص اعتباري، مستقل عن السلطة العمومية، يقتصر عمله على ممارسة النشاط المشار إليه في الفقرة (ج) من المادة 477، والمتمثل في تشغيل أجراء بهدف وضعهم، مؤقتاً، رهن إشارة شخص ثالث يسمى "المستعمل"، يحدد مهامهم ويراقب تنفيذها»}[cite: 2].

وانطلاقاً من هذا النص الصريح، يتأسس مفهوم الشغل المؤقت على أركان قانونية آمرة تهدف إلى ضبط هذه المهنة ومنع انحرافها:
\begin{enumerate}
    \item \textbf{الشخصية الاعتبارية المسبقة:} حصر المشرع مزاولة هذا النشاط في الأشخاص الاعتباريين (شركات تجارية)[cite: 2, 3]، مانعاً الأشخاص الذاتيين (الطبيعيين) من ممارسته[cite: 3]. والغاية من هذا الشرط هي ضمان ملاءة مالية ومسك محاسبة منتظمة تتيح مراقبة التدفقات المالية وقدرة الشركة على الوفاء بأجور الأجراء واشتراكاتهم الاجتماعية[cite: 2, 3].
    \item \textbf{مبدأ حصرية النشاط (\textit{Principe d'exclusivité}):} أوجبت المادة 495 أن «يقتصر عملها» على نشاط وضع الأجراء[cite: 2]. فلا يجوز قانوناً لشركة التشغيل المؤقت أن تجمع بين هذا النشاط وأنشطة تجارية أخرى، كالمقاولات العامة أو الخدمات أو التوريدات؛ وذلك لقطع الطريق أمام استغلال الأجراء في أنشطة موازية أو خلط الذمم المالية والمصالح التجارية[cite: 2].
    \item \textbf{الطابع الاستثنائي والعارض للعمل («المهمة»):} لا يمارس الأجير المؤقت عملاً دائماً ومستمراً، بل يعين لإنجاز «مهمة» (\textit{Mission}) محددة الأجل والهدف، ومحصورة في الحالات التي عددها المشرع استثناءً في المادة 496 من المدونة[cite: 2, 3].
\end{enumerate}

\subsection{العلاقة الثلاثية وتجزئة التبعية القانونية}
تتميز علاقة الشغل المؤقت بتركيبتها الثلاثية غير النمطية، إذ تنفصم فيها الثنائية الكلاسيكية لعقد الشغل (أجير - مشغل) لتتحول إلى شبكة تعاقدية تضم ثلاثة فاعلين أساسيين تحكمهم رابطتان عقديتان[cite: 2]:

\begin{itemize}
    \item \textbf{الرابطة العقدية الأولى (عقد الشغل المؤقت):} وهو عقد يربط بين مقاولة التشغيل المؤقت والأجير[cite: 2]. وبموجبه تعد المقاولة هي \textbf{المشغل القانوني الحصري}؛ فهي التي تتحمل عبء أداء الأجر، والتعويض عن العطل المؤدى عنها، والتصريح الإجباري بالأجير لدى الصندوق الوطني للضمان الاجتماعي، وممارسة السلطة التأديبية والجزائية كالفصل أو الإنذار[cite: 2, 3].
    \item \textbf{الرابطة العقدية الثانية (عقد وضع الأجير رهن الإشارة):} وهو عقد تجاري مدني يربط مقاولة التشغيل المؤقت بالمقاولة المستعملة[cite: 2, 3]. وقد اشترطت المادة 499 كتابته وبيان سبب اللجوء، ومدة المهمة، ومكانها، والمقابل المالي المتفق عليه لتوفير اليد العاملة[cite: 2, 3].
\end{itemize}

وينجم عن هذه الهندسة ما يسميه الفقه الاجتماعي بـ \textbf{«تجزئة التبعية القانونية»} (\textit{Démembrement du lien de subordination})؛ حيث تنشطر عناصر التبعية إلى شقين:
\begin{itemize}
    \item \textbf{تبعية قانونية وإدارية:} تظل لصيقة بمقاولة التشغيل المؤقت (إبرام العقد، الأجر، التأديب، إنهاء العلاقة التعاقدية)[cite: 2].
    \item \textbf{تبعية تقنية ووظيفية وميدانية:} تنتقل فعلياً إلى المقاولة المستعملة؛ إذ يمارس مسؤولوها سلطة الإشراف والمراقبة اليومية، وتوجيه التعليمات التقنية، وتحديد وتائر العمل وأوقاته، وإلزام الأجير بقواعد النظام الداخلي وتدابير السلامة والوقاية المعمول بها داخل المنشأة[cite: 2, 3].
\end{itemize}

% ---------------------------------------------------------------------
\section{مفهوم المناولة (المقاولة من الباطن)}

\subsection{التأصيل القانوني لعقد المقاولة من الباطن}
نصت المادة 86 من مدونة الشغل على أن المقاولة من الباطن هي: \textit{«عقد مكتوب يكلف بمقتضاه مقاول أصلي مقاولاً من الباطن بأن يقوم له بشغل من الأشغال، أو بأن ينجز له خدمة من الخدمات، كلما كان ذلك في صالح المقاولة الأصلية ودون إخلال بحقوق أجرائها»}[cite: 2, 3].

والمناولة في جوهرها الاقتصادي والقانوني هي أسلوب من أساليب التنظيم الصناعي والإنتاجي تتوخى من خلاله المقاولة الأصلية تفويض أجزاء من نشاطها إلى مقاولين مستقلين يمتلكون المعرفة التقنية والوسائل اللوجستية التي تفتقر إليها المنشأة الأصلية، وتتخذ في التطبيق شكلين:
\begin{itemize}
    \item \textbf{مناولة التخصص (\textit{Sous-traitance de spécialité}):} حيث تفتقر المقاولة الأصلية للخبرة التقنية أو التجهيزات الدقيقة، فتعمد إلى إسناد أشغال الصيانة المعقدة، أو الأنظمة المعلوماتية، أو الأمن والحراسة والنظافة إلى شركات مختصة تمتلك الدراية والمهارات اللازمة.
    \item \textbf{مناولة السعة (\textit{Sous-traitance de capacité}):} وتلجأ إليها المنشأة عندما تعجز وسائلها الذاتية عن استيعاب حجم الطلبات المتزايدة مؤقتاً، فتسند جزءاً من التصنيع لمقاول خارجي ينجزه في ورشاته الخاصة أو داخل منشأة المقاول الأصلي بوسائل إشراف مستقلة.
\end{itemize}

\subsection{أركان العقد وشروط مشروعيته}
لكي يظل عقد المناولة صحيحاً ومشروعاً ولا ينحرف إلى عملية اتجار غير مشروع بقوة العمل، يجب أن تتوافر فيه جملة من الشروط والأركان الجوهرية[cite: 2]:
\begin{enumerate}
    \item \textbf{شكلية الكتابة:} فرضت المادة 86 إفراغ العقد في محرر مكتوب لبيان حدود العملية وطبيعة المهام[cite: 2, 3].
    \item \textbf{محل العقد (تحقيق نتيجة متخصصة):} يلتزم المقاول من الباطن بإنجاز شغل أو أداء خدمة محددة المعالم تقنياً تنتهي بتحقيق نتيجة ملموسة ومستقلة، وليس مجرد تقديم وتوريد أجراء أو وضع ساعات عمل رهن تصرف المقاول الأصلي[cite: 2].
    \item \textbf{الاستقلال الإداري والتقني الكامل:} يظل المقاول من الباطن هو المشغل الحصري لأجرائه قانونياً وفعلياً؛ فهو الذي يمارس عليهم سلطة الإدارة والتوجيه والإشراف والتأديب عبر أطره ورؤساء ورشاته، دون أن يتدخل المقاول الأصلي في توجيه عمال المناولة بصورة مباشرة[cite: 2].
    \item \textbf{استخدام وسائل ومعدات ذاتية:} يفترض في المقاول من الباطن أن يقدم أدوات العمل ومعداته ومواده الأولية، ما لم تتطلب طبيعة الخدمة استعمال آلات المنشأة الأصلية بمقتضى بند تعاقدي صريح ومبرر تقنياً.
    \item \textbf{شرط الحفاظ على حقوق الأجراء الأصليين:} علق المشرع مشروعية اللجوء إلى المناولة بشرط \textit{«ألا يكون مخلاً بحقوق أجراء المقاولة الأصلية»}[cite: 2, 3]، مانعاً استعمالها لتجاوز مكتسبات الأجراء، أو كوسيلة للإغلاق الجزئي للمنشأة لتسريح العمال وتعويضهم بمقاولات باطن أقل تكلفة[cite: 2].
\end{enumerate}

% ---------------------------------------------------------------------
\section{التمييز بين الشغل المؤقت والمناولة وعقد الشغل المحدد المدة}

\subsection{المقارنة المنهجية بين الشغل المؤقت والمناولة}
لتوضيح الحدود الفاصلة بدقة بين المؤسستين، يقدم الجدول المقارن الآتي تحليلاً تفصيلياً لعناصر التمايز بينهما:

\begin{table}[h]
\centering
\small
\renewcommand{\arraystretch}{1.55}
\begin{tabularx}{\textwidth}{|p{2.8cm}|X|X|}
\hline
\textbf{وجه المقارنة} & \textbf{الشغل المؤقت (\textit{Intérim})} & \textbf{المناولة (\textit{Sous-traitance})} \\
\hline
\textbf{الأساس التشريعي} & المواد 495 وما يليها من مدونة الشغل[cite: 2, 3] & المواد 86 وما يليها من مدونة الشغل[cite: 2, 3] \\
\hline
\textbf{محل الالتزام} & وضع قوة عمل بشرية رهن إشارة المستعمل[cite: 2] & أداء شغل فني أو تقديم خدمة متخصصة متكاملة[cite: 2] \\
\hline
\textbf{طبيعة التزام المقاول} & التزام ببذل عناية (توفير عمال مؤهلين)[cite: 2] & التزام بتحقيق نتيجة محددة ومستقلة[cite: 2] \\
\hline
\textbf{السلطة والإشراف الميداني} & سلطة التوجيه والمراقبة للمقاولة المستعملة[cite: 2] & السلطة الفنية والتوجيهية حصرية لمقاول الباطن[cite: 2] \\
\hline
\textbf{وسائل العمل والتجهيزات} & تعود للمقاولة المستعملة وتوضع تحت تصرف الأجير & يوفرها المقاول من الباطن كأصل عام[cite: 2] \\
\hline
\textbf{المقابل المالي والفوترة} & فاتورة تحتسب على أساس عدد ساعات العمل والأجور & فاتورة إجمالية أو جزافية تحدد بحسب إنجاز الخدمة \\
\hline
\textbf{مدة التعاقد} & محددة بمدد قصوى حصرية بموجب المادة 500[cite: 2, 3] & ترتبط بالمدة الضرورية لإنجاز الصفقة أو المشروع[cite: 2] \\
\hline
\textbf{المسؤولية القانونية} & مسؤولية مباشرة وتضامنية في السلامة والتأمين (المادة 504)[cite: 3] & مسؤولية احتياطية مشروطة بحالات الإعسار (المادتان 89 و90)[cite: 2, 3] \\
\hline
\end{tabularx}
\caption{المقارنة النسقية والمفاهيمية بين الشغل المؤقت وعقد المناولة}
\label{tab:comparatif-detaille}
\end{table}

\subsection{التمييز عن عقد الشغل المحدد المدة (CDD)}
يمثل عقد الشغل المحدد المدة، المنظم بمقتضى المواد 16 وما يليها من مدونة الشغل، الإطار القانوني التقليدي للمرونة المباشرة[cite: 2, 3]:
\begin{itemize}
    \item \textbf{رابطة ثنائية مباشرة:} يبرم العقد مباشرة بين المشغل المستفيد والأجير، وتجتمع في المشغل كافة عناصر التبعية القانونية والفنية دون أي وساطة[cite: 2].
    \item \textbf{الاستثناء من الأصل:} يظل العقد غير محدد المدة (CDI) هو الأصل العام والمهيمن في مدونة الشغل، ولا يجوز إبرام عقد محدد المدة إلا في الحالات المحددة حصراً بنص المادة 16 والمرسوم رقم 2.19.793[cite: 3] (إحلال أجير، تزايد نشاط مؤقت، أو قطاعات موسمية محددة).
    \item \textbf{الجزاء المترتب عن التجاوز:} إذا انتهت المدة المحددة واستمر الأجير في أداء عمله دون اعتراض المشغل، يتحول العقد بقوة القانون وتلقائياً إلى عقد غير محدد المدة، وهو جزاء حمائي صارم يفتقر إليه الشغل المؤقت في بعض صوره التطبيقية بالمغرب.
\end{itemize}

\subsection{إسقاط التمييز المفاهيمي على نازلة الدراسة}
بإنزال هذه المعايير الفقهية والضوابط القانونية على واقعة النازلة التي تشكل منطلق بحثنا — والمتمثلة في تعاقد شركة صناعية كبرى مع شركة مناولة لتوفير عمال يشتغلون داخل خطوط إنتاج المصنع لسنوات تحت إشراف أطر الشركة الأصلية[cite: 1, 2] — تبرز جملة من الخلاصات التحليلية الحاسمة:

\begin{enumerate}
    \item \textbf{انتفاء الصفة الحقيقية لعقد المناولة:} إن قصر موضوع التعاقد على «توفير اليد العاملة والعمالة اليدوية» يجرد العقد من جوهر المناولة القائم على التخصص والنتيجة المستقلة[cite: 2]؛ فالشركة الأصلية لم تسند خدمة متكاملة (كنظافة المنشأة أو حراستها أو صيانة آلاتها)، بل استعارت عمالاً لدمجهم في صلب الدورة الإنتاجية للمصنع.
    \item \textbf{ثبوت ممارسة التبعية المباشرة:} إن خضوع العمال للأوامر اليومية والتوجيه الفني الصادر عن أطر ومهندسي الشركة الأصلية ينفي الاستقلال الإداري والفني لشركة المناولة[cite: 2]، مما يجعل هذه الأخيرة مجرد «مشغل صوري» أو وسيط مالي لتوزيع الأجور.
    \item \textbf{مخالفة قواعد الشغل المؤقت والمدد القانونية:} حتى لو حاول الأطراف تكييف العقد على أنه وضع يد عاملة، فإن المشرع يحظر هذا النشاط على غير مقاولات التشغيل المؤقت المرخصة[cite: 2, 3]، ويمنع استخدام الأجراء لسنوات متواصلة، إذ تحصر المادة 500 المدد القصوى للمهام في مدد لا تتجاوز 3 أو 6 أشهر[cite: 3].
    \item \textbf{إعمال مبدأ أولوية الواقع على المظهر (\textit{Primauté de la réalité}):} بموجب هذا المبدأ المستقر في الاجتهاد القضائي المقارن والمغربي والمكرس في التوصية رقم 198 لمنظمة العمل الدولية، فإن القاضي الاجتماعي لا يقف عند الأوصاف التي أطلقها الأطراف على عقودهم، بل يستنبط الحقيقة من الممارسة اليومية. وبما أن العمال يلبون حاجة قارة ودائمة، فإن العلاقة التعاقدية الحقيقية هي عقد شغل مباشر وغير محدد المدة (CDI) يربط العمال بالشركة الصناعية الأصلية[cite: 1, 2].
\end{enumerate}

\section*{خلاصة الفصل}
يخلص هذا الفصل إلى أن المشرع المغربي، وإن فتح باب المرونة من خلال الشغل المؤقت والمناولة، فإنه أحاطهما بضوابط موضوعية وشكلية دقيقة لمنع استغلالهما كأدوات للالتفاف على استقرار الشغل[cite: 1, 2]. وتكشف دراسة الواقعة الميدانية أن استمرار العمال لسنوات داخل المصنع ينفي عن العقد صفة المناولة ويحوله إلى وساطة مشبوهة تفتقد للسند القانوني. ولتبيان النصوص المنظمة والواجبة التطبيق بتفصيل، ننتقل في الفصل الثاني لدراسة الإطار القانوني للشغل المؤقت والمناولة في مدونة الشغل المغربية[cite: 2, 3].
\cleardoublepage

%% ─── الفصل الثاني: الإطار القانوني ───────────────────────
%% =====================================================================
%% الفصل الثاني: الإطار القانوني في مدونة الشغل المغربية
%% =====================================================================

\chapter{الإطار القانوني في مدونة الشغل المغربية}
\label{chap:chapitre2}

%% ─── مقدمة الفصل ──────────────────────────────────────────
\section*{مقدمة الفصل}
\phantomsection\addcontentsline{toc}{section}{مقدمة الفصل}

أفرد المشرع المغربي في مدونة الشغل (القانون رقم \textenglish{65.99}) منظومة معيارية خاصة لتقنين العلاقات المهنية غير النمطية، مستهدفاً تحقيق التوازن بين متطلبات المقاولة في المرونة التسييرية والقدرة التنافسية من جهة، وبين صيانة النظام العام الاجتماعي وضمان الأمن المهني للأجراء من جهة ثانية. غير أن استقراء النصوص القانونية المنظمة لتقنيتي الشغل المؤقت والمناولة يكشف عن تفاوت هيكلي ملحوظ؛ إذ في الوقت الذي أحاط فيه المشرع الشغل المؤقت بترسانة تفصيلية من القيود الشكلية والموضوعية والمدد الزمنية الصارمة، اكتفى في المقابل بصياغة قواعد عامة ومقتضبة في تنظيم المقاولة من الباطن، مما فتح ثغرات للتحايل في الواقع الميداني.

وللإحاطة بهذا البناء التشريعي، يتناول هذا الفصل أربعة محاور مترابطة:
\begin{itemize}[leftmargin=*, itemsep=4pt]
    \item \textbf{أولاً:} المرجعيات القانونية وهيكلتها في الدستور والمعايير الدولية والمدونة؛
    \item \textbf{ثانياً:} أحكام التشغيل المؤقت في مدونة الشغل (المواد \textenglish{495} إلى \textenglish{506})؛
    \item \textbf{ثالثاً:} أحكام المناولة وعقود الشغل المحددة المدة (المواد \textenglish{86}--\textenglish{91} و\textenglish{16}--\textenglish{17})؛
    \item \textbf{رابعاً:} النصوص التنظيمية وقراءة نقدية في مكامن القصور التشريعي.
\end{itemize}

%% ---------------------------------------------------------------------
\section{المرجعيات القانونية}

\subsection{المرجعيات الدستورية والمعايير الدولية}
تتأصل القواعد الحاكمة للتشغيل غير المباشر في المنظومة الدستورية للمملكة المغربية لسنة \textenglish{2011}؛ حيث ينص الفصل \textenglish{31} على التزام الدولة والمؤسسات العمومية والجماعات الترابية بتعبئة كل الوسائل المتاحة لتيسير أسباب استفادة المواطنات والمواطنين على قدم المساواة من الحق في الشغل والحماية الاجتماعية، فيما يكفل الفصل \textenglish{35} حرية المبادرة والمقاولة والتنافس الحر.

وتتكامل هذه المرجعية الدستورية الوطنية مع معايير منظمة العمل الدولية (\textenglish{OIT})، وعلى رأسها \textbf{الاتفاقية رقم \textenglish{181} بشأن وكالات الاستخدام الخاصة (\textenglish{1997})} التي تُلزم بتقييد الوساطة وإخضاعها للترخيص وحظر اقتطاع أي عمولات من أجور العمال، و\textbf{التوصية رقم \textenglish{198} بشأن علاقة العمل (\textenglish{2006})} التي تكرس مبدأ «أولوية الواقع» عند تكييف العقود الشغلية.

\subsection{التوزيع البنيوي للنصوص داخل مدونة الشغل}
اعتمد المشرع المغربي منهجية الفصل الهيكلي بين مؤسستي الشغل المؤقت والمناولة داخل مدونة الشغل (القانون رقم \textenglish{65.99} الصادر بتنفيذه الظهير الشريف رقم \textenglish{1.03.194} بتاريخ \textenglish{11} أغسطس \textenglish{2003}):
\begin{itemize}[leftmargin=*, itemsep=5pt]
    \item \textbf{أحكام المناولة (المقاولة من الباطن):} أُدرجت ضمن الكتاب الأول المتعلق بالاتفاقيات المتعلقة بالشغل (المواد \textenglish{86} إلى \textenglish{91})، بوصفها عقداً تجارياً ومهنياً يكلف فيه مقاول أصلي مقاولاً من الباطن بإنجاز شغل أو خدمة.
    \item \textbf{أحكام الشغل المؤقت:} أُفردت لها المواد \textenglish{495} إلى \textenglish{506} ضمن \textbf{الباب الثاني من الكتاب الرابع} المتعلق بـ «الوساطة في الاستخدام وتشغيل الأجراء»، مما يعكس فلسفة المشرع في التعامل معها كنشاط استثنائي للوساطة خاضع لأحكام التراخيص والمراقبة الإدارية المقررة في المواد \textenglish{477} وما يليها.
\end{itemize}

%% ---------------------------------------------------------------------
\section{أحكام التشغيل المؤقت (المواد 495 وما يليها)}

\subsection{شروط ممارسة النشاط ونظام الترخيص والكفالة المالية}
انطلاقاً من الطبيعة الحساسة لنشاط التشغيل المؤقت وما يكتنفه من مخاطر المضاربة بقوة العمل، أخضع المشرع تأسيس وتسيير هذه المقاولات لضوابط مالية وإدارية صارمة:
\begin{itemize}[leftmargin=*, itemsep=5pt]
    \item \textbf{الترخيص الإداري المسبق والرأسمال الأدنى (المادتان \textenglish{477} و\textenglish{481}):} لا يجوز فتح مقاولة للتشغيل المؤقت إلا بعد الحصول على إذن مسبق من السلطة الحكومية المكلفة بالشغل، ويشترط ألا يقل رأسمالها عن \textenglish{100.000} درهم مع انتفاء الأحكام الماسة بالشرف أو الأمانة في حق المسيرين.
    \item \textbf{إيداع الكفالة المالية بصندوق الإيداع والتدبير (المادتان \textenglish{482} و\textenglish{488}):} ألزم المشرع في المادة \textenglish{482} كل مقاولة تشغيل مؤقت بإيداع كفالة لدى صندوق الإيداع والتدبير (\textenglish{CDG}) تعادل \textenglish{50} مرة القيمة الإجمالية السنوية للحد الأدنى القانوني للأجر؛ وتسمح المادة \textenglish{488} للمحكمة المختصة، عند إعسار الوكالة أو سحب ترخيصها، بتخصيص هذه الكفالة لأداء مستحقات الأجراء وواجبات الصندوق الوطني للضمان الاجتماعي (\textenglish{CNSS}).
    \item \textbf{سحب الترخيص واستمرار الالتزامات (المادتان \textenglish{487} و\textenglish{503}):} يجوز للسلطة الحكومية سحب الإذن بالممارسة بقرار معلل ودون تعويض (المادة \textenglish{487})، غير أن السحب لا يُعفي مسؤولي المقاولة من الوفاء بالتزاماتهم تجاه الأجراء و\textenglish{CNSS} (المادة \textenglish{503}).
\end{itemize}

\subsection{الحالات الحصرية والمدد القصوى والوقاية والسلامة (المواد 496--505)}
رسخت المادة \textenglish{496} مبدأً جوهرياً يقضي بأن اللجوء إلى الشغل المؤقت هو \textbf{استثناء محصور في أربع حالات فقط} بعد استشارة الهيئات التمثيلية للأجراء: إحلال أجير محل آخر (في غير حالة الإضراب)، أو التزايد المؤقت في النشاط، أو الأشغال الموسمية، أو الأشغال التي استقر العرف على عدم اللجوء فيها لعقد غير محدد المدة.

كما حظرت المادة \textenglish{497} اللجوء إلى الأجراء المؤقتين لإنجاز أشغال تكتسي خطورة خاصة، ومنعت المادة \textenglish{498} استخدامهم لمواجهة تزايد مؤقت في المناصب التي شملها فصل لأسباب اقتصادية خلال السنة الموالية للفصل. وألزمت المادة \textenglish{504} المقاولة المستعملة باتخاذ كافة تدابير الصحة والسلامة والتأمين عن حوادث الشغل والأمراض المهنية (وفق القانون رقم \textenglish{18.12}). ويوضح الجدول الآتي ضوابط المدة وفترة التجربة والعقوبة الزجرية:

\begin{table}[htbp]
\centering
\small
\renewcommand{\arraystretch}{1.55}
\begin{tabularx}{\textwidth}{|p{4.5cm}|p{4.5cm}|X|}
\hline
\textbf{حالة اللجوء (المادة \textenglish{496})} & \textbf{المدة القصوى للمهمة (المادة \textenglish{500})} & \textbf{فترة التجربة (المادة \textenglish{502})} \\
\hline
\textbf{إحلال أجير محل آخر} & طيلة مدة توقف عقد الأجير المتغيب & يومان إذا كان العقد مبرماً لأقل من شهر \\
\cline{1-2}
\textbf{تزايد مؤقت في نشاط المقاولة} & \textenglish{3} أشهر قابلة للتجديد مرة واحدة فقط & \textenglish{3} أيام إذا تراوحت مدة العقد بين شهر وشهرين \\
\cline{1-2}
\textbf{أشغال موسمية وأعراف مهنية} & \textenglish{6} أشهر غير قابلة للتجديد & \textenglish{5} أيام إذا فاقت مدة العقد شهرين \\
\hline
\multicolumn{3}{|r|}{\textbf{الجزاء الجنائي (المادة \textenglish{505}):} غرامة مالية من \textenglish{2.000} إلى \textenglish{5.000} درهم عن مخالفة أحكام هذا الباب.} \\
\hline
\end{tabularx}
\caption{الضوابط الإجرائية والمدد الزمنية وفترة التجربة في الشغل المؤقت}
\label{tab:durees-temporaire}
\end{table}

%% ---------------------------------------------------------------------
\section{أحكام المناولة وعقود الشغل المحددة المدة}

\subsection{أحكام عقد المقاولة من الباطن (المواد 86 إلى 91)}
نصت المادة \textenglish{87} من مدونة الشغل على أن المقاول من الباطن يُعتبر مشغلاً خاضعاً لكافة أحكام المدونة، ونصوص الضمان الاجتماعي، والتعويض عن حوادث الشغل والأمراض المهنية. وأوجبت المادة \textenglish{88} تعليق اسم المقاول الأصلي وعنوانه في كل ورش أو منشأة وتحديد بيانات الأداء وبطاقة الشغل وفق الشكليات القانونية متى اشتغل الأجراء في منشآت المقاول الأصلي.

أما بشأن \textbf{المسؤولية الاحتياطية للمقاول الأصلي}، فقد نصت المادتان \textenglish{89} و\textenglish{90} على أنه إذا كان المقاول من الباطن:
\begin{itemize}[leftmargin=*, itemsep=4pt]
    \item \textbf{غير مقيد في السجل التجاري}؛
    \item \textbf{ولا يتوفر على أصل تجاري}؛
    \item \textbf{وكان في حالة إعسار}؛
\end{itemize}
أصبح المقاول الأصلي مسؤولاً—في حدود المبالغ المترتبة بذمته لفائدة المقاول من الباطن—عن أداء الأجور، والتعويض عن العطلة السنوية، وتعويضات الفصل، واشتراكات \textenglish{CNSS}، والرسم الخاص بالتكوين المهني، مع تخويل الأجراء و\textenglish{CNSS} حق إقامة دعوى مباشرة ضده (المادة \textenglish{90}). كما يسأل المقاول الأصلي بصفة مباشرة عن التعويض عن حوادث الشغل والأمراض المهنية إذا أنجزت الأشغال أو الخدمات داخل مؤسساته أو ملحقاتها (الفقرة الأخيرة من المادة \textenglish{89}).

\subsection{ضوابط عقد الشغل المحدد المدة (المادتان 16 و17)}
ينظم القسم الخاص بانعقاد الشغل عقد الشغل المحدد المدة (\textenglish{CDD}) كآلية مرونة ثنائية مباشرة بين المشغل والأجير دون وسيط:
\begin{itemize}[leftmargin=*, itemsep=4pt]
    \item \textbf{شروط الصحة (المادة \textenglish{16}):} لا يُبرم عقد محدد المدة إلا في حالات استثنائية (إحلال أجير، ازدياد مؤقت في النشاط، عمل موسمي)، أو في القطاعات والحالات الاستثنائية المحددة بنص تنظيمي، أو عند فتح مقاولة لأول مرة أو إطلاق منتوج جديد (لمدة أقصاها سنة قابلة للتجديد مرة واحدة في القطاعات غير الفلاحية).
    \item \textbf{التحويل التلقائي (المادة \textenglish{17}):} ينص القانون صراحة على أن عقد الشغل المحدد المدة يتحول بقوة القانون إلى عقد غير محدد المدة إذا استمر الأجير في عمله بعد انتهاء المدة المحددة فيه، وهي قاعدة صريحة في الـ \textenglish{CDD} لكنها غير منصوص عليها بنص مماثل في باب التشغيل المؤقت.
\end{itemize}

%% ---------------------------------------------------------------------
\section{النصوص التنظيمية وقراءة نقدية في مكامن القصور}

\subsection{المراسيم التطبيقية والاتفاقيات الجماعية}
تعزز الإطار التشريعي للمدونة بمراسيم تطبيقية، إلى جانب دور الاتفاقيات الجماعية (المواد \textenglish{104} وما يليها) في تحسين ظروف الشغل القطاعية:
\begin{enumerate}[leftmargin=*, itemsep=5pt]
    \item \textbf{المرسوم رقم \textenglish{2.19.456} (الصادر في \textenglish{30} يناير \textenglish{2020}، الجريدة الرسمية عدد \textenglish{6855}):} المحدد لتكوين وطريقة عمل اللجنة المختصة بمقاولات التشغيل المؤقت المنصوص عليها في المادة \textenglish{496}، والتي تتألف من \textenglish{18} عضواً وفق تركيبة ثلاثية (الإدارة، المنظمات المهنية للمشغلين، والمنظمات النقابية للأجراء الأكثر تمثيلاً)، وقد حل محل المرسوم السابق رقم \textenglish{2.04.464} الصادر سنة \textenglish{2004}.
    \item \textbf{المرسوم رقم \textenglish{2.19.793} (الصادر في \textenglish{23} يوليو \textenglish{2020}):} المحدد للقطاعات والحالات الاستثنائية التي يمكن فيها إبرام عقد الشغل محدد المدة استناداً إلى المادة \textenglish{16} من المدونة.
\end{enumerate}

\subsection{مكامن القصور التشريعي والثغرات الهيكلية}
تكشف القراءة النقدية لأحكام مدونة الشغل عن ثغرات بنيوية تتكرر آثارها في الممارسة الميدانية:
\begin{itemize}[leftmargin=*, itemsep=5pt]
    \item \textbf{اختلال التوازن بين الشغل المؤقت والمناولة:} مقابل التقييد الصارم للشغل المؤقت بمدد لا تتعدى \textenglish{3} إلى \textenglish{6} أشهر (المادة \textenglish{500})، وردت أحكام المناولة (المواد \textenglish{86}--\textenglish{91}) فضفاضة دون تحديد مدد قصوى أو حالات حصرية، مما شجع على التحايل عبر «مناولة مقنعة».
    \item \textbf{ضيق شروط المادة \textenglish{89}:} يعفي اشتراط عدم تقييد مقاول الباطن في السجل التجاري وعدم امتلاكه أصلاً تجارياً المقاولَ الأصلي من المسؤولية بمجرد إنشاء شركة شكلية مقيدة بالسجل التجاري، كما يتيح حصر الأداء «في حدود ما تبقى في الذمة» التحايل بتصفية الفواتير مسبقاً.
    \item \textbf{غياب المساواة الأجرية الصريحة وضعف العقوبات:} تخلو المدونة من نص صريح يُلزم بالمساواة في الأجر بين الأجير المؤقت والأجير الدائم (عكس القانون الفرنسي)، كما أن غرامة المادة \textenglish{505} (\textenglish{2.000} إلى \textenglish{5.000} درهم) ضعيفة ولا توازيها قاعدة تشريعية صريحة تقرر إعادة التكييف التلقائي.
\end{itemize}

%% ─── خلاصة الفصل ──────────────────────────────────────────
\section*{خلاصة الفصل}
\phantomsection\addcontentsline{toc}{section}{خلاصة الفصل}

يتضح مما سبق أن الإطار القانوني المغربي يضع قواعد دقيقة للشغل المؤقت وأحكاماً عامة للمناولة، غير أنه يشكو من ثغرات تشريعية أبرزها ضيق نطاق المادة \textenglish{89}، وهزالة غرامة المادة \textenglish{505}، وسكوت النص عن المساواة الأجرية وعن إعادة التكييف التلقائي. وللوقوف على الضوابط التي قيد بها المشرع اللجوء إلى هذه الآليات، ننتقل في الفصل الثالث إلى تحليل الحالات الحصرية التي يسمح فيها القانون باللجوء إلى الشغل المؤقت.

\cleardoublepage

%% ─── الفصل الثالث: الحالات المسموح بها ───────────────────
%% ============================================================
%%  chapitre3.tex — الفصل الثالث: الحالات التي يسمح فيها القانون
%%                  باللجوء إلى الشغل المؤقت
%% ============================================================

\chapter{الحالات التي يسمح فيها القانون باللجوء إلى الشغل المؤقت}
\label{chap:chapitre3}

%% ─── مقدمة الفصل ──────────────────────────────────────────
\section*{مقدمة الفصل}
\phantomsection\addcontentsline{toc}{section}{مقدمة الفصل}

الأصل في مدونة الشغل أن تقوم علاقة الشغل على الاستقرار، وأن يكون العقد غير محدد المدة (\textenglish{CDI}) هو القاعدة. ولذلك لم يُجز المشرع اللجوء إلى الأجير المؤقت إلا استثناءً، وفي أربع حالات أوردتها المادة \textenglish{496} على سبيل الحصر، لا يجوز التوسع فيها ولا القياس عليها. ويُشترط فضلاً عن ذلك أن تستشير المقاولة المستعملة الهيئات التمثيلية لأجرائها قبل اللجوء إلى هذه الصيغة.

ويجيب هذا الفصل عن سؤالين محوريين من أسئلة النقاش في النازلة: \textbf{السؤال الثالث} (ما الحالات التي يسمح فيها القانون باللجوء إلى شركات التشغيل المؤقت؟) و\textbf{السؤال الخامس} (هل يمكن استعمال الشغل المؤقت لتلبية حاجيات دائمة ومستمرة للمقاولة؟)، عبر أربعة محاور متوازنة:
\begin{itemize}[leftmargin=*, itemsep=4pt]
    \item \textbf{أولاً:} استبدال أجير توقف عقده؛
    \item \textbf{ثانياً:} ارتفاع مؤقت في نشاط المقاولة؛
    \item \textbf{ثالثاً:} الأشغال الموسمية والأشغال التي استقر العرف بشأنها (الظرفية)؛
    \item \textbf{رابعاً:} الحالات الممنوعة ومدى جواز تلبية الحاجيات الدائمة وإسقاطها على المجموعة \textenglish{8}.
\end{itemize}

\begin{table}[htbp]
\centering
\small
\renewcommand{\arraystretch}{1.5}
\begin{tabular}{|p{4.4cm}|p{4.6cm}|p{4.4cm}|}
\hline
\textbf{حالة اللجوء (المادة \textenglish{496})} & \textbf{الشرط الأساسي} & \textbf{المدة القصوى (المادة \textenglish{500})} \\
\hline
استبدال أجير توقف عقده & توقف مؤقت لعقد أجير محدد، لغير سبب الإضراب & مدة توقف عقد الأجير المعوَّض \\
\hline
ارتفاع مؤقت في النشاط & زيادة فعلية وعابرة في حجم العمل & \textenglish{3} أشهر قابلة للتجديد مرة واحدة \\
\hline
الأشغال الموسمية & نشاط يتكرر في فترات محددة من السنة & \textenglish{6} أشهر غير قابلة للتجديد \\
\hline
الأشغال التي جرى العرف بعدم إبرام عقود غير محددة المدة بشأنها & طبيعة القطاع أو النشاط تجعل العمل مؤقتاً بطبيعته & \textenglish{6} أشهر غير قابلة للتجديد \\
\hline
\end{tabular}
\caption{الحالات الحصرية للجوء إلى الشغل المؤقت وشروطها ومددها}
\label{tab:cas-recours}
\end{table}

%% ─────────────────────────────────────────────────────────
\section{استبدال أجير توقف عقده}

\subsection{مضمون الحالة ونطاق التوقف المشروع}
تسمح المادة \textenglish{496} بوضع أجير مؤقت رهن إشارة المقاولة المستعملة \textbf{لتعويض أجير غائب أو توقف عقد شغله}. ويقتضي ذلك أن يكون عقد الأجير الأصلي قائماً لكنه معلَّق مؤقتاً، بحيث يعود صاحبه إلى منصبه عند زوال سبب التوقف. ومن أبرز صور هذا التوقف المشروع:
\begin{itemize}[leftmargin=*, itemsep=4pt]
    \item الغياب بسبب المرض غير المهني أو حادثة الشغل أو المرض المهني.
    \item إجازة الولادة أو التمديد المرتبط بها.
    \item العطلة السنوية المؤدى عنها أو التغيبات المرخص بها قانوناً أو اتفاقياً.
    \item الإيقاف المؤقت للعقد لأي سبب آخر يقره القانون (كالخدمة أو التفرغ النقابي المرخص).
\end{itemize}

\subsection{شروطها وحدودها القانونية}
\begin{enumerate}[leftmargin=*, itemsep=6pt]
    \item \textbf{استثناء الإضراب:} لا يجوز اللجوء إلى الأجير المؤقت إذا كان توقف العقد ناتجاً عن ممارسة حق الإضراب (المادة \textenglish{496})، حتى لا يُستعمل الشغل المؤقت أداةً لكسر الإضرابات وإفراغ هذا الحق الدستوري من محتواه.
    \item \textbf{ارتباط المدة بالغياب:} لا تتجاوز مدة المهمة مدة توقف عقد الأجير المعوَّض (المادة \textenglish{500})، فإذا عاد هذا الأخير انتهى مبرر المهمة فوراً.
    \item \textbf{تحديد الأجير المعوَّض:} يجب أن يُذكر سبب اللجوء في العقد المكتوب بين مقاولة التشغيل المؤقت والمقاولة المستعملة (المادة \textenglish{499})، مما يقتضي عملياً تحديد الأجير الذي يُعوَّض وسبب غيابه، حتى يتسنى لمفتش الشغل التحقق من صحة السبب.
\end{enumerate}

%% ─────────────────────────────────────────────────────────
\section{ارتفاع مؤقت في النشاط}

\subsection{مضمون التزايد المؤقت ومعايير تحققه}
تسمح المادة \textenglish{496} باللجوء إلى الأجير المؤقت لمواجهة \textbf{تزايد مؤقت في نشاط المقاولة}، كطلبية استثنائية، أو ذروة غير معتادة في الإنتاج، أو مشروع عابر يتجاوز طاقة الأجراء الدائمين. والمعيار هنا مزدوج: أن تكون الزيادة \textbf{حقيقية} يمكن إثباتها محاسبياً وتجارياً، وأن تكون \textbf{مؤقتة} بطبيعتها، أي أن ينتظر زوالها في أجل منظور، لا أن تمثل توسعاً دائماً في حجم الإنتاج.

\subsection{القيود الزمنية والمنع بعد الفصل الاقتصادي}
\begin{enumerate}[leftmargin=*, itemsep=6pt]
    \item \textbf{مدة محدودة:} ثلاثة أشهر قابلة للتجديد \textbf{مرة واحدة فقط} (المادة \textenglish{500})، أي ستة أشهر كحد أقصى. فإذا استمرت «الزيادة» بعد ذلك، فهي في الحقيقة مستوى عادي من النشاط يستوجب تشغيلاً دائماً.
    \item \textbf{منع اللجوء بعد الفصل لأسباب اقتصادية:} لا يجوز للمقاولة التي قامت بفصل أجراء لأسباب اقتصادية أن تلجأ إلى أجراء مؤقتين لمواجهة تزايد مؤقت في النشاط، وذلك خلال السنة الموالية للفصل وبالنسبة للمناصب التي شملها (المادة \textenglish{498})، مع مراعاة أحكام الأولوية في إعادة التشغيل وفق المادة \textenglish{508}.
    \item \textbf{عدم التحايل بالتعاقب:} لا يجوز تحويل هذه الحالة إلى غطاء دائم بتعاقب مهام قصيرة على المنصب نفسه أو بتدوير الأجراء، لأن العبرة بطبيعة الحاجة لا بشكل العقود المتتالية.
\end{enumerate}

%% ─────────────────────────────────────────────────────────
\section{الأشغال الموسمية والظرفية}

تجمع هذه الفقرة حالتين أوردتهما المادة \textenglish{496} تشتركان في أن العمل فيهما مؤقت بطبيعته الموضوعية، لا بقرار إداري من المشغل.

\subsection{مضمون الأشغال الموسمية والأشغال التي استقر العرف بشأنها}
\begin{itemize}[leftmargin=*, itemsep=6pt]
    \item \textbf{الأشغال الموسمية:} هي الأشغال التي تتكرر في فترات محددة من السنة تبعاً لإيقاع طبيعي أو اقتصادي، كجني المحاصيل الفلاحية وتلفيفها، أو ذروة الموسم السياحي، أو بعض الصناعات الغذائية المرتبطة بمواسم الإنتاج. ويقوم معيارها على التكرار الدوري للنشاط وانقطاعه خارج الموسم.
    \item \textbf{الأشغال التي جرى العرف بعدم إبرام عقود غير محددة المدة بشأنها (الأشغال الظرفية):} وهي أنشطة يكون العمل فيها بالمهمة أو بالمشروع عرفاً مستقراً في القطاع بالنظر إلى طبيعتها المؤقتة، كالإنتاج السمعي البصري والسينمائي، وتنظيم التظاهرات الثقافية والرياضية والمعارض.
\end{itemize}

\subsection{النظام الزمني الصارم (ستة أشهر غير قابلة للتجديد)}
تحدد المادة \textenglish{500} مدة المهمة في هاتين الحالتين في \textbf{ستة أشهر غير قابلة للتجديد}. ويعكس هذا السقف الصارم منطق المشرع: فالموسم أو التظاهرة الظرفية محدودة بطبيعتها في الزمان، وما يتجاوز ستة أشهر لم يعد موسماً أو نشاطاً عرضياً بل عملاً مستمراً يفرض اللجوء إلى التعاقد الدائم.

%% ─────────────────────────────────────────────────────────
\section{الحالات الممنوعة (الحاجيات الدائمة والمستمرة)}

\subsection{الموانع التشريعية الصريحة ومدى جواز تلبية الحاجيات الدائمة}
أقر المشرع جملة من الموانع الآمرة التي يُحظر فيها اللجوء إلى الشغل المؤقت:
\begin{tcolorbox}[colback=red!5, colframe=accentcolor, arc=4pt,
    title={\textbf{حالات يُمنع فيها اللجوء إلى الشغل المؤقت}}]
\begin{itemize}[leftmargin=*, itemsep=4pt]
    \item \textbf{تعويض الأجراء المضربين} (المادة \textenglish{496}).
    \item \textbf{الأشغال التي تكتسي خطورة خاصة} (المادة \textenglish{497}).
    \item \textbf{مواجهة تزايد مؤقت في النشاط داخل المناصب التي شملها فصل لأسباب اقتصادية}، خلال السنة الموالية للفصل (المادة \textenglish{498}).
    \item \textbf{كل لجوء خارج الحالات الأربع للمادة \textenglish{496}}، أو دون استشارة الهيئات التمثيلية للأجراء.
\end{itemize}
\end{tcolorbox}

وبشأن التساؤل حول \textbf{إمكانية استعمال الشغل المؤقت لتلبية حاجيات دائمة ومستمرة للمقاولة}، فإن الجواب هو \textbf{بالنفي القاطع}، وإن لم تصُغه مدونة الشغل في جملة منفردة على غرار المادة \textfrench{L. 1251-5} من قانون العمل الفرنسي. ويُستخلص هذا المنع في القانون المغربي من ثلاثة أسس متكاملة:
\begin{enumerate}[leftmargin=*, itemsep=5pt]
    \item \textbf{الطابع الحصري للحالات:} كل حالات المادة \textenglish{496} مبنية على فكرة الحاجة العابرة، ولا مكان بينها للحاجة الدائمة.
    \item \textbf{المدد القصوى:} لا تتجاوز المهمة في أي حالة ستة أشهر (المادة \textenglish{500}) باستثناء الإحلال المرتبط بمدة غياب محدد، ومن ثم يستحيل قانوناً أن يشغل أجير مؤقت منصباً قاراً لسنوات.
    \item \textbf{مبدأ استقرار الشغل:} باعتبار العقد غير محدد المدة هو الأصل العام، يُفسَّر كل استثناء تفسيراً ضيقاً. ويترتب عن خرق هذا المنع جزاء زجري في المادة \textenglish{505} (غرامة \textenglish{2.000} إلى \textenglish{5.000} درهم)، وجزاء مدني يستخلصه القضاء بإعادة التكييف إلى عقد غير محدد المدة.
\end{enumerate}

\subsection{إسقاط الضوابط الحصرية على وضعية المجموعة 8}
في الوضعية المدروسة للمجموعة \textenglish{8}، يشتغل العمال الذين توفرهم شركة المناولة داخل المصنع \textbf{منذ سنوات}، ويؤدّون المهام ذاتها التي يقوم بها أجراء الشركة الأصلية. ولا تندرج هذه الوضعية في أي من الحالات الأربع للمادة \textenglish{496}:
\begin{itemize}[leftmargin=*, itemsep=4pt]
    \item ليس هناك أجير غائب يُعوَّض، بل مناصب قارة في خط الإنتاج.
    \item لا يمكن الحديث عن زيادة مؤقتة في النشاط تستمر لسنوات، والحد الأقصى هو ستة أشهر.
    \item لا يتعلق الأمر بنشاط موسمي ولا بقطاع جرى عرفه على العمل المؤقت.
\end{itemize}
فالأمر يتعلق إذن بتلبية حاجيات دائمة ومستمرة، وهو ما يحظره القانون. ويفسر ذلك لجوء بعض المقاولات إلى تسمية العقد «مناولة» بدل «شغل مؤقت»، غير أن العبرة في القضاء بحقيقة العلاقة لا بتسميتها.

%% ─── خلاصة الفصل ──────────────────────────────────────────
\section*{خلاصة الفصل}
\phantomsection\addcontentsline{toc}{section}{خلاصة الفصل}

حصر المشرع المغربي اللجوء إلى الشغل المؤقت في أربع حالات تجمعها فكرة الحاجة العابرة، وأحاطها بمدد قصوى لا تتجاوز ستة أشهر في الغالب، وبموانع صريحة تتعلق بالإضراب والأشغال الخطيرة والفصل الاقتصادي. ويُستخلص أن الشغل المؤقت لا يمكن قانوناً أن يُستعمل لتلبية حاجيات دائمة ومستمرة، وهو ما يجعل وضعية عمال المجموعة \textenglish{8} خارج كل إطار مشروع. ويبقى بعد ذلك تحديد الالتزامات الملقاة على مقاولة التشغيل المؤقت والمقاولة المستعملة، وهو موضوع الفصل الرابع.

\cleardoublepage

%% ─── الفصل الرابع: الالتزامات ────────────────────────────
%% ============================================================
%%  chapitre4.tex — الفصل الرابع: التزامات شركة التشغيل المؤقت
%%                  والمقاولة المستعملة
%% ============================================================

\chapter{التزامات شركة التشغيل المؤقت والمقاولة المستعملة}
\label{chap:chapitre4}

\section*{مقدمة الفصل}

تقوم علاقة الشغل المؤقت على بنية ثلاثية تتوزع فيها صفة المشغل بين طرفين: فمقاولة التشغيل المؤقت هي المشغل القانوني الذي يبرم العقد ويؤدي الأجر، بينما تمارس المقاولة المستعملة سلطة الإدارة والتوجيه اليومي داخل منشأتها. ولتفادي أن يضيع الأجير بين الطرفين، وزّع المشرع المغربي الالتزامات بينهما توزيعاً دقيقاً في المواد 477 إلى 506 من مدونة الشغل.

ويجيب هذا الفصل عن السؤال الرابع من أسئلة النقاش: \textbf{ما الالتزامات القانونية لشركة التشغيل المؤقت تجاه الأجراء؟} فنعرض التزامات مقاولة التشغيل المؤقت (أولاً)، ثم التزامات المقاولة المستعملة (ثانياً)، ثم الضوابط الشكلية للعقود ومدتها وتجديدها (ثالثاً)، لنختم بجرد حقوق الأجير المؤقت وضماناته (رابعاً). وتفادياً للتكرار، نحيل على الفصل الثاني بخصوص التفاصيل التي سبق عرضها.

%% ─────────────────────────────────────────────────────────
\section{التزامات شركة التشغيل المؤقت}

تعرّف المادة 495 مقاولة التشغيل المؤقت بأنها شخص معنوي مستقل عن السلطة العمومية، يقتصر نشاطه على وضع أجراء رهن إشارة الغير بصفة مؤقتة، وهي التي \textbf{تشغّل هؤلاء الأجراء وتلتزم بأداء أجورهم وبالوفاء بالالتزامات القانونية الناشئة عن عقد الشغل}. ويترتب عن ذلك أن صفة المشغل القانوني تبقى لها طيلة مدة المهمة، ولو كان الأجير يعمل فعلياً داخل منشأة أخرى.

\subsection{الالتزامات السابقة لمزاولة النشاط}
لا يجوز لأي شركة أن تضع أجراء رهن إشارة الغير إلا إذا استوفت شروطاً مسبقة، سبق تفصيلها في الفصل الثاني، ونذكّر بأهمها:
\begin{itemize}[leftmargin=*, itemsep=4pt]
    \item \textbf{الترخيص المسبق} من السلطة الحكومية المكلفة بالشغل (المادة 477).
    \item \textbf{رأسمال لا يقل عن 100.000 درهم}، وانتفاء الأحكام الماسة بالشرف أو بالأمانة في حق المسيرين (المادة 481).
    \item \textbf{إيداع كفالة} لدى صندوق الإيداع والتدبير تعادل خمسين مرة القيمة الإجمالية السنوية للحد الأدنى القانوني للأجر (المادة 482).
    \item \textbf{حصر النشاط} في التشغيل المؤقت دون غيره (المادة 495).
\end{itemize}

\subsection{الالتزامات تجاه الأجير المؤقت}
\begin{enumerate}[leftmargin=*, itemsep=6pt]
    \item \textbf{أداء الأجر والوفاء بالالتزامات القانونية للمشغل} (المادة 495)، بما فيها التصريح بالأجير لدى الصندوق الوطني للضمان الاجتماعي (\textfrench{CNSS}).
    \item \textbf{إبرام عقد مكتوب مع الأجير} يتضمن البيانات الإلزامية التي تحددها المادة 501 (انظر القسم الثالث من هذا الفصل).
    \item \textbf{احترام المدد القصوى لفترة التجربة}: يومان إذا كان العقد أقل من شهر، وثلاثة أيام إذا تراوحت مدته بين شهر وشهرين، وخمسة أيام إذا تجاوز شهرين (المادة 502).
    \item \textbf{مجانية الخدمة}: يُمنع على الوكالة أن تتقاضى من الباحث عن الشغل أي أجر أو أتعاب، بصفة مباشرة أو غير مباشرة، كلياً أو جزئياً (المادة 480). كما يُمنع على مسيريها تسلّم أي ودائع أو ضمانات بمناسبة التشغيل (المادة 485).
    \item \textbf{عدم التمييز}: يُحظر كل تمييز مبني على العرق أو اللون أو الجنس أو الدين أو الرأي السياسي أو الأصل الوطني أو الاجتماعي أو الانتماء النقابي (المادة 478).
    \item \textbf{حماية المعطيات الشخصية}: لا تعالج الوكالة من معطيات الباحثين عن الشغل إلا ما يتعلق بمؤهلاتهم وخبرتهم المهنية، مع احترام حياتهم الخاصة (المادة 479).
\end{enumerate}

\subsection{الالتزامات تجاه الإدارة}
\begin{itemize}[leftmargin=*, itemsep=4pt]
    \item توجيه بيان مفصّل عن خدماتها إلى المصالح الإقليمية للتشغيل عند نهاية كل نصف سنة (المادة 484).
    \item مسك سجل وفق النموذج الذي تحدده السلطة المكلفة بالشغل، يُمكّن مفتش الشغل من مراقبة احترام القانون (المادة 486).
    \item بقاء التزاماتها تجاه الأجراء والصندوق الوطني للضمان الاجتماعي قائمة حتى بعد سحب الترخيص (المادة 503).
\end{itemize}

%% ─────────────────────────────────────────────────────────
\section{التزامات المقاولة المستعملة}

لا يعفي وجود وسيط مرخّص المقاولةَ المستعملة من كل التزام؛ فبما أنها هي التي تستفيد من العمل وتمارس سلطة التوجيه، أخضعها المشرع لالتزامات تسبق المهمة وأخرى تواكبها.

\subsection{الالتزامات السابقة للجوء إلى الأجير المؤقت}
\begin{enumerate}[leftmargin=*, itemsep=6pt]
    \item \textbf{استشارة الهيئات التمثيلية للأجراء} داخل المقاولة قبل اللجوء إلى التشغيل المؤقت (المادة 496).
    \item \textbf{التقيد بالحالات الحصرية الأربع} المنصوص عليها في المادة 496 (موضوع الفصل الثالث)، وعدم استعمال الأجير المؤقت لسدّ حاجيات دائمة.
    \item \textbf{عدم إسناد أشغال تكتسي خطورة خاصة} إلى الأجراء المؤقتين (المادة 497).
    \item \textbf{عدم اللجوء إلى الأجراء المؤقتين لمواجهة تزايد مؤقت في النشاط} داخل المناصب التي شملها فصل لأسباب اقتصادية، وذلك خلال السنة الموالية للفصل (المادة 498).
    \item \textbf{إبرام عقد مكتوب مع مقاولة التشغيل المؤقت} يبيّن سبب اللجوء ومدة المهمة ومكان تنفيذها ومقابل وضع الأجير رهن الإشارة (المادة 499).
\end{enumerate}

\subsection{الالتزامات أثناء تنفيذ المهمة}
\begin{itemize}[leftmargin=*, itemsep=6pt]
    \item \textbf{الصحة والسلامة}: اتخاذ جميع تدابير الوقاية والحماية لفائدة الأجراء المؤقتين، والتأمين عليهم ضد حوادث الشغل والأمراض المهنية (المادة 504). وهذا هو الالتزام الوحيد الذي حمّله المشرع صراحةً للمقاولة المستعملة بصفة مباشرة.
    \item \textbf{احترام المدة القصوى للمهمة} المحددة في المادة 500، وعدم الإبقاء على الأجير بعد انقضائها تحت غطاء تجديدات متتالية.
    \item \textbf{تطبيق قواعد تنظيم العمل السارية في المنشأة} (ساعات العمل، الراحة، حفظ الصحة)، بحكم أن الأجير يشتغل تحت إدارتها اليومية.
\end{itemize}

\begin{tcolorbox}[colback=lightgray, colframe=maincolor, arc=4pt,
    title={\textbf{تنبيه: المساواة في الأجر}}]
لا تتضمن مدونة الشغل نصاً صريحاً يُلزم المقاولة المستعملة بأن تمنح الأجير المؤقت أجراً مساوياً لأجر الأجير الدائم الذي يشغل المنصب نفسه، خلافاً للتشريع الفرنسي (المادة \textfrench{L1251-18} من قانون الشغل الفرنسي). أما المادة 9 فتحظر التمييز في الأجر على أسس محددة حصراً (العرق، اللون، الجنس، الإعاقة، الحالة الزوجية، الدين، الرأي السياسي، الانتماء النقابي، الأصل الوطني أو الاجتماعي)، وليس من بينها طبيعة العقد. لذلك يبقى الفارق في الأجر، كما في وضعية المجموعة 8، مسألةً يُعالجها الاجتهاد القضائي (الفصل التاسع) أو إعادة تكييف العلاقة (الفصل الخامس).
\end{tcolorbox}

%% ─────────────────────────────────────────────────────────
\section{العقود الكتابية ومدتها وتجديدها}

يقوم الشغل المؤقت على عقدين مكتوبين مترابطين: عقد وضع الأجير رهن الإشارة بين المقاولتين، وعقد المهمة بين مقاولة التشغيل المؤقت والأجير. ويوضح الجدول الآتي البيانات الإلزامية لكل منهما.

\begin{table}[h]
\centering
\small
\renewcommand{\arraystretch}{1.5}
\begin{tabular}{|p{3.2cm}|p{5.2cm}|p{5.2cm}|}
\hline
 & \textbf{عقد وضع الأجير رهن الإشارة (المادة 499)} & \textbf{عقد المهمة مع الأجير (المادة 501)} \\
\hline
\textbf{الأطراف} & مقاولة التشغيل المؤقت والمقاولة المستعملة & مقاولة التشغيل المؤقت والأجير \\
\hline
\textbf{الشكل} & مكتوب وجوباً & مكتوب وجوباً \\
\hline
\textbf{البيانات الإلزامية} & سبب اللجوء إلى الأجير المؤقت؛ مدة المهمة؛ مكان تنفيذها؛ مقابل وضع الأجير رهن الإشارة & بيانات المادة 499؛ المؤهلات المطلوبة؛ الأجر وكيفية أدائه؛ فترة التجربة؛ خصائص المنصب؛ رقم انخراط المقاولة ورقم تسجيل الأجير في الضمان الاجتماعي؛ شرط الترحيل إذا نُفّذت المهمة خارج المغرب \\
\hline
\textbf{بند خاص} & -- & إمكانية تشغيل الأجير من طرف المقاولة المستعملة بعد انتهاء المهمة \\
\hline
\end{tabular}
\caption{البيانات الإلزامية لعقدي الشغل المؤقت}
\label{tab:contrats-temporaire}
\end{table}

\subsection{المدة والتجديد}
ربطت المادة 500 مدة المهمة بسبب اللجوء إليها (انظر الجدول \ref{tab:durees-temporaire} في الفصل الثاني):
\begin{itemize}[leftmargin=*, itemsep=4pt]
    \item \textbf{إحلال أجير محل آخر}: لا تتجاوز المهمة مدة توقف عقد الأجير المعوَّض.
    \item \textbf{تزايد مؤقت في النشاط}: ثلاثة أشهر قابلة للتجديد \textbf{مرة واحدة فقط}، أي ستة أشهر كحد أقصى.
    \item \textbf{الأشغال الموسمية والأشغال التي جرى العرف بعدم إبرام عقود غير محددة المدة بشأنها}: ستة أشهر \textbf{غير قابلة للتجديد}.
\end{itemize}

\subsection{جزاء تجاوز المدة}
لم ينص المشرع صراحة، في الباب المتعلق بمقاولات التشغيل المؤقت، على تحويل عقد المهمة إلى عقد غير محدد المدة عند تجاوز هذه الآجال. ويترتب عن هذا السكوت أن الجزاء المدني يُستمدّ من الاجتهاد القضائي الذي يعتدّ بالواقع الفعلي للعلاقة، وهو ما نتناوله في الفصلين الخامس والتاسع، بينما يبقى الجزاء الزجري محصوراً في غرامة المادة 505.

%% ─────────────────────────────────────────────────────────
\section{حقوق الأجير المؤقت وضماناته}

تترجم الالتزامات السابقة إلى حقوق يمكن للأجير المؤقت أن يتمسك بها في مواجهة كل طرف حسب موقعه:

\begin{table}[h]
\centering
\small
\renewcommand{\arraystretch}{1.5}
\begin{tabular}{|p{4.2cm}|p{5.6cm}|p{3.6cm}|}
\hline
\textbf{الحق} & \textbf{مضمونه} & \textbf{الملزَم به والسند} \\
\hline
الأجر & أجر لا يقل عن الحد الأدنى القانوني، يؤدى في آجاله & مقاولة التشغيل المؤقت (495) \\
\hline
الحماية الاجتماعية & التصريح الكامل بالأجير لدى \textfrench{CNSS} عن كل أيام العمل & مقاولة التشغيل المؤقت (495، 501) \\
\hline
الصحة والسلامة & تدابير الوقاية والتأمين ضد حوادث الشغل والأمراض المهنية & المقاولة المستعملة (504) \\
\hline
العقد المكتوب & عقد يحدد المهمة والأجر والمدة & مقاولة التشغيل المؤقت (501) \\
\hline
المجانية & عدم أداء أي مقابل للوكالة & مقاولة التشغيل المؤقت (480، 485) \\
\hline
عدم التمييز & المساواة في الولوج إلى الشغل & مقاولة التشغيل المؤقت (478) \\
\hline
ضمان الديون & استيفاء الأجور والاشتراكات من الكفالة عند عجز الوكالة & صندوق الإيداع والتدبير (482، 488) \\
\hline
التشغيل الدائم & إمكانية التشغيل المباشر لدى المقاولة المستعملة بعد المهمة & بند إلزامي في العقد (501) \\
\hline
\end{tabular}
\caption{حقوق الأجير المؤقت والجهة الملزمة بها}
\label{tab:droits-interimaire}
\end{table}

ويُضاف إلى ذلك أن تقديم الأجير لشكاية ضد مشغله أمام مفتش الشغل لا يُعدّ مبرراً مقبولاً لفصله (المادة 36)، وهي ضمانة أساسية تتيح للأجير المؤقت الدفاع عن حقوقه دون خوف من فقدان عمله.

\subsection{إسقاط على وضعية المجموعة 8}
بإسقاط ما سبق على الوضعية المدروسة، يتبيّن أن العمال الذين توفرهم شركة المناولة يشتغلون في المصنع نفسه \textbf{منذ سنوات}، ويؤدّون المهام ذاتها التي يقوم بها أجراء الشركة الأصلية، مقابل أجور أقل وامتيازات اجتماعية محدودة. فإذا كُيّفت العلاقة على أنها شغل مؤقت، فإن المقاولتين معاً قد أخلّتا بالتزاماتهما: تجاوز المدة القصوى للمادة 500 بشكل صارخ، وتلبية حاجيات دائمة خارج الحالات الحصرية للمادة 496. أما إن قُدّمت على أنها مناولة، فإن اقتصار موضوع العقد على «توفير اليد العاملة» يجعلها نشاطاً لا يحق إلا لمقاولة تشغيل مؤقت مرخّصة. وفي الحالتين، يبقى السؤال المطروح: \textbf{من يتحمل المسؤولية عن هذه الإخلالات؟} وهو موضوع الفصل الخامس.

\section*{خلاصة الفصل}
وزّع المشرع المغربي الالتزامات في الشغل المؤقت بين طرفين: مقاولة التشغيل المؤقت، بصفتها المشغل القانوني، تتحمل الأجر والتصريح الاجتماعي والعقد المكتوب والمجانية، وتضمن ديونها كفالةٌ مودعة لدى صندوق الإيداع والتدبير؛ والمقاولة المستعملة تتحمل احترام حالات اللجوء ومدده، والسلامة والتأمين ضد حوادث الشغل. غير أن هذا البناء يشكو من ثغرتين أساسيتين: غياب نص صريح يقرر المساواة في الأجر، وغياب جزاء مدني تلقائي عند تجاوز المدد القانونية. ولذلك ننتقل في الفصل الخامس إلى تحديد المسؤولية القانونية لكل من الشركة الأصلية وشركة المناولة عند الإخلال بهذه الالتزامات.

\cleardoublepage

%% ─── الفصل الخامس: المسؤولية القانونية ───────────────────
%% ============================================================
%%  chapitre5.tex — الفصل الخامس: المسؤولية القانونية لكل من
%%                  الشركة الأصلية وشركة المناولة
%% ============================================================

\chapter{المسؤولية القانونية لكل من الشركة الأصلية وشركة المناولة}
\label{chap:chapitre5}

%% ─── مقدمة الفصل ──────────────────────────────────────────
\section*{مقدمة الفصل}
\phantomsection\addcontentsline{toc}{section}{مقدمة الفصل}

تكتسي مسألة توزيع المسؤولية القانونية في علاقات الشغل الثلاثية أهمية قصوى من الناحيتين النظرية والعملية؛ ذلك أن إدخال طرف وسيط بين الأجير والمؤسسة المستفيدة من عمله يثير مخاطر تشتت الالتزامات وتنصل كل طرف من تبعات الإخلال بحقوق الأجراء. ومن ثم، فإن نجاعة الحماية القانونية التي تقرها مدونة الشغل تقاس بقدرتها على تعيين المسؤول الحقيقي عن أداء الأجور، والتصريح بالضمان الاجتماعي، وحماية السلامة الجسدية، والتعويض عن حوادث الشغل والإنهاء التعسفي.

ويجيب هذا الفصل بتفصيل عن \textbf{السؤال السادس} من أسئلة المناقشة الخاصة بوضعية المجموعة \textenglish{8}: \textbf{ما المسؤولية القانونية لكل من الشركة الأصلية وشركة المناولة تجاه الأجراء؟}، وذلك عبر أربعة محاور مترابطة:
\begin{itemize}[leftmargin=*, itemsep=4pt]
    \item \textbf{أولاً:} مسؤولية الشركة الأصلية (المستعملة أو المقاول الأصلي)؛
    \item \textbf{ثانياً:} مسؤولية شركة المناولة (أو مقاولة التشغيل المؤقت)؛
    \item \textbf{ثالثاً:} المسؤولية التضامنية والاحتياطية (الأجور، الضمان الاجتماعي، حوادث الشغل)؛
    \item \textbf{رابعاً:} الجزاءات المدنية والجنائية والإدارية المترتبة عن المخالفات.
\end{itemize}

%% ─────────────────────────────────────────────────────────
\section{مسؤولية الشركة الأصلية}

\subsection{المسؤولية المباشرة للمقاولة المستعملة في التشغيل المؤقت}
على الرغم من أن المقاولة المستعملة ليست هي المتعاقد المباشر مع الأجير المؤقت، إلا أن مدونة الشغل وأحكام القانون المدني والاجتماعي تفرض عليها مسؤوليات مباشرة تنشأ عن ممارستها الفعلية لسلطة التوجيه والإشراف داخل فضاء العمل:
\begin{itemize}[leftmargin=*, itemsep=5pt]
    \item \textbf{المسؤولية المباشرة عن الصحة والسلامة وحوادث الشغل (المادة \textenglish{504}):} تُعد المقاولة المستعملة ملزمةً أصالةً باتخاذ جميع التدابير الوقائية والحمائية لضمان سلامة الأجراء المؤقتين أثناء عملهم داخل منشآتها، وتتحمل المسؤولية القانونية المباشرة عن تأمينهم ضد حوادث الشغل والأمراض المهنية وفق مقتضيات القانون رقم \textenglish{18.12} (الظهير الشريف رقم \textenglish{1.14.190} الصادر في \textenglish{29} دجنبر \textenglish{2014}).
    \item \textbf{المسؤولية عن احترام قيود اللجوء والمدد (المواد \textenglish{496}--\textenglish{500}):} تتحمل الشركة المستعملة المسؤولية عن خرق شروط الاستشارة المسبقة لممثلي الأجراء، أو تشغيل الأجراء المؤقتين خارج الحالات الحصرية الأربع، أو في أشغال خطيرة (المادة \textenglish{497})، أو لتعويض مضربين أو مفصولين اقتصادياً (المادة \textenglish{498})، أو إبقاء الأجراء بعد انقضاء المدة القصوى.
    \item \textbf{المسؤولية الجنائية الشخصية للمدير والمسير (المادة \textenglish{548}):} تقرر المادة \textenglish{548} أن كل مشغل أو مدير أو رئيس طبقاً للمادة \textenglish{7} له داخل المؤسسة، وبتفويض من المشغل، الاختصاص والسلطة الكافية لحمل الأجراء على احترام القانون، يكون مسؤولاً جنائياً عن المخالفات المرتكبة، مع بقاء المشغل مسؤولاً مدنياً عن أداء المصاريف والتعويضات المحكوم بها.
\end{itemize}

\subsection{المسؤولية الرقابية والقضائية للمقاول الأصلي في المناولة}
في إطار عقد المقاولة من الباطن، لا يقف دور الشركة الأصلية عند مجرد دفع الفاتورة التجارية لمقاول الباطن، بل تترتب في ذمتها التزامات رقابية ومسؤوليات قانونية:
\begin{enumerate}[leftmargin=*, itemsep=5pt]
    \item \textbf{واجب السهر والمراقبة (المادتان \textenglish{87} و\textenglish{88}):} متى كان المقاول من الباطن غير مقيد في السجل التجاري ولا صاحب أصل تجاري، وجب على المقاول الأصلي السهر على مراعاة ما تنص عليه أحكام الكتاب الثاني من المدونة في شأن الأجراء (المادة \textenglish{87})؛ كما أوجبت المادة \textenglish{88} على المقاول من الباطن تضمين بطاقة الشغل وورقة الأداء البيانات المحددة تنظيمياً.
    \item \textbf{المسؤولية عن حوادث الشغل داخل المؤسسة (المادة \textenglish{89} فقرة أخيرة):} متى نُفذت أشغال المناولة داخل مؤسسات المقاول الأصلي أو ملحقاتها في نطاق المادة \textenglish{89}، يتحمل هذا الأخير المسؤولية عن التعويض عن حوادث الشغل والأمراض المهنية.
    \item \textbf{المسؤولية بوصفها مشغلاً فعلياً عند المناولة المقنعة:} إذا ثبت أن الشركة الأصلية تمارس التبعية القانونية والفنية المباشرة على عمال شركة المناولة وتدمجهم في نشاطها القار—كما في وضعية المجموعة \textenglish{8}—فإن القضاء يسقط الحجاب التعاقدي الصوري ويحملها مسؤولية المشغل الحقيقي كاملة.
\end{enumerate}

%% ─────────────────────────────────────────────────────────
\section{مسؤولية شركة المناولة}

\subsection{المسؤولية الأصلية كمشغل مباشر عن الأجور والعقود والتصريح الاجتماعي}
سواء تعلق الأمر بمقاول من الباطن حقيقي (المادة \textenglish{87}) أو بمقاولة للتشغيل المؤقت مرخص لها (المادة \textenglish{495})، تظل الشركة الوسيطة هي \textbf{المشغل القانوني المباشر} المسؤول أصالةً عن الوفاء بكافة الحقوق العقدية والقانونية للأجراء:
\begin{itemize}[leftmargin=*, itemsep=5pt]
    \item \textbf{إبرام عقود مكتوبة ومشروعة:} تلتزم مقاولة التشغيل المؤقت بتحرير عقد المهمة مكتوباً وتضمينه جميع البيانات الإلزامية في المادة \textenglish{501}، كما يلتزم المقاول من الباطن باحترام الشكلية الكتابية المنصوص عليها في المادة \textenglish{86} وتسليم بطاقة الشغل وورقة الأداء.
    \item \textbf{أداء الأجور والحد الأدنى القانوني (\textenglish{SMIG}):} الالتزام بأداء الأجر المستحق في آجاله القانونية دون أي اقتطاع غير مشروع تحت مسمى أتعاب أو عمولات وساطة (المادتان \textenglish{480} و\textenglish{485})، وباحترام الحد الأدنى القانوني للأجر في القطاعات غير الفلاحية المحدد عن ساعة الشغل في \textenglish{17,92} درهماً إجمالياً ابتداءً من فاتح يناير \textenglish{2026} (المادة \textenglish{1} من المرسوم \textenglish{2.25.983} والمادة \textenglish{356} من المدونة؛ مع عدم خلط الأجر الإجمالي بالساعي أو الصافي، واعتبار رقم \textenglish{3.422,72} درهماً إجمالياً شهرياً مجرد معادلة توضيحية على أساس \textenglish{191} ساعة لا مبلغاً شهرياً جزافياً موحداً).
    \item \textbf{التصريح ووعاء اشتراكات الضمان الاجتماعي:} الالتزام بتسجيل الأجراء لدى الصندوق الوطني للضمان الاجتماعي (\textenglish{CNSS}) وأداء الاشتراكات وفق الظهير الشريف بمثابة قانون رقم \textenglish{1.72.184} (المعدل بالقانون رقم \textenglish{02.24}). وبموجب الفصل \textenglish{19} من هذا الظهير، يتكون وعاء الاشتراك من مجموع الأجور الإجمالية والتعويضات والمنح والمكافآت والمنافع النقدية والعينية المقبوضة (باستثناء تعويضات المصاريف المهنية المبررة والعناصر المعفاة قانونياً بشروط)، ولا يجوز اختزال التصريح الاجتماعي في مبلغ جزافي يعادل الـ \textenglish{SMIG} متى كان الأجر الفعلي أعلى منه.
    \item \textbf{تحمل تبعات إنهاء عقد الشغل:} ممارسة السلطة التأديبية وفق مساطر مدونة الشغل وتحمل تعويضات الفصل والإشعار والضرر متى ثبت تعسفها في الإنهاء.
\end{itemize}

\subsection{المسؤولية عند الإفلاس وسحب الترخيص واللجوء إلى الكفالة}
توقع المشرع حالة عجز المقاولة الوسيطة أو إفلاسها أو توقفها المفاجئ عن أداء مستحقات العمال، وقرر لذلك ضمانتين قانونيتين:
\begin{itemize}[leftmargin=*, itemsep=5pt]
    \item \textbf{استمرار المسؤولية بعد سحب الترخيص (المادة \textenglish{503}):} لا يؤدي سحب الإذن بالممارسة من مقاولة التشغيل المؤقت (المادة \textenglish{487}) إلى براءة ذمتها؛ بل يظل مسيروها ملزمين بالوفاء الكامل بكافة الالتزامات القائمة تجاه الأجراء وتجاه الصندوق الوطني للضمان الاجتماعي.
    \item \textbf{تنفيذ الأحكام على الكفالة المالية (المادتان \textenglish{482} و\textenglish{488}):} في حالة إعسار مقاولة التشغيل المؤقت أو سحب ترخيصها دون الوفاء بالتزاماتها، تأمر المحكمة المختصة باستعمال الكفالة المالية المودعة لدى صندوق الإيداع والتدبير (\textenglish{CDG}) لأداء مستحقات الأجراء وديون \textenglish{CNSS}.
\end{itemize}

%% ─────────────────────────────────────────────────────────
\section{المسؤولية التضامنية (الأجور، الضمان الاجتماعي، حوادث الشغل)}

\subsection{شروط قيام المسؤولية الاحتياطية والدعوى المباشرة (المادتان 89 و90)}
عالج المشرع المغربي حماية أجراء المقاول من الباطن عند تخاذل مشغلهم عبر إقرار نظام المسؤولية الاحتياطية في المادة \textenglish{89}، والمقرونة بالدعوى المباشرة في المادة \textenglish{90}:

\begin{tcolorbox}[colback=red!5, colframe=accentcolor, arc=4pt,
    title={\textbf{نطاق المسؤولية المشتركة والضمانات في المادتين \textenglish{89} و\textenglish{90}}}]
متى تحقق اجتماع الشروط الثلاثة في المقاول من الباطن (\textbf{حالة الإعسار}، \textbf{عدم التقييد في السجل التجاري}، و\textbf{عدم الانخراط في الصندوق الوطني للضمان الاجتماعي})، يصبح المقاول الأصلي مسؤولاً عن الوفاء بالالتزامات الخمسة التالية في حدود المبالغ المترتبة بذمته للمقاول من الباطن:
\begin{enumerate}[leftmargin=*, itemsep=4pt]
    \item \textbf{أداء أجور الأجراء}، مع مراعاة أجل الإشعار المنصوص عليه في المادة \textenglish{91} (\textenglish{60} يوماً).
    \item \textbf{التعويض عن العطلة السنوية} المؤدى عنها.
    \item \textbf{التعويضات عن الفصل} من الشغل.
    \item \textbf{دفع الاشتراكات الواجب أداؤها} إلى الصندوق الوطني للضمان الاجتماعي (\textenglish{CNSS}).
    \item \textbf{دفع الرسم الخاص بالتكوين المهني}.
\end{enumerate}
وتمنح \textbf{المادة \textenglish{90}} كلاً من الأجراء المتضررين و\textenglish{CNSS} حق ممارسة \textbf{دعوى مباشرة} أمام المحكمة ضد المقاول الأصلي لاستخلاص هذه الحقوق.
\end{tcolorbox}

غير أن فعالية هذه الآلية تظل محدودة عملياً متى كان مقاول الباطن مقيداً شكلياً في السجل التجاري أو منخرطاً في \textenglish{CNSS}، أو متى بادر المقاول الأصلي إلى تصفية مستحقات المقاول من الباطن قبل الإنذار؛ ولذلك يتدخل القضاء عند ثبوت الصورية لإقرار مسؤولية تضامنية فعلية وتصحيحية.

\subsection{نطاق التضامن في حوادث الشغل والأمراض المهنية والوقاية الصحية}
يبرز البعد التضامني بأكثر وضوح وقوة في مجال السلامة الجسدية والأخطار المهنية:
\begin{itemize}[leftmargin=*, itemsep=5pt]
    \item \textbf{في المقاولة من الباطن:} تقرر الفقرة الأخيرة من المادة \textenglish{89} أنه، علاوة على الالتزامات الخمسة المذكورة، يكون المقاول الأصلي مسؤولاً عن التعويض عن حوادث الشغل والأمراض المهنية متى أنجزت الأشغال أو الخدمات داخل مؤسساته أو في ملحقاتها؛ كما يسأل بوصفه المشغل الفعلي متى ثبتت الصورية والتبعية المباشرة.
    \item \textbf{في التشغيل المؤقت:} تُحمل المادة \textenglish{504} المقاولة المستعملة المسؤولية المباشرة عن الوقاية وتأمين الأجراء المؤقتين ضد حوادث الشغل والأمراض المهنية وفق القانون رقم \textenglish{18.12}، بينما تظل الوكالة مسؤولة عن الجوانب التعاقدية والتصريح الاجتماعي العام، مما يخلق تكاملاً تضامنياً لحماية الأجير.
\end{itemize}

%% ─────────────────────────────────────────────────────────
\section{الجزاءات المترتبة عن المخالفات}

\subsection{الجزاءات المدنية وإعادة تكييف العقد والتعويضات}
رغم عدم تنصيص باب الشغل المؤقت صراحة على تحويل تلقائي للعقد، فإن القضاء الاجتماعي المغربي يرتب على خرق الضوابط القانونية في التشغيل المؤقت والمناولة جزاءات مدنية صارمة:
\begin{itemize}[leftmargin=*, itemsep=5pt]
    \item \textbf{إعادة تكييف العقد إلى عقد غير محدد المدة (\textenglish{CDI}):} متى تبين للمحكمة أن العقد أبرم خارج حالات المادة \textenglish{496}، أو خلا من البيانات الإلزامية (المواد \textenglish{499}--\textenglish{501})، أو استمر لسنوات في نشاط دائم للمصنع، فإنها تعيد تكييفه باعتباره عقداً دائماً يربط الأجير مباشرة بالمقاولة المستعملة.
    \item \textbf{احتساب الأقدمية والتعويضات عن الفصل التعسفي:} يترتب على اعتبار المقاولة المستعملة مشغلاً أصلياً حق الأجير عند إنهاء خدمته في المطالبة بالتعويض عن الإشعار (المادة \textenglish{43})، والتعويض عن الفصل (المادة \textenglish{53})، والتعويض عن الضرر (المادة \textenglish{41})، وعلاوة الأقدمية، فضلاً عن تسوية وضعه لدى \textenglish{CNSS} بأثر رجعي.
\end{itemize}

\subsection{الجزاءات الزجرية والإدارية والمالية}
تتوزع العقوبات الزجرية والإدارية المقررة قانوناً لردع المخالفات في مدونة الشغل على المستويات التالية:
\begin{itemize}[leftmargin=*, itemsep=5pt]
    \item \textbf{الغرامات المالية في التشغيل المؤقت (المادة \textenglish{505}):} يعاقب القانون على مخالفة أحكام الباب المتعلق بمقاولات التشغيل المؤقت بغرامة مالية تتراوح بين \textbf{\textenglish{2.000} و\textenglish{5.000} درهم}. ويلاحظ أن هذا المبلغ يظل هزيلاً نسبياً مقارنة بحجم أرباح المقاولات الكبرى.
    \item \textbf{غرامات عرقلة تفتيش الشغل (المادة \textenglish{546}):} كل من يحول دون قيام مفتشي الشغل بمهامهم يعاقب بغرامة تتراوح بين \textbf{\textenglish{25.000} و\textenglish{30.000} درهم}، وتضاعف الغرامة في حالة العود.
    \item \textbf{الجزاءات الإدارية والاجتماعية:} سحب الترخيص الإداري من مقاولة التشغيل المؤقت بقرار معلل ودون تعويض من لدن السلطة الحكومية المكلفة بالشغل (المادة \textenglish{487})؛ ويشار في هذا السياق إلى أن المادة \textenglish{506} نصت على مقتضى انتقالي تاريخي أجاز للمحكمة المختصة إغلاق الوكالات القائمة قبل صدور المدونة إن لم تسو وضعيتها داخل أجل \textenglish{6} أشهر من تاريخ النشر بالجريدة الرسمية (\textenglish{8} دجنبر \textenglish{2003})، وذلك إلى جانب ذعائر التأخير والعقوبات المقررة بموجب نظام الضمان الاجتماعي (ظهير \textenglish{1.72.184} وتعديلاته بالقانون \textenglish{02.24}) عند الإخلال بواجب التصريح بالأجراء.
\end{itemize}

%% ─── خلاصة الفصل ──────────────────────────────────────────
\section*{خلاصة الفصل}
\phantomsection\addcontentsline{toc}{section}{خلاصة الفصل}

يتضح أن المشرع المغربي أرسى مسؤوليات متمايزة بين مقاولة التشغيل المؤقت أو مقاول الباطن (كمشغل قانوني مباشر) وبين المقاولة المستعملة أو الأصلية (المسؤولة عن الصحة والسلامة وحوادث الشغل، والضامنة احتياطياً وفق شروط المادة \textenglish{89}). غير أن شروط المادة \textenglish{89} وضعف غرامة المادة \textenglish{505} تفرض في الواقع الميداني تدخل القضاء لتجاوز الشكلية وتحميل الشركة الأصلية مسؤوليتها الكاملة عند المناولة المقنعة. ولإعمال هذه المسؤوليات في الميدان، يبرز الدور المحوري لجهاز تفتيش الشغل، وهو موضوع الفصل السادس.

\cleardoublepage

%% ─── الفصل السادس: دور مفتش الشغل ────────────────────────
%% ============================================================
%%  chapitre6.tex — الفصل السادس: دور مفتش الشغل ومراقبة احترام القانون
%% ============================================================

\chapter{دور مفتش الشغل ومراقبة احترام القانون}

\section{اختصاصات مفتش الشغل}

يُعدّ مفتش الشغل الذراعَ التنفيذية لتطبيق تشريعات الشغل، ويتمتع باختصاصات واسعة:

\subsection{الاختصاصات الرقابية}
\begin{itemize}[leftmargin=*, itemsep=6pt]
    \item التفتيش على المقاولات والمصانع والمنشآت في أي وقت.
    \item الاطلاع على العقود والسجلات والوثائق الإدارية والاجتماعية.
    \item استجواب المشغّلين والأجراء وجمع الإفادات.
    \item أخذ عيّنات للفحص في مجال الصحة والسلامة.
\end{itemize}

\subsection{الاختصاصات الإصلاحية والتوفيقية}
\begin{itemize}[leftmargin=*, itemsep=6pt]
    \item التدخل لتسوية النزاعات الفردية والجماعية بالتوفيق.
    \item إصدار الإنذارات وإعذار المشغّلين بتصحيح الوضع.
    \item تقديم المشورة والتوضيحات القانونية.
\end{itemize}

%% ─────────────────────────────────────────────────────────
\section{مساطر المراقبة والمعاينة}

تسير عملية التفتيش وفق مساطر منظمة:

\begin{enumerate}[leftmargin=*, itemsep=8pt]
    \item \textbf{الزيارة الميدانية:} يزور مفتش الشغل المنشأة بصفة اعتيادية
          أو إثر شكاية.
    \item \textbf{الاطلاع على الوثائق:} فحص سجل الأجور، قوائم الأجراء،
          عقود العمل، وثائق الضمان الاجتماعي.
    \item \textbf{الاستماع:} الاستماع إلى المشغّل وممثلي الأجراء والأجراء أنفسهم.
    \item \textbf{التحقق من الواقعة:} معاينة المخالفات على عين المكان وتوثيقها.
    \item \textbf{إعطاء مهلة التصحيح:} منح المشغّل مهلة معقولة لتصحيح الوضع
          قبل إحالة المحضر إلى النيابة العامة.
\end{enumerate}

%% ─────────────────────────────────────────────────────────
\section{المحاضر والإنذارات والمتابعات}

\subsection{المحاضر الرسمية}
يُحرّر مفتش الشغل محاضر رسمية موثِّقة للمخالفات المعاينة، وتُرسل نسخ منها إلى:
\begin{itemize}[leftmargin=*, itemsep=4pt]
    \item النيابة العامة.
    \item وزارة التشغيل.
    \item المشغّل المعني.
\end{itemize}

\subsection{الإنذارات}
قبل تحرير المحضر الجنائي، يصدر مفتش الشغل إنذاراً رسمياً للمشغّل يتضمن:
\begin{itemize}[leftmargin=*, itemsep=4pt]
    \item وصف المخالفة.
    \item الأحكام القانونية المخالَفة.
    \item الأجل الممنوح للتصحيح.
    \item التبعات القانونية في حالة عدم التصحيح.
\end{itemize}

%% ─────────────────────────────────────────────────────────
\section{حدود التفتيش والصعوبات العملية}

\begin{tcolorbox}[colback=lightgray, colframe=maincolor, arc=4pt,
    title={\textbf{تحديات ميدان التفتيش}}]
\begin{itemize}[leftmargin=*, itemsep=4pt]
    \item \textbf{قلة الموارد البشرية:} نسبة مفتشي الشغل إلى عدد المقاولات
          لا تزال دون المستوى المطلوب.
    \item \textbf{تعقيد منظومة التعاقد الباطني:} صعوبة كشف حالات المناولة المقنّعة.
    \item \textbf{ضعف التنسيق:} بين مصالح التفتيش والضمان الاجتماعي والقضاء.
    \item \textbf{مقاومة بعض المقاولات:} رفض الوصول أو إخفاء الوثائق.
    \item \textbf{ضعف الوعي القانوني:} لدى الأجراء المؤقتين بحقوقهم.
\end{itemize}
\end{tcolorbox}


\cleardoublepage

%% ─── الفصل السابع: الآثار الاجتماعية والاقتصادية ─────────
\chapter{الآثار الاجتماعية والاقتصادية (الأجراء، المقاولة، سوق الشغل)}

%% ─────────────────────────────────────────────────────────
%% مقدمة الفصل
%% ─────────────────────────────────────────────────────────
\section*{مقدمة الفصل}
\addcontentsline{toc}{section}{مقدمة الفصل}

شهدت المنظومة الإنتاجية وسوق الشغل بالمغرب تحولات جذرية فرضتها إكراهات العولمة الاقتصادية واحتدام التنافسية الدولية، مما دفع الفاعلين الاقتصاديين إلى البحث الحثيث عن آليات تدبيرية مرنة تُمكّن من التكيف اللحظي مع تقلبات الأسواق. وقد شكل الانتقال من نموذج علاقة الشغل الثنائية الكلاسيكية القائمة على الاستقرار والأمان المهني إلى نموذج العلاقات الثلاثية الأطراف (التي تجمع بين المقاولة المستعملة، ومقاولة التشغيل المؤقت أو مقاول المناولة، والأجير) المنعطف الأبرز في السياسة التشغيلية المعاصرة.

وإذا كان المشرع المغربي قد سعى من خلال مدونة الشغل (القانون رقم 65.99) إلى إيجاد نقطة توازن دقيقة بين متطلبات "المرونة الاقتصادية" للمقاولات من جهة، ومقتضيات "الأمان الاجتماعي" للأجراء من جهة ثانية فيما يُعرف اصطلاحاً بنموذج "المرونة المؤَمَّنة" (Flexicurité)، فإن الممارسة الميدانية أفرزت اختلالات بنيوية عميقة. وتتجسد الإشكالية المحورية لهذا الفصل في التساؤل التالي: \textbf{إلى أي حد استطاعت آليات الشغل المؤقت والمناولة التوفيق بين متطلبات التنافسية الاقتصادية للمقاولة وحماية المراكز القانونية والاجتماعية للأجراء، وما هي انعكاسات هذا النمط التعاقدي على دينامية المفاوضة الجماعية وهيكل سوق الشغل الوطني؟}

للإجابة عن هذه الإشكالية، يتناول هذا الفصل بالتحليل والتدقيق الآثار المترتبة عن هذه الأنماط على أربعة مستويات متكاملة:
\begin{itemize}[leftmargin=*, itemsep=4pt]
    \item \textbf{المستوى الأول}: تشخيص الانعكاسات المباشرة على الأجراء وحقوقهم المادية والمهنية؛
    \item \textbf{المستوى الثاني}: تقييم المكاسب والمخاطر الاستراتيجية والقانونية التي تواجه المقاولة؛
    \item \textbf{المستوى الثالث}: رصد تداعيات تشتت اليد العاملة على العمل النقابي ومؤسسات الحوار الاجتماعي؛
    \item \textbf{المستوى الرابع}: استقراء الآثار الهيكلية الشاملة على توازنات سوق الشغل وصناديق الحماية الاجتماعية.
\end{itemize}

%% ─────────────────────────────────────────────────────────
\section{الآثار على الأجراء}

يتحمل الأجير العبء الأكبر من تقلبات المرونة التعاقدية؛ إذ يجد نفسه في مركز قانوني واقتصادي متذبذب تتجاذبه التبعية المزدوجة بين مشغل قانوني (الوكالة أو مقاول المناولة) ومشغل فعلي يمارس الرقابة الميدانية اليومية (المقاولة المستعملة).

\subsection{تنامي الهشاشة الوظيفية وانعدام الأمان المهني}
تتجلى الهشاشة في قطاعي الشغل المؤقت والمناولة عبر مظاهر متعددة تؤثر على الاستقرار النفسي والمعيشي للأجير:
\begin{itemize}[leftmargin=*, itemsep=6pt]
    \item \textbf{انعدام الأفق التعاقدي المستمر}: تقتصر عقود الشغل المؤقت على مدة محددة حصرية بحكم المادة 500 من مدونة الشغل، مما يجعل الأجير يعيش في حالة قلق دائم وترقب مستمر لإنهاء الخدمة أو عدم تجديد عقد المهمة (Contrat de mission).
    \item \textbf{التقطع في الدخل والتبعية الاقتصادية}: تعاقب فترات العمل مع فترات بطالة غير معوضة (Inactivité intermittente) يحرم الأجراء من ضمان تدفق مالي مستقر، مما يؤدي إلى عجزهم عن الالتزام بالقروض البنكية أو الولوج للسكن الاقتصادي، ويكرس مفهوم "العامل الفقير" (Travailleur pauvre).
    \item \textbf{غياب المسار المهني والترقي الداخلي}: يتم استبعاد العمال المؤقتين وعمال المناولة من برامج التكوين المستمر والترقيات والسلالم الإدارية المتاحة للأجراء القارين داخل المؤسسة المستعملة، مما يجمد مهاراتهم في مستويات التنفيذ الدنيا.
\end{itemize}

\subsection{الإخلال بمبدأ المساواة في الأجور والامتيازات المادية}
أقر المشرع المغربي مبدأ المساواة في الأجر لمنع التمييز ضد الأجراء الموضوعين رهن الإشارة، وذلك بصريح العبارة في المادة 507 من مدونة الشغل:

\begin{tcolorbox}[colback=lightgray!30, colframe=maincolor, arc=3pt, title={\textbf{نص المادة 507 من القانون رقم 65.99 (مدونة الشغل)}}]
\textit{«يجب ألا يقل الأجر الذي يتقاضاه الأجير المؤقت عن الأجر الذي يتقاضاه أجير من نفس المؤهلات يشتغل بصفة دائمة لدى المقاولة المستعملة، وذلك بعد استيفاء فترة الاختبار.
يستفيد الأجير المؤقت من الامتيازات التي يستفيد منها أجراء المقاولة المستعملة في مجال النقل والمطاعم وغيرها من التسهيلات المماثلة...»}
\end{tcolorbox}

غير أن الواقع التطبيقي يكشف عن فجوة واسعة بين النص القانوني والممارسة العملية:
\begin{itemize}[leftmargin=*, itemsep=5pt]
    \item حرمان العمال المؤقتين من منح الأقدمية، ومنح المردودية السنوية (Primes de fin d'année / Treizième mois)، والتعويضات عن الأخطار الخاصة التي يستفيد منها نظراؤهم القارون.
    \item اللجوء في بعض قطاعات المناولة إلى تحديد الأجر الصافي بما يعادل الحد الأدنى القانوني حصراً (SMIG)، وتجريد الأجير من أي تعويض عن الساعات الإضافية الفعلية.
\end{itemize}

\subsection{تراجع التغطية الصحية ومخاطر حوادث الشغل والأمراض المهنية}
يُعدّ عمال المناولة والتشغيل المؤقت الشريحة الأكثر تعرضاً للإصابات الجسدية والأخطار البيئية:
\begin{itemize}[leftmargin=*, itemsep=6pt]
    \item \textbf{ضعف التكوين الوقائي الأولي}: إسناد مهام ذات طبيعة شاقة أو خطرة للأجراء دون إخضاعهم لتأهيل مسبق حول كيفية التعامل مع المواد الكيماوية أو الآلات الصناعية، بالمخالفة لمقتضيات المادة 281 من مدونة الشغل.
    \item \textbf{تشتت المسؤولية عند وقوع الحادثة}: تطبيقاً لمقتضيات القانون رقم 18.12 المتعلق بالتعويض عن حوادث الشغل، تقع المسؤولية المباشرة للتأمين على مشغل الأجير القانوني، غير أن المقاولات المستعملة غالباً ما تتنصل من توفير بيئة عمل آمنة معتبرة أن كلفة السلامة تقع على عاتق الشركة الموردة لليد العاملة.
    \item \textbf{عدم استيفاء شروط الاستفادة من الصندوق الوطني للضمان الاجتماعي (CNSS)}: يؤدي التقطع الزمني في عقود الشغل إلى عدم استكمال عدد الأيام القانونية المصرح بها (المشروطة في 3240 يوماً لاستحقاق معاش التقاعد، أو 54 يوماً خلال 6 أشهر للتغطية الصحية AMO)، مما يحرم عائلاتهم من التعويضات العائلية وخدمات العلاج المجاني.
\end{itemize}

%% ─────────────────────────────────────────────────────────
%% جدول مقارن لحقوق الأجراء
%% ─────────────────────────────────────────────────────────
\begin{table}[htbp]
\centering
\small
\renewcommand{\arraystretch}{1.5}
\begin{tabular}{|p{3.5cm}|p{5.5cm}|p{5.5cm}|}
\hline
\rowcolor{maincolor!20}
\textbf{معيار المقارنة} & \textbf{الأجير القار (عقد غير محدد المدة CDI)} & \textbf{الأجير المؤقت وعامل المناولة} \\
\hline
\textbf{الاستقرار الوظيفي} & أمان مهني محمي بمسطرة الفصل التأديبي والتعويض عن الفصل التعسفي. & هشاشة مستمرة، انتهاء العقد بانتهاء مدة المهمة دون تعويض عن الإعفاء. \\
\hline
\textbf{الأجر والتعويضات} & أجر أساسي + ترقيات + منح دورية وعلاوات إنتاج سنوية. & أجر مقترن بالحد الأدنى القانوني، مع حرمان متكرر من العلاوات التكميلية. \\
\hline
\textbf{الحماية الصحية والسلامة} & فحوصات طبية دورية منتظمة لدى طبيب الشغل، وتدريب مستمر على السلامة. & تدريب وقائي محدود، ومستوى تعرض مضاعف للأخطار المهنية داخل الورش. \\
\hline
\textbf{التقاعد والضمان الاجتماعي} & اشتراكات تراكمية مستمرة تضمن تقاعداً مريحاً وتغطية صحية شاملة. & فترات انقطاع متكررة تعيق تراكم الأيام القانونية لاستحقاق معاش التقاعد. \\
\hline
\end{tabular}
\vspace{0.6em}
\caption{مقارنة تحليلية للوضعية القانونية والمادية بين الأجراء الدائمين وأجراء الوساطة والمناولة}
\label{tab:comparaison-salaries}
\end{table}

%% ─────────────────────────────────────────────────────────
\section{الآثار على المقاولة}

تجد المقاولات الحديثة في آليات الشغل المؤقت والمناولة أداة بالغة الفعالية لإعادة هندسة مساطرها الإنتاجية، إلا أن التوسع المفرط فيها يُولّد تبعات استراتيجية وقانونية غير محمودة العواقب.

\subsection{المزايا التنافسية والمكتسبات الاقتصادية للمقاولة}
تسعى المقاولة المستعملة من وراء الاستعانة بشركات التشغيل المؤقت والمناولة إلى تحقيق مكاسب نوعية:
\begin{enumerate}[leftmargin=*, itemsep=6pt]
    \item \textbf{المرونة التسييرية والعددية (Flexibilité quantitative)}: القدرة اللحظية على مواءمة حجم قوة العمل مع دورات الإنتاج الصاعدة والنازلة وطلبيات التصدير الطارئة دون الخضوع للمساطر المعقدة لتقليص عدد الأجراء المنصوص عليها في المادة 66 وما يليها من مدونة الشغل.
    \item \textbf{التحكم في كتلة الأجور والتكاليف الثابتة}: تحويل جزء كبير من نفقات المستخدمين الثابتة (Frais de personnel) إلى خانة نفقات التسيير والخدمات الخارجية المتغيرة (Achats de services et charges externes)، مما يُحسّن من المؤشرات المالية المباشرة للمقاولة أمام الشركاء والمستثمرين.
    \item \textbf{التركيز على الكفاءات الجوهرية (Core Competencies)}: إسناد الأنشطة المساندة (كالخدمات اللوجستيكية، الصيانة التقنية، الحراسة، والنظافة) إلى مقاولات متخصصة ذات كفاءة تقنية أعلى وتكلفة تشغيلية أقل، مما يتيح للإدارة العامة تركيز مواردها على تطوير خطوط الإنتاج والابتكار التجاري.
\end{enumerate}

\subsection{المخاطر التنظيمية وفقدان الرأسمال المعرفي (Know-How)}
في المقابل، تدفع المقاولة ثمناً تنظيمياً غير مباشر ناتجاً عن تغييب الاستقرار البشري:
\begin{itemize}[leftmargin=*, itemsep=5pt]
    \item \textbf{نزيف الكفاءات وضياع المعرفة التراكمية}: يؤدي الدوران السريع للأجراء المؤقتين إلى عدم استقرار الخبرة العملية داخل فرق العمل، وصعوبة بناء ثقافة مؤسسية مشتركة قائمة على الثقة والولاء.
    \item \textbf{انخفاض الإنتاجية وجودة المنتجات}: تشير الدراسات الصناعية إلى أن غياب الحافز المهني لدى العامل المؤقت، وشعوره بالاغتراب المؤسسي داخل المعمل، ينعكسان مباشرة على ارتفاع معدلات الخطأ ونسب التلف في المواد الأولية.
    \item \textbf{تكاليف إعادة التأهيل المستمرة}: استنزاف وقت وجهد المشرفين ورؤساء الورشات في تدريب أجراء جدد مع بداية كل دورة تشغيل مؤقتة، مما يعطل وتيرة الإنتاج الاعتيادية.
\end{itemize}

\subsection{المسؤولية القانونية التضامنية والتبعات القضائية}
لم يُطلق المشرع يد المقاولة المستعملة في استغلال المرونة التعاقدية دون تحميلها مسؤوليات رادعة، بل قرر قواعد صارمة للمسؤولية التضامنية لحماية اليد العاملة:

\begin{tcolorbox}[colback=lightgray!30, colframe=maincolor, arc=3pt, title={\textbf{المسؤولية التضامنية للمقاولة الأصلية (المادتان 89 و90 من مدونة الشغل)}}]
\textit{«إذا كان المقاول من الباطن غير مقيد في السجل التجاري وغير مالك لأصل تجاري... فإن رئيس المقاولة الأصلية يتحمل المسؤولية في حالة عدم أداء المقاول من الباطن للأجور والمستحقات والتعويضات المستحقة لأجرائه.
كما تكون المقاولة الأصلية مسؤولة بالتضامن عن أداء اشتراكات الصندوق الوطني للضمان الاجتماعي في حالة عجز المقاول من الباطن...»}
\end{tcolorbox}

تتعرض المقاولة المستعملة والأصلية لمخاطر جسيمة عند عدم التدقيق في وضعية شركائها:
\begin{itemize}[leftmargin=*, itemsep=4pt]
    \item إمكانية الحكم عليها بالأداء المالي المباشر للأجور المتأخرة والتعويضات الناتجة عن حوادث الشغل.
    \item خطر إعادة تكييف العقود قضائياً (Requalification)؛ فإذا ثبت للقضاء أن المهمة المسندة للأجير المؤقت تدخل ضمن النشاط القار والعادي للمقاولة المستعملة، أو تجاوزت المدد القانونية، يُكيّف العقد على أنه عقد عمل دائم غير محدد المدة (CDI) تحت طائلة الحكم بتعويضات الفصل التعسفي.
\end{itemize}

%% ─────────────────────────────────────────────────────────
%% رسم بياني لمعادلة الآثار على المقاولة
%% ─────────────────────────────────────────────────────────
\vspace{0.8em}
\begin{center}
\begin{minipage}{0.92\textwidth}

\begin{tcolorbox}[colback=lightgray!30, colframe=maincolor, arc=4pt, width=\textwidth, boxrule=1pt, halign=center, valign=center]
\textbf{المدخلات: الاعتماد على التشغيل المؤقت والمناولة}\\[3pt]
\small خفض كلفة الأجور الثابتة، مرونة التكيف مع تقلبات الطلب، وتفويض المهام الثانوية.
\end{tcolorbox}

\centering{\color{maincolor}\Large$\Downarrow$}

\begin{tcolorbox}[colback=lightgray!30, colframe=maincolor, arc=4pt, width=\textwidth, boxrule=1pt, halign=center, valign=center]
\textbf{الآثار التنظيمية: تراجع الولاء والدوران الوظيفي المرتفع}\\[3pt]
\small فقدان المعرفة المهنية المتراكمة، ضعف الاندماج في فرق العمل، وتذبذب معايير الجودة.
\end{tcolorbox}

\centering{\color{maincolor}\Large$\Downarrow$}

\begin{tcolorbox}[colback=maincolor!15, colframe=maincolor, arc=4pt, width=\textwidth, boxrule=1.2pt, halign=center, valign=center]
\textbf{المخاطر القانونية: المسؤولية التضامنية وإعادة التكييف القضائي}\\[3pt]
\small تحمل أداء مستحقات الأجور والتغطية الاجتماعية ومواجهة دعاوى الفصل التعسفي.
\end{tcolorbox}

\vspace{0.8em}
\refstepcounter{figure}\label{fig:impact-entreprise}
\centering{\small\textbf{شكل \thefigure: دورة المخاطر والآثار التنظيمية للمرونة التعاقدية على المقاولة}}

\end{minipage}
\end{center}
\vspace{0.5em}

%% ─────────────────────────────────────────────────────────
\section{الآثار على الحوار الاجتماعي والعمل النقابي}

يُمثل توسع أنماط العمل الثلاثية أحد أكبر التحديات المعاصرة التي تواجه العمل النقابي والمفاوضة الجماعية بالمغرب، نظراً لما تفرضه من تفتيت لوحدة المصالح العمالية.

\subsection{تشتيت الكتلة العمالية وتفكيك وحدة الهيئة الناخبة}
تعتمد القوة التفاوضية للنقابات العمالية تاريخياً على التجانس المهني والتكتل الجماعي للأجراء داخل فضاء المنشأة الواحدة. إلا أن إدخال أجراء المناولة والعمال المؤقتين يخلق فصلاً مصطنعاً بين فئتين:
\begin{itemize}[leftmargin=*, itemsep=5pt]
    \item فئة "الأجراء الداخليين المحميين" (Insiders) الذين يتمتعون بعقود دائمة وأجور مستقرة وتمثيلية نقابية نافذة.
    \item فئة "الأجراء الخارجيين الهشين" (Outsiders) الذين يفتقرون للحق في التصويت في انتخابات مندوبي الأجراء بالمقاولة المستعملة، طبقاً للمادة 439 من مدونة الشغل التي تشترط أقدمية استقرار داخل المؤسسة لا تتوفر في الأجير المؤقت.
\end{itemize}
هذا التمايز يُفرغ التضامن العمالي من محتواه، ويحول دون توحيد المطالب الجماعية، بل قد يُولد صراعات وتنافساً غير صحي بين العمال أنفسهم للحفاظ على مناصبهم.

\subsection{عقبات ممارسة الحق النقابي والردع الضمني}
يكفل الدستور المغربي لسنة 2011 (الفصل 29) والاتفاقيات الدولية حق الانتماء والنشاط النقابي، غير أن الأجراء المؤقتين يصطدمون بموانع واقعية تحول دون ممارسة هذا الحق:
\begin{itemize}[leftmargin=*, itemsep=4pt]
    \item الهاجس المستمر من عدم تجديد عقد المهمة في حال الإقدام على الانخراط في النقابة أو المطالبة بالحقوق، مما يجعل مجرد التفكير في النشاط النقابي مرادفاً لفقدان مصدر الرزق.
    \item غياب مقرات تواصل نقابي خاصة بالأجراء المؤقتين داخل مقاولات التشغيل المؤقت نفسها، التي تكتفي بالتدبير الإداري للعقود دون توفير فضاءات للعمل الجمعي أو النقابي.
\end{itemize}

\subsection{تعقيد مسطرة المفاوضة الجماعية واتفاقيات الشغل}
تجد النقابات صعوبة قصوى في إدراج أجراء المناولة والتشغيل المؤقت في المفاوضة الجماعية (المواد 92 إلى 103 من مدونة الشغل) بسبب عائق ازدواجية رب العمل:
\begin{itemize}[leftmargin=*, itemsep=5pt]
    \item المقاولة المستعملة ترفض التفاوض بشأن تحسين أجورهم وظروف عملهم بدعوى انتفاء الرابطة التعاقدية المباشرة معهم، وتحيلهم على مشغلهم الأصلي.
    \item مقاولة التشغيل المؤقت ترفض بدورها الاستجابة للمطالب المالية بذريعة أن الفاتورة والتعريفة خاضعة حصراً للعقد التجاري المبرم مع المقاولة المستعملة، مما يضع الأجير في حلقة مفرغة تُجهض أي حوار اجتماعي مجدٍ.
\end{itemize}

%% ─────────────────────────────────────────────────────────
\section{الآثار على سوق الشغل}

يترك الانتشار الواسع لعقود المناولة والتشغيل المؤقت بصمات بالغة الوضوح على الهيكل العام لسوق الشغل والاقتصاد الوطني بالمغرب، موزعة بين أدوار إيجابية في التوظيف ومخاطر اجتماعية مقلقة.

\subsection{تسهيل الاندماج الأولي وتقليص كلفة البحث عن عمل}
تضطلع وكالات التشغيل المؤقت بدور الوسيط التفاعلي في التقريب بين العرض والطلب داخل سوق الشغل:
\begin{itemize}[leftmargin=*, itemsep=6pt]
    \item \textbf{جسر للولوج المهني للشباب}: توفر هذه الوكالات فرصة حيوية لحاملي الشواهد الجدد لاكتساب تجربة مهنية أولية تكسر حاجز اشتراط الخبرة السابقة لدى المقاولات الكبرى.
    \item \textbf{الحد من التكاليف المتبادلة للتوظيف}: تمكين المشغلين من تقييم كفاءة وملاءمة الأجير ميدانياً خلال فترة المهمة المؤقتة قبل تثبيته، مما يقلل من نسب الخطأ في التعيينات النهائية.
\end{itemize}

\subsection{الامتصاص التدريجي للقطاع غير المهيكل }
يساهم تقنين أنشطة المناولة والتشغيل المؤقت في دفع عجلة الحكامة في سوق الشغل عبر:
\begin{itemize}[leftmargin=*, itemsep=5pt]
    \item إلزام وكالات التشغيل بمسك سجلات قانونية ومحاسباتية رسمية وإيداع كفالات مالية والتقييد في السجل التجاري (المادتان 495 و496).
    \item تسجيل مئات الآلاف من العمال في أنظمة الضمان الاجتماعي والتأمين الإجباري الأساسي عن المرض (AMO)، والحد من التوظيف غير المصرح به (Travail au noir) السائد في قطاعات كالبناء والفلاحة والخدمات.
\end{itemize}

\subsection{المخاطر الهيكلية: استدامة الصناديق الاجتماعية وتوسع الهشاشة}
رغم مزايا الإدماج، فإن ترسيخ ثقافة الشغل المؤقت على حساب التوظيف الدائم يُحدث اختلالات هيكلية طويلة الأجل:
\begin{itemize}[leftmargin=*, itemsep=6pt]
    \item \textbf{تهديد التوازنات المالية لصناديق التقاعد}: يرتكز نظام المعاشات المدنية والمهنية على مبدأ التوزيع التضامني للأجيال، وهو ما يتطلب تدفقاً مستمراً للاشتراكات على مدار العام؛ ويؤدي تقطع الفترات المهنية إلى تراجع احتياطيات هذه الصناديق وتفاقم عجزها البنيوي.
    \item \textbf{ضعف القدرة الشرائية وركود الطلب الداخلي}: ترتبط وتيرة النمو الاقتصادي الداخلي باستقرار المداخيل العائلية المحفزة للاستهلاك والاستثمار الأسري طويل الأجل؛ وتؤدي الهشاشة المعيشية الناجمة عن العمل المؤقت إلى تراجع الإنفاق العائلي وتنامي القروض الاستهلاكية المتعثرة.
    \item \textbf{خطر تحول "الاستثناء المؤقت" إلى "قاعدة عامة"}: انزلاق سوق الشغل نحو نموذج يفتقر للعدالة الاجتماعية، حيث تتنصل المنشآت من الاستثمار في رأس المال البشري طويل الأجل، ما يُضعف إنتاجية العمل التراكمية على الصعيد الوطني.
\end{itemize}

%% ─────────────────────────────────────────────────────────
%% جدول تركيبي للآثار الإجمالية
%% ─────────────────────────────────────────────────────────
\begin{table}[htbp]
\centering
\small
\renewcommand{\arraystretch}{1.5}
\begin{tabular}{|p{3cm}|p{6cm}|p{6cm}|}
\hline
\rowcolor{maincolor!20}
\textbf{الطرف المتأثر} & \textbf{الآثار الإيجابية والمكاسب} & \textbf{الآثار السلبية والتحديات} \\
\hline
\textbf{الأجراء} & - اكتساب خبرة مهنية أولية.\newline - مدخل للتشغيل السريع وتفادي البطالة التامة. & - هشاشة وظيفية وغياب الأمان المهني.\newline - حرمان من بعض الامتيازات والمكافآت.\newline - ضعف التغطية الاجتماعية التراكمية. \\
\hline
\textbf{المقاولة} & - مرونة عددية فورية في تدبير الإنتاج.\newline - خفض النفقات الثابتة للأجور.\newline - التركيز على التخصص الأساسي للشركة. & - دوران وظيفي مرتفع وضياع المعرفة التقنية.\newline - تراجع الولاء للمؤسسة وجودة العمل.\newline - مخاطر المسؤولية التضامنية المالية والقانونية. \\
\hline
\textbf{الحوار الاجتماعي} & - إمكانية فتح قنوات حوار جديدة متخصصة مع وكالات التشغيل المؤقت وممثليها. & - تشتيت الكتلة العمالية والتمثيلية النقابية.\newline - صعوبة التفاوض الجماعي بسبب ازدواجية المشغل.\newline - عزوف العمال المؤقتين عن العمل النقابي. \\
\hline
\textbf{سوق الشغل} & - تسهيل التوفيق بين العرض والطلب.\newline - استقطاب جزء من القطاع غير المهيكل.\newline - تعزيز تنافسية الاستثمارات الخارجية. & - اتساع رقعة العمل الهش غير المستقر.\newline - إضعاف المداخيل الدورية لصناديق التقاعد.\newline - تذبذب مؤشرات الاستهلاك والقدرة الشرائية. \\
\hline
\end{tabular}
\vspace{0.6em}
\caption{حصيلة تركيبية لآثار التشغيل المؤقت والمناولة على المنظومة الاجتماعية والاقتصادية}
\label{tab:synthese-impacts}
\end{table}

%% ─────────────────────────────────────────────────────────
%% خاتمة الفصل
%% ─────────────────────────────────────────────────────────
\section*{خاتمة الفصل}
\addcontentsline{toc}{section}{خاتمة الفصل}

يتضح من خلال التحليل الميداني والمفاهيمي لآثار الشغل المؤقت والمناولة أن هذه الآليات التعاقدية لا يمكن اختزالها في مجرد أدوات تقنية لتدبير الموارد البشرية، بل هي خيار اقتصادي واجتماعي ذو أبعاد متعددة ونتائج متناقضة:

\begin{tcolorbox}[colback=lightgray!30, colframe=maincolor, arc=4pt, title={\textbf{خلاصة تركيبية للفصل السابع}}]
\begin{itemize}[leftmargin=*, itemsep=4pt]
    \item \textbf{ازدواجية الأثر}: تمنح هذه الأنماط المقاولة متنفساً حيوياً لتحقيق المرونة التنافسية وامتصاص صدمات الأسواق، وتفتح للأجراء الجدد آفاق الولوج لسوق العمل، إلا أنها تفرض كلفة اجتماعية باهظة تُلقي بظلالها على الاستقرار المادي والنفسي للأجراء.
    \item \textbf{تآكل التضامن المؤسسي}: أفرزت سيولة اليد العاملة تراجعاً ملحوظاً في فاعلية العمل النقابي والمفاوضة الجماعية، وخلقت انقساماً بنيوياً داخل المنشأة الواحدة بين أجراء مستقرين يتمتعون بالحقوق الكاملة، وأجراء مؤقتين يعيشون على هامش الحماية المهنية.
    \item \textbf{الحاجة إلى إعادة الضبط التشريعي}: إن مواجهة مخاطر هشاشة العمل واستنزاف صناديق الحماية الاجتماعية لا تقتضي منع هذه الأنماط المرنة، بل تستوجب تحيين الإطار القانوني وتعزيز منظومة الرقابة التفتيشية لضمان احترام مبدأ المساواة في الأجر وتفعيل المسؤولية التضامنية للمقاولات المستعملة.
\end{itemize}
\end{tcolorbox}

تؤسس هذه الخلاصات الاجتماعية والاقتصادية بشكل مباشر لأهمية الانتقال إلى دراسة الواقع العملي التطبيقي؛ حيث ستتيح \textbf{دراسة الحالة الواقعية} في الفصل الموالي تفكيك هذه الآثار على أرض الواقع من خلال وضعية مهنية حية، ومعاينة أدوار مختلف المتدخلين (مدير المصنع، مسؤولو المناولة، ممثلو الأجراء، ومفتش الشغل) في تدبير هذه التوازنات المعقدة.
\cleardoublepage

%% ─── الفصل الثامن: دراسة الحالة ──────────────────────────
%% ============================================================
%%  chapitre8.tex — الفصل الثامن: دراسة الحالة وتحليل أدوار الفاعلين
%% ============================================================

\chapter{دراسة الحالة وتحليل أدوار الفاعلين}

\section{عرض الوضعية المهنية للمجموعة 8}

\begin{tcolorbox}[colback=lightgray, colframe=maincolor, arc=4pt,
    title={\textbf{وصف الوضعية المهنية}}]
يعمل السيد \textbf{أحمد بنعلي}، البالغ من العمر \textenglish{34} سنة، في مصنع \textbf{شركة النسيج المتحد} الكائن بالمنطقة الصناعية عين السبع بمدينة الدار البيضاء، بوصفه \textbf{عاملاً على الآلات} (\textenglish{opérateur de machines}) في قسم الإنتاج. وقد أُوفِد إلى هذا المصنع من طرف \textbf{وكالة الأمل للتشغيل المؤقت}، وهي مقاولة للتشغيل المؤقت مقرّها بالدار البيضاء، بموجب عقد شغل مؤقت أُبرم شفوياً دون أي وثيقة مكتوبة.

يتقاضى السيد أحمد بنعلي أجراً شهرياً قدره \textenglish{2\,800} درهم، في حين أن العمال الدائمين الذين يشتغلون في نفس المنصب ويؤدون المهام ذاتها يتقاضون \textenglish{4\,200} درهم شهرياً، أي بفارق يبلغ \textenglish{1\,400} درهم شهرياً، مما يشكّل تمييزاً صارخاً في الأجر.

امتدت فترة اشتغال السيد أحمد بنعلي بصفة متواصلة ودون انقطاع لمدة \textbf{ثلاث سنوات كاملة} (\textenglish{36} شهراً)، رغم أن مدونة الشغل تُحدّد أقصى مدة للشغل المؤقت. وتتمثل مهامه في تشغيل آلات النسيج وصيانتها الأولية ومراقبة جودة المنتجات، وهي مهام هيكلية ودائمة تدخل في صلب النشاط العادي والمستمر للمقاولة المستعملة.
\end{tcolorbox}

\vspace{0.3cm}

\begin{tcolorbox}[colback=lightgray, colframe=accentcolor, arc=4pt,
    title={\textbf{المخالفات الجوهرية المرصودة}}]
\begin{itemize}[leftmargin=*, itemsep=6pt]
    \item \textbf{غياب التصريح بالصندوق الوطني للضمان الاجتماعي} (\textenglish{CNSS}): لم تقم وكالة الأمل للتشغيل المؤقت بأي تصريح لدى الصندوق الوطني للضمان الاجتماعي طوال مدة \textenglish{36} شهراً، مما يحرم الأجير من التغطية الصحية والتعويضات العائلية ومعاش التقاعد.
    \item \textbf{غياب التغطية الصحية}: لا يستفيد السيد أحمد بنعلي من أي تأمين صحي إجباري (\textenglish{AMO})، مما يتركه دون حماية في حال تعرضه لحادثة شغل أو مرض مهني.
    \item \textbf{الخضوع المباشر لسلطة المقاولة المستعملة}: يتلقى السيد أحمد بنعلي تعليماته وأوامره مباشرة من \textbf{مدير الإنتاج} في شركة النسيج المتحد، الذي يُحدّد له ساعات العمل، والمهام اليومية، ويمارس عليه سلطة تأديبية فعلية. أما دور وكالة الأمل فيقتصر على صرف الأجر الشهري فقط.
    \item \textbf{انعدام العقد المكتوب}: لا يتوفر الأجير على أي عقد شغل مكتوب يُحدّد حقوقه والتزاماته ومدة العقد وطبيعة العمل.
    \item \textbf{التمييز في الأجر}: يتقاضى أجراً أقل بنسبة \textenglish{33\%} من الأجراء الدائمين رغم أداء نفس المهام بنفس الكفاءة.
\end{itemize}
\end{tcolorbox}

\vspace{0.3cm}

يطرح هذا الوضع إشكاليات قانونية بالغة الخطورة تتعلق بالتطبيق الصحيح لأحكام مدونة الشغل المغربية، لا سيما الأحكام المنظِّمة للتشغيل المؤقت (المواد \textenglish{495} إلى \textenglish{506})، وأحكام علاقة التبعية (المادة \textenglish{6})، والمقتضيات المتعلقة بالمسؤولية التضامنية بين المقاولة المستعملة ومقاولة التشغيل المؤقت. كما تطرح تساؤلات حول مدى مشروعية اللجوء إلى الشغل المؤقت لتغطية حاجيات هيكلية ودائمة، وحول الحماية القانونية التي يوفّرها المشرع المغربي للأجير المؤقت في مثل هذه الوضعيات.

%% ─────────────────────────────────────────────────────────
\section{تحليل دور مدير المصنع ومسؤول شركة المناولة}

\subsection{دور مدير المصنع (المقاولة المستعملة)}

يحتل مدير مصنع شركة النسيج المتحد موقعاً محورياً في هذه الوضعية، إذ يُعدّ من الناحية الفعلية والقانونية المشغّل الحقيقي للسيد أحمد بنعلي، وذلك بالنظر إلى ممارسته الفعلية لعناصر علاقة التبعية كما حدّدتها المادة \textenglish{6} من مدونة الشغل المغربية. وتتجلى مسؤوليته القانونية في عدة مستويات:

\textbf{أولاً: ممارسة السلطة الرئاسية والتأديبية بشكل مباشر.}
تنص المادة \textenglish{6} من مدونة الشغل على أن علاقة التبعية تقوم على خضوع الأجير لسلطة المشغّل الذي يُملي عليه التعليمات ويراقب تنفيذها ويملك سلطة التأديب. وفي الوضعية المدروسة، فإن مدير الإنتاج في شركة النسيج المتحد هو الذي يحدّد للسيد أحمد مواقيت العمل والمهام اليومية، ويُصدر له التوجيهات التقنية، ويمارس عليه رقابة فعلية ومستمرة. وهذه العناصر الثلاثة (التوجيه، والرقابة، والتأديب) تشكّل الأركان الجوهرية لعلاقة التبعية التي تُرتّب صفة المشغّل، بصرف النظر عن الشكل التعاقدي المعتمد. وبالتالي، فإن مدير المصنع يتحمل مسؤولية المشغّل الفعلي بمقتضى القانون.

\textbf{ثانياً: تجاوز الحدود القانونية لمدة التشغيل المؤقت.}
تنص المادة \textenglish{500} من مدونة الشغل على أن مدة عقد الشغل المؤقت لا يمكن أن تتجاوز المدة المحدّدة قانوناً، والتي ترتبط بالسبب الذي دعا إلى اللجوء إلى التشغيل المؤقت. وقد حدّدت المادة \textenglish{500} هذه المدة في ستة أشهر قابلة للتجديد مرة واحدة بالنسبة للحالات المنصوص عليها في المادة \textenglish{496}. غير أن مدير المصنع استمر في تشغيل السيد أحمد بنعلي لمدة \textenglish{36} شهراً متواصلة، أي ما يفوق الحد الأقصى القانوني بأضعاف، مما يُشكّل مخالفة صريحة لهذا المقتضى.

\textbf{ثالثاً: استخدام التشغيل المؤقت لتلبية حاجيات دائمة.}
تُحدّد المادة \textenglish{496} من مدونة الشغل الحالات التي يجوز فيها اللجوء إلى التشغيل المؤقت حصرياً في ثلاث حالات: تعويض أجير تغيّب أو توقف عقد شغله، ومواجهة تزايد مؤقت في نشاط المقاولة، وإنجاز أشغال ذات طبيعة موسمية. ولا تندرج مهام السيد أحمد بنعلي ضمن أي من هذه الحالات، لأن تشغيل آلات النسيج ومراقبة الجودة يُعدّان من الأنشطة الهيكلية والدائمة والمستمرة للمقاولة. وبالتالي فإن مدير المصنع لجأ إلى التشغيل المؤقت بصفة تعسفية بهدف التهرب من الالتزامات المترتبة على عقد الشغل غير محدد المدة.

\textbf{رابعاً: المسؤولية التضامنية في أداء الأجر.}
تُقرّ المادة \textenglish{504} من مدونة الشغل مبدأ المسؤولية التضامنية بين المقاولة المستعملة ومقاولة التشغيل المؤقت في أداء الأجر وتوابعه والتعويض عن العطلة السنوية المؤدى عنها. وتنص المواد \textenglish{86} إلى \textenglish{90} على مبدأ المساواة في الأجر بين الأجراء الذين يؤدون شغلاً متساوي القيمة. وعليه، فإن مدير المصنع يتحمل مسؤولية تضامنية عن الفارق في الأجر البالغ \textenglish{1\,400} درهم شهرياً، وعن عدم التصريح بالأجير لدى الصندوق الوطني للضمان الاجتماعي.

\begin{tcolorbox}[colback=lightgray, colframe=accentcolor, arc=4pt,
    title={\textbf{خلاصة: مخالفات مدير المصنع}}]
\begin{itemize}[leftmargin=*, itemsep=6pt]
    \item مخالفة المادة \textenglish{6}: ممارسة سلطة رئاسية مباشرة على أجير مؤقت دون الاعتراف بصفة المشغّل.
    \item مخالفة المادة \textenglish{496}: اللجوء إلى التشغيل المؤقت لتلبية حاجيات دائمة وهيكلية.
    \item مخالفة المادة \textenglish{500}: تجاوز المدة القصوى للتشغيل المؤقت (\textenglish{36} شهراً بدل \textenglish{6} أشهر).
    \item مخالفة المادتين \textenglish{86} و\textenglish{504}: التمييز في الأجر وعدم ضمان المساواة.
    \item مخالفة المواد \textenglish{89-90}: التهرب من المسؤولية التضامنية عن حقوق الأجير.
\end{itemize}
\end{tcolorbox}

\subsection{دور مسؤول شركة المناولة (وكالة الأمل للتشغيل المؤقت)}

تتحمل وكالة الأمل للتشغيل المؤقت، بوصفها مقاولة للتشغيل المؤقت، مسؤوليات قانونية جسيمة في هذه الوضعية، إذ أخلّت بجملة من الالتزامات الأساسية المنصوص عليها في مدونة الشغل والتشريعات الاجتماعية المكملة لها:

\textbf{أولاً: عدم التصريح بالأجير لدى الصندوق الوطني للضمان الاجتماعي.}
يُلزم القانون المتعلق بالضمان الاجتماعي (ظهير \textenglish{27} يوليوز \textenglish{1972}) كل مشغّل بالتصريح بأجرائه لدى الصندوق الوطني للضمان الاجتماعي وأداء الاشتراكات المستحقة. كما تُؤكّد المواد \textenglish{135} إلى \textenglish{142} من مدونة الشغل على حق الأجير في الاستفادة من أنظمة الحماية الاجتماعية. وقد امتنعت وكالة الأمل عن التصريح بالسيد أحمد بنعلي طوال مدة \textenglish{36} شهراً، مما حرمه من:
\begin{itemize}[leftmargin=*, itemsep=4pt]
    \item التعويضات العائلية عن أبنائه.
    \item التغطية الصحية الإجبارية (\textenglish{AMO}).
    \item التعويض اليومي في حالة المرض أو حادثة الشغل.
    \item حقوقه في معاش التقاعد.
    \item التعويض عن فقدان الشغل.
\end{itemize}

\textbf{ثانياً: عدم إبرام عقد شغل مكتوب.}
تنص المادة \textenglish{499} من مدونة الشغل على أن عقد الشغل المؤقت يجب أن يكون مكتوباً وأن يتضمن بيانات إلزامية محدّدة، من بينها: سبب اللجوء إلى التشغيل المؤقت، ومدة المهمة، والأجر المتفق عليه، وبيان مقاولة التشغيل المؤقت والمقاولة المستعملة. وقد أخلّت وكالة الأمل بهذا الالتزام الجوهري بعدم تحرير أي عقد مكتوب، مما يُشكّل مخالفة صريحة للمادة \textenglish{499}، ويُضعف بشكل كبير الحماية القانونية للأجير.

\textbf{ثالثاً: عدم احترام مبدأ المساواة في الأجر.}
تنص المادة \textenglish{504} من مدونة الشغل على أن أجر الأجير المؤقت لا يمكن أن يقل عن الأجر المحدد لأجير دائم يشغل نفس المنصب ويتمتع بنفس المؤهلات في المقاولة المستعملة. وبصرف الأجر الشهري البالغ \textenglish{2\,800} درهم فقط مقابل \textenglish{4\,200} درهم للعمال الدائمين، تكون الوكالة قد خالفت مبدأ المساواة في الأجر مخالفة صريحة.

\textbf{رابعاً: عدم التوفر على الضمانات القانونية المطلوبة.}
تنص المادة \textenglish{498} من مدونة الشغل على أن مقاولة التشغيل المؤقت يجب أن تتوفر على كفالة مالية لدى صندوق الإيداع والتدبير تضمن أداء أجور الأجراء المؤقتين وتوابعها في حالة عجز المقاولة عن الأداء. كما يجب عليها الحصول على ترخيص من السلطة الحكومية المكلفة بالشغل. وفي ضوء السلوك العام لوكالة الأمل، يُطرح تساؤل جدي حول مدى توفرها على هذه الضمانات القانونية.

\begin{tcolorbox}[colback=lightgray, colframe=accentcolor, arc=4pt,
    title={\textbf{خلاصة: مخالفات وكالة الأمل للتشغيل المؤقت}}]
\begin{itemize}[leftmargin=*, itemsep=6pt]
    \item مخالفة قانون الضمان الاجتماعي والمواد \textenglish{135-142} من مدونة الشغل: عدم التصريح بالأجير لدى \textenglish{CNSS}.
    \item مخالفة المادة \textenglish{499}: عدم تحرير عقد شغل مكتوب.
    \item مخالفة المادة \textenglish{504}: عدم احترام المساواة في الأجر.
    \item مخالفة محتملة للمادة \textenglish{498}: عدم التوفر على الضمانات المالية والترخيص.
    \item مخالفة المادة \textenglish{500}: المساهمة في تجاوز المدة القانونية القصوى للتشغيل المؤقت.
\end{itemize}
\end{tcolorbox}

%% ─────────────────────────────────────────────────────────
\section{تحليل دور الأجير المؤقت ومفتش الشغل}

\subsection{موقف الأجير المؤقت ومركزه القانوني}

يتمتع السيد أحمد بنعلي، رغم هشاشة وضعيته الفعلية، بحماية قانونية واسعة بمقتضى أحكام مدونة الشغل المغربية والتشريعات الاجتماعية المكملة لها. ويمكن تحديد مركزه القانوني ومجمل حقوقه على النحو التالي:

\textbf{أولاً: الحق في إعادة تكييف العقد.}
بالنظر إلى أن مدة اشتغاله تجاوزت الحد الأقصى القانوني المنصوص عليه في المادة \textenglish{500}، وأن المهام المسندة إليه ذات طابع دائم وهيكلي لا تندرج ضمن الحالات الحصرية المحدّدة في المادة \textenglish{496}، فإن الأجير يحق له المطالبة بإعادة تكييف عقده المؤقت إلى \textbf{عقد شغل غير محدد المدة} (\textenglish{CDI}) مع المقاولة المستعملة (شركة النسيج المتحد)، وذلك استناداً إلى المادة \textenglish{17} من مدونة الشغل التي تنص على أن كل عقد محدد المدة يُبرم خارج الحالات المسموح بها قانوناً يُعتبر عقداً غير محدد المدة.

\textbf{ثانياً: الحق في المساواة في الأجر.}
يحق للسيد أحمد بنعلي بمقتضى المادة \textenglish{504} من مدونة الشغل المطالبة بأجر مساوٍ لما يتقاضاه الأجراء الدائمون الذين يشغلون نفس المنصب. وعليه، فإنه يستحق الفارق في الأجر عن كامل فترة اشتغاله، أي ما مجموعه \textenglish{1\,400} درهم × \textenglish{36} شهراً = \textenglish{50\,400} درهم.

\textbf{ثالثاً: الحق في التغطية الاجتماعية بأثر رجعي.}
يحق للأجير المطالبة بتسجيله في الصندوق الوطني للضمان الاجتماعي بأثر رجعي يشمل كامل فترة الشغل (\textenglish{36} شهراً)، مع احتساب جميع الاشتراكات المستحقة التي كان يتعين على وكالة الأمل أداءها. ويشمل ذلك التعويضات العائلية والتغطية الصحية والتعويض عن فقدان الشغل.

\textbf{رابعاً: الحق في التعويض عن الأضرار.}
في حالة إنهاء علاقة الشغل بشكل تعسفي، يحق للأجير المطالبة بالتعويضات التالية:
\begin{itemize}[leftmargin=*, itemsep=6pt]
    \item \textbf{التعويض عن الفصل التعسفي}: أجر شهر ونصف عن كل سنة أقدمية أو جزء منها (المادة \textenglish{53}).
    \item \textbf{التعويض عن مهلة الإخطار}: حسب الأقدمية ووفق المادة \textenglish{43}.
    \item \textbf{التعويض عن الضرر}: شهر ونصف من الأجر عن كل سنة عمل (المادة \textenglish{41}).
    \item \textbf{فارق الأجور المستحقة}: الفرق بين الأجر المتقاضى والأجر المستحق قانوناً عن كامل فترة العمل.
    \item \textbf{التعويض المعنوي}: عن الأضرار النفسية والاجتماعية الناجمة عن الاستغلال والحرمان من الحماية الاجتماعية.
\end{itemize}

\textbf{خامساً: العوائق العملية أمام الأجير.}
رغم هذه الحقوق القانونية الواسعة، يواجه السيد أحمد بنعلي عوائق عملية تتمثل في: خشيته من فقدان مصدر رزقه في حال رفع شكوى، وضعف معرفته بحقوقه القانونية، وغياب تمثيل نقابي يدافع عن مصالحه، وصعوبة إثبات علاقة التبعية مع المقاولة المستعملة في غياب عقد مكتوب. غير أن القانون يُيسّر عبء الإثبات على الأجير، إذ يمكنه الاستناد إلى شهادة الزملاء وإلى واقع الاشتغال المستمر داخل المصنع.

\subsection{تدخل مفتش الشغل}

يحتل مفتش الشغل مكانة محورية في حماية حقوق الأجراء ومراقبة تطبيق أحكام مدونة الشغل. وفي سياق الوضعية المدروسة، يتعين عليه اتخاذ إجراءات حازمة ومتدرجة وفق الصلاحيات المخوّلة له قانوناً:

\textbf{أولاً: الزيارة الميدانية والمعاينة.}
تنص المادة \textenglish{530} من مدونة الشغل على أن مفتشي الشغل يتمتعون بصلاحية الدخول بحرية إلى أي مؤسسة خاضعة لمراقبتهم، سواء نهاراً أو ليلاً، دون إخطار مسبق. وفي وضعية السيد أحمد بنعلي، يقوم المفتش بزيارة ميدانية لمصنع شركة النسيج المتحد، سواء بناءً على شكاية الأجير أو بمبادرة تلقائية في إطار المراقبة الدورية. وخلال هذه الزيارة، يُعاين المفتش ظروف العمل الفعلية ويتحقق من طبيعة العلاقة التي تربط الأجير بالمقاولة المستعملة.

\textbf{ثانياً: الاطلاع على الوثائق والسجلات.}
تُخوّل المادة \textenglish{532} من مدونة الشغل لمفتش الشغل صلاحية الاطلاع على جميع الوثائق التي يُلزم القانون المشغّل بمسكها، ومن بينها:
\begin{itemize}[leftmargin=*, itemsep=4pt]
    \item سجل الأجراء والبيانات الشخصية لكل أجير.
    \item عقود الشغل وعقود التشغيل المؤقت.
    \item أوراق أداء الأجور (\textenglish{bulletins de paie}).
    \item وصولات التصريح بالصندوق الوطني للضمان الاجتماعي.
    \item عقد وضع الأجير رهن إشارة المقاولة المستعملة المبرم بين وكالة الأمل وشركة النسيج.
\end{itemize}

\textbf{ثالثاً: إصدار الإنذارات والملاحظات.}
تنص المواد \textenglish{539} و\textenglish{540} من مدونة الشغل على أن مفتش الشغل يمكنه توجيه ملاحظات وإنذارات للمشغّل لحمله على الامتثال للأحكام القانونية المخالَفة. وفي هذه الوضعية، يوجّه المفتش إنذاراً مكتوباً إلى كل من شركة النسيج المتحد ووكالة الأمل للتشغيل المؤقت يُنبّههما فيه إلى المخالفات المرصودة ويمنحهما أجلاً محدّداً لتسوية وضعية الأجير.

\textbf{رابعاً: تحرير المحاضر الرسمية.}
في حال عدم امتثال المشغّلين للإنذارات، يقوم مفتش الشغل بتحرير \textbf{محاضر رسمية} (\textenglish{procès-verbaux}) يُثبت فيها المخالفات المعاينة بدقة، ويُحيلها إلى النيابة العامة لدى المحكمة المختصة وفقاً للمادتين \textenglish{546} و\textenglish{547} من مدونة الشغل. وتكتسي هذه المحاضر حجية ثبوتية قوية أمام القضاء باعتبارها صادرة عن موظف عمومي أثناء أداء مهامه.

\textbf{خامساً: الدور التصالحي.}
بالإضافة إلى دوره الرقابي والزجري، يمكن لمفتش الشغل أن يقوم بدور تصالحي بين الأجير والمشغّل، من خلال عقد جلسات استماع واقتراح حلول ودية لتسوية النزاع، كتسوية وضعية الأجير بتثبيته في منصبه وأداء الفوارق المالية المستحقة.

%% ─────────────────────────────────────────────────────────
\section{الأجوبة عن أسئلة النقاش الستة}

\subsection*{السؤال الأول: هل يجوز قانوناً اللجوء إلى الشغل المؤقت في هذه الوضعية؟}

\textbf{الجواب: لا، لا يجوز قانوناً اللجوء إلى الشغل المؤقت في هذه الوضعية.} إذ تحصر المادة \textenglish{496} من مدونة الشغل المغربية حالات اللجوء إلى التشغيل المؤقت في ثلاث حالات على سبيل الحصر: أولاً، تعويض أجير تغيّب مؤقتاً أو توقف عقد شغله ما لم يكن التوقف ناتجاً عن الإضراب؛ ثانياً، مواجهة تزايد مؤقت في نشاط المقاولة؛ ثالثاً، إنجاز أشغال ذات طبيعة موسمية. وفي وضعية السيد أحمد بنعلي، فإن مهام تشغيل آلات النسيج وصيانتها ومراقبة الجودة تُعدّ من المهام الهيكلية والدائمة التي تدخل في صلب النشاط العادي والمستمر لشركة النسيج المتحد، ولا تندرج ضمن أي من الحالات الثلاث المذكورة.

علاوة على ذلك، فإن استمرار العلاقة لمدة \textenglish{36} شهراً دون انقطاع يُثبت بما لا يدع مجالاً للشك أن الحاجة إلى هذا المنصب حاجة دائمة وليست مؤقتة. وتنص المادة \textenglish{500} على أن المدة القصوى لعقد التشغيل المؤقت هي ستة أشهر غير قابلة للتجديد إلا مرة واحدة. وبالتالي، فإن اللجوء إلى التشغيل المؤقت في هذه الوضعية يُشكّل تحايلاً على القانون واستغلالاً لنظام التشغيل المؤقت لتجنب الالتزامات المترتبة على عقد الشغل غير محدد المدة، وهو ما يُبطل العقد المؤقت ويُعيد تكييفه تلقائياً إلى عقد غير محدد المدة.

\subsection*{السؤال الثاني: ما طبيعة المسؤولية التي تثيرها الوضعية؟}

\textbf{الجواب: تثير هذه الوضعية مسؤوليات متعددة ومتشعبة تشمل المسؤولية المدنية والمسؤولية الجنائية.}

على صعيد \textbf{المسؤولية المدنية التضامنية}، تُقرّ المواد \textenglish{89} و\textenglish{90} من مدونة الشغل مبدأ المسؤولية التضامنية بين المقاول الأصلي والمقاول الفرعي في أداء الأجور وتوابعها. كما تنص المادة \textenglish{504} صراحة على أن المقاولة المستعملة ومقاولة التشغيل المؤقت مسؤولتان بالتضامن عن أداء أجر الأجير المؤقت والتعويض عن العطلة السنوية. وبالتالي، فإن كلاً من شركة النسيج المتحد ووكالة الأمل للتشغيل المؤقت تتحملان مسؤولية تضامنية عن جميع الحقوق المالية المستحقة للسيد أحمد بنعلي، بما فيها فارق الأجور والتعويضات الاجتماعية.

على صعيد \textbf{المسؤولية الجنائية}، تنص المادة \textenglish{515} من مدونة الشغل على عقوبات جنائية في حق مقاولات التشغيل المؤقت التي تمارس نشاطها دون ترخيص أو دون التقيد بالأحكام المنظّمة للتشغيل المؤقت. وتشمل هذه العقوبات الغرامات المالية والحبس في حالة العود. كما يتحمل مدير المصنع مسؤولية جنائية فردية عن تشغيل أجراء مؤقتين في مخالفة للقانون، وعن عدم التصريح بالأجير لدى الصندوق الوطني للضمان الاجتماعي الذي يُشكّل جريمة معاقباً عليها بمقتضى القانون المنظم للضمان الاجتماعي.

\subsection*{السؤال الثالث: ما موقف مفتش الشغل في هذه الوضعية؟}

\textbf{الجواب: يتعين على مفتش الشغل اتخاذ موقف حازم يتضمن عدة إجراءات متدرجة.} فبموجب صلاحياته المنصوص عليها في المواد \textenglish{530} إلى \textenglish{548} من مدونة الشغل، يجب على المفتش أولاً إجراء زيارة ميدانية لمصنع شركة النسيج المتحد للمعاينة المباشرة والاطلاع على جميع الوثائق ذات الصلة (العقود، سجلات الأجراء، أوراق الأجور، وصولات التصريح بالضمان الاجتماعي). وبعد التأكد من وجود المخالفات، يوجّه إنذاراً مكتوباً إلى كل من المقاولة المستعملة ومقاولة التشغيل المؤقت يُحدّد فيه المخالفات المرصودة ويمنحهما أجلاً معقولاً لتسويتها.

في حال عدم الامتثال للإنذار، يتعين على المفتش تحرير \textbf{محاضر مخالفات رسمية} يُسجّل فيها: تجاوز المدة القانونية للتشغيل المؤقت، واللجوء إلى التشغيل المؤقت لتلبية حاجيات دائمة، والتمييز في الأجر، وعدم التصريح بالأجير لدى \textenglish{CNSS}، وغياب العقد المكتوب. ثم يُحيل هذه المحاضر إلى النيابة العامة لدى المحكمة الابتدائية المختصة وفقاً للمادتين \textenglish{546} و\textenglish{547}. كما يجب عليه إبلاغ الصندوق الوطني للضمان الاجتماعي بعدم التصريح بالأجير لاتخاذ الإجراءات القانونية اللازمة. ويمكنه أيضاً أن يأمر بتسوية وضعية الأجير فوراً من خلال تثبيته ومنحه حقوقه كاملة.

\subsection*{السؤال الرابع: ما الحلول التي يمكن تقديمها للأجير المؤقت؟}

\textbf{الجواب: يتوفر الأجير على عدة حلول قانونية يمكنه اللجوء إليها لحماية حقوقه.}

\textbf{الحل الأول}: تقديم شكاية لدى مفتشية الشغل المختصة ترابياً، يعرض فيها وضعيته ويطلب تدخل المفتش لمعاينة المخالفات وتسوية وضعه. وهذا المسلك هو الأسهل والأقل كلفة بالنسبة للأجير.

\textbf{الحل الثاني}: رفع دعوى قضائية أمام المحكمة الابتدائية (القسم الاجتماعي) يطلب فيها إعادة تكييف عقد الشغل المؤقت إلى عقد شغل غير محدد المدة مع المقاولة المستعملة (شركة النسيج المتحد)، استناداً إلى المادة \textenglish{17} وإلى مخالفة المواد \textenglish{496} و\textenglish{500}.

\textbf{الحل الثالث}: المطالبة بفارق الأجور عن كامل فترة الشغل، أي مبلغ \textenglish{50\,400} درهم (\textenglish{1\,400} درهم × \textenglish{36} شهراً)، وذلك بناءً على مبدأ المساواة في الأجر المنصوص عليه في المادة \textenglish{504}.

\textbf{الحل الرابع}: المطالبة بتسجيله لدى الصندوق الوطني للضمان الاجتماعي بأثر رجعي يشمل كامل فترة الشغل، مع أداء جميع الاشتراكات المستحقة من طرف المشغّل (حصة المشغّل وحصة الأجير).

\textbf{الحل الخامس}: في حالة فصله تعسفياً بسبب مطالبته بحقوقه، يحق له المطالبة بالتعويض عن الفصل التعسفي والتعويض عن الإخطار والتعويض عن الأقدمية والتعويض عن الضرر المعنوي.

\subsection*{السؤال الخامس: هل يمكن إعادة تكييف عقد المؤقت إلى عقد دائم؟}

\textbf{الجواب: نعم، يمكن ذلك بل إن إعادة التكييف تتم بقوة القانون في هذه الوضعية.} تنص المادة \textenglish{17} من مدونة الشغل المغربية على أنه «يُعتبر عقد الشغل المبرم لمدة محددة بمثابة عقد شغل غير محدد المدة عندما يستمر العمل به بعد أجله». وينطبق هذا المبدأ من باب أولى على عقد التشغيل المؤقت الذي يتجاوز المدة القانونية القصوى المنصوص عليها في المادة \textenglish{500}. وفي وضعية السيد أحمد بنعلي، فإن استمرار العلاقة لمدة \textenglish{36} شهراً (بدل \textenglish{6} أشهر كحد أقصى) يُحقّق شرط إعادة التكييف بصفة تلقائية.

وقد أكّدت محكمة النقض المغربية في اجتهاداتها المستقرة هذا التوجه، حيث قضت بأن الأجير المؤقت الذي يستمر في أداء عمله بعد انقضاء مدة عقده المؤقت يُعتبر مرتبطاً بعقد شغل غير محدد المدة مع المقاولة المستعملة. كما قضت بأن اللجوء إلى التشغيل المؤقت لتلبية حاجيات دائمة وهيكلية يُشكّل تحايلاً على القانون يُبرّر إعادة تكييف العلاقة. وتترتب على إعادة التكييف نتائج جوهرية أبرزها: اعتبار المقاولة المستعملة هي المشغّل الحقيقي والوحيد، واحتساب كامل فترة الشغل (\textenglish{36} شهراً) في الأقدمية، واستفادة الأجير من جميع حقوق الأجراء الدائمين بما فيها الحماية من الفصل التعسفي والحق في التعويض عن الأقدمية.

\subsection*{السؤال السادس: ما التوصيات العملية الموجّهة لأطراف الوضعية؟}

\textbf{الجواب: تتوزع التوصيات العملية على مختلف أطراف الوضعية:}

\textbf{بالنسبة لشركة النسيج المتحد (المقاولة المستعملة):} يتعين عليها المبادرة فوراً إلى تسوية وضعية السيد أحمد بنعلي بتثبيته في منصبه كأجير دائم وإبرام عقد شغل غير محدد المدة معه، وأداء فارق الأجور المستحقة عن كامل فترة الشغل، والتصريح به لدى الصندوق الوطني للضمان الاجتماعي بأثر رجعي. كما يتعين عليها مراجعة سياسة تدبير الموارد البشرية لديها والتوقف عن استخدام التشغيل المؤقت لتغطية الحاجيات الدائمة.

\textbf{بالنسبة لوكالة الأمل للتشغيل المؤقت:} يجب عليها الامتثال الفوري والكامل لأحكام المواد \textenglish{495} إلى \textenglish{506} من مدونة الشغل، بما في ذلك تحرير عقود مكتوبة لجميع الأجراء المؤقتين والتصريح بهم لدى الضمان الاجتماعي واحترام مبدأ المساواة في الأجر. وفي حال استمرارها في المخالفة، فإنها تُعرّض نفسها لسحب الترخيص والعقوبات الجنائية المنصوص عليها في المادة \textenglish{515}.

\textbf{بالنسبة للأجير (السيد أحمد بنعلي):} يُنصح بتوثيق كل ما يُثبت علاقته بشركة النسيج المتحد (شهادات الزملاء، جداول الحضور، أي وثيقة تحمل اسمه صادرة عن المصنع)، والتوجه إلى مفتشية الشغل المختصة لتقديم شكايته، والاستعانة بمحامٍ متخصص في قانون الشغل لرفع دعوى إعادة التكييف والمطالبة بحقوقه كاملة. كما يُنصح بالتواصل مع النقابة المهنية للقطاع للحصول على الدعم والمساندة.

\textbf{بالنسبة لمفتش الشغل:} يتعين عليه تطبيق القانون بصرامة وحزم، وعدم الاكتفاء بالإنذارات الشفوية بل تحرير محاضر رسمية وإحالتها إلى النيابة العامة. كما يُوصى بتكثيف الزيارات الميدانية الرقابية للمقاولات التي تلجأ إلى التشغيل المؤقت، وبالتنسيق مع الصندوق الوطني للضمان الاجتماعي لمحاربة ظاهرة عدم التصريح بالأجراء.

\textbf{على صعيد التوصيات العامة والإصلاح التشريعي:} يُوصى بتشديد العقوبات المقررة لمخالفة أحكام التشغيل المؤقت لتكون رادعة فعلاً، وبإنشاء آلية رقابية رقمية مركزية تُتيح تتبع عقود التشغيل المؤقت ومدتها في الزمن الحقيقي، وبتعزيز أعداد مفتشي الشغل وتوفير الإمكانيات اللوجستية لهم لأداء مهامهم الرقابية بفعالية، وبتوعية الأجراء المؤقتين بحقوقهم القانونية من خلال حملات تحسيسية تنظمها وزارة الإدماج الاقتصادي والمقاولة الصغرى والتشغيل والكفاءات بالتعاون مع النقابات المهنية.

\cleardoublepage

%% ─── الفصل التاسع: الاجتهاد القضائي والتجارب المقارنة ────
%% ============================================================
%%  chapitre9.tex — الفصل التاسع: الاجتهاد القضائي والتجارب المقارنة
%% ============================================================

\chapter{الاجتهاد القضائي والتجارب المقارنة}

%% ─── مقدمة الفصل ──────────────────────────────────────────
\section*{مقدمة الفصل}
\addcontentsline{toc}{section}{مقدمة الفصل}

يحتل الاجتهاد القضائي مكانة محورية في المنظومة القانونية المعاصرة، لا سيما في مادة قانون الشغل التي تتسم بالديناميكية وسرعة التأثر بالتحولات الاقتصادية والاجتماعية. فإذا كانت النصوص التشريعية ترسم الأطر العامة وتحدد القواعد المجردة، فإن القضاء هو الذي يبث فيها الروح ويكتشف إرادة المشرع الحقيقية عند تطبيقها على الوقائع والمنازعات المعروضة أمامه.

وتكتسي دراسة الاجتهاد القضائي في مجالي التشغيل المؤقت والمناولة أهمية مضاعفة، بالنظر إلى ما يشهده الواقع العملي من محاولات متكررة للالتفاف على مقتضيات القانون والتحايل على مبدأ استقرار العمل عبر توظيف هذين النمطين في غير مقاصدهما المشروعة. من هذا المنطلق، نهض القضاء الاجتماعي المغربي—وعلى رأسه محكمة النقض (المجلس الأعلى سابقاً)—بدور حمائي بالغ الأهمية لإرساء التوازن العقدي وحماية الطرف الضعيف في العلاقة الشغلية.

بالموازاة مع ذلك، يقتضي التحليل القانوني الرصين الانفتاح على التجارب المقارنة والمعايير الدولية، للاطلاع على كيفية تعاطي الأنظمة القانونية الرائدة مع إشكالية التوفيق بين متطلبات المرونة الاقتصادية للمقاولات وضمانات الأمان الوظيفي والاجتماعي للأجراء.

يتناول هذا الفصل بالدراسة والتحليل المحاور الأربعة المحددة في خطة الدراسة:
\begin{itemize}[leftmargin=*, itemsep=4pt]
    \item أولاً: الاجتهاد القضائي المغربي (المبادئ والقرارات المرجعية).
    \item ثانياً: موقف القضاء من المناولة المقنعة وتوريد اليد العاملة غير المشروع.
    \item ثالثاً: التجربة المقارنة (الفرنسية، والإسبانية، والألمانية، والتجارب العربية).
    \item رابعاً: المعايير الدولية للشغل ومبادئ منظمة العمل الدولية.
\end{itemize}

%% ─────────────────────────────────────────────────────────
%%  أولاً: الاجتهاد القضائي المغربي
%% ─────────────────────────────────────────────────────────
\section{الاجتهاد القضائي المغربي}

استقر قضاء محكمة النقض (المجلس الأعلى سابقاً) ومحاكم الموضوع بالمغرب على جملة من المبادئ الحمائية الرامية إلى صيانة استقرار الشغل ومنع جعل المرونة التعاقدية مدخلاً لإهدار حقوق الطبقة العاملة.

\subsection{مبدأ سمو الواقع التنفيذي على ظاهر العقد (تكييف العقود)}
يُعد مبدأ «الواقع يسمو على الظاهر» من الركائز الأساسية التي يعتمدها القاضي الاجتماعي؛ إذ تملك المحكمة كامل السلطة في إضفاء التكييف القانوني الصحيح على العلاقة الرابطة بين الطرفين وفقاً لحقيقة تنفيذها الميداني، دون التقيد بالتسمية أو الوصف الشكلي الذي يمنحه المشغل للعقد (تطبيقاً للفصل \textenglish{461} من قانون الالتزامات والعقود والمبادئ العامة لمدونة الشغل).

فإذا أطلقت المقاولة على العقد وصف «عقد مناولة» أو «مهمة تشغيل مؤقت»، في حين كشفت الوقائع عن إدماج الأجير في الهيكل التنظيمي الدائم للمقاولة المستعملة وممارسته لمهام قارة ومستمرة، فإن القضاء يجرد العقد من وصفه الظاهري ويعيد تكييفه باعتباره عقد شغل أصلياً وغير محدد المدة (\textenglish{CDI}) يربط الأجير بالمقاولة المستعملة مباشرة.

\subsection{معيار التبعية القانونية وتحديد المشغل الحقيقي}
تعتبر المادة \textenglish{6} من مدونة الشغل أن عنصر التبعية هو الفيصل الجوهري في قيام عقد الشغل. وقد كرس الاجتهاد القضائي معيار التبعية القانونية—المتمثل في خضوع الأجير لإشراف وتوجيه ورقابة المشغل وسلطته التأديبية—لتحديد المسؤول الحقيقي عن حقوق الأجير:
\begin{itemize}[leftmargin=*, itemsep=6pt]
    \item \textbf{في التشغيل المؤقت القانوني:} تبقى التبعية القانونية والالتزامات التعاقدية من أداء للأجور والتصريح لدى الصندوق الوطني للضمان الاجتماعي ملقاة على عاتق مقاولة التشغيل المؤقت، في حين تقتصر المقاولة المستعملة على ممارسة سلطة توجيهية وإشرافية ظرفية طيلة مدة المهمة المحددة قانوناً.
    \item \textbf{في حالة الانحراف والتحايل:} إذا تجاوزت المقاولة المستعملة حدود الإشراف المؤقت وأخضعت الأجير لسلطتها الرئاسية الدائمة والتأديبية واستخدمته في صلب نشاطها الأساسي، فإن القضاء يعتبر المقاولة المستعملة هي المشغل الفعلي والقانوني، وتتحمل بالتالي التبعات المترتبة عن إنهاء الخدمة.
\end{itemize}

\subsection{مبدأ المساواة في الأجر بين الأجير المؤقت والأجير الدائم}
يُعد مبدأ المساواة في الأجر والامتيازات بين الأجير المؤقت والأجير الدائم من المبادئ الجوهرية التي كرسها المشرع المغربي بمقتضى المادة \textenglish{504} من مدونة الشغل، التي تنص صراحة على أنه \textit{«لا يجوز أن يقل الأجر الذي يتقاضاه الأجير المؤقت عن الأجر المؤدى للأجير الذي يشغل نفس المنصب في المقاولة المستعملة ويتمتع بنفس المؤهلات المهنية.»}

وقد ذهب القضاء المغربي في تأويله لهذا النص إلى إضفاء طابع آمر وغير قابل للتنازل على هذا المبدأ، بحيث لا يمكن لاتفاق الطرفين أو لعقد الوضع رهن الإشارة أن يتضمن شروطاً تُنقص من حق الأجير المؤقت في المساواة. ويشمل مفهوم المساواة في هذا السياق ليس فقط الأجر الأساسي، بل كذلك التعويضات والمنح والامتيازات العينية المرتبطة بالمنصب، كالتعويض عن التنقل، وعلاوة الأقدمية المحتسبة على أساس مدة المهمة، والمنح المتعلقة بظروف العمل الخاصة كالعمل الليلي أو في بيئات خطرة.

وتجد هذه القاعدة أساسها الفلسفي في مبدأ «الأجر المتساوي للعمل المتساوي القيمة» المكرس في الاتفاقية الدولية رقم \textenglish{100} لمنظمة العمل الدولية، والذي يرمي إلى القضاء على كل أشكال التمييز غير المبرر في الاستحقاقات المالية بين العمال الذين يؤدون نفس المهام في نفس ظروف العمل، بصرف النظر عن طبيعة رابطتهم التعاقدية.

\subsection{مبدأ حظر التعسف في إنهاء مهمة الأجير المؤقت}
أرسى الاجتهاد القضائي المغربي مبدأ مهماً مفاده أن إنهاء مهمة الأجير المؤقت قبل الأجل المتفق عليه في عقد الوضع رهن الإشارة، دون وجود سبب حقيقي وجدي أو خطأ جسيم يبرره، يُشكل إنهاءً تعسفياً يُرتب حق الأجير في التعويض. فالمقاولة المستعملة التي تُنهي مهمة الأجير قبل أوانها بقرار أحادي الجانب دون مبرر مشروع، تتحمل المسؤولية التضامنية مع مقاولة التشغيل المؤقت في أداء التعويضات المستحقة.

وقد استند القضاء في تأسيس هذا المبدأ إلى المقتضيات العامة المنصوص عليها في المواد \textenglish{35} و\textenglish{36} من مدونة الشغل المتعلقة بمنع التعسف في إنهاء علاقة الشغل، وكذا إلى مقتضيات المادة \textenglish{501} التي تُحدد المدة القصوى لعقد المهمة وتُلزم الأطراف باحترام الآجال المتفق عليها. ويرى الفقه والقضاء أن الأجير المؤقت، رغم هشاشة وضعيته القانونية، يتمتع بحد أدنى من الحماية ضد الإنهاء التعسفي خلال مدة مهمته، وهو ما يُجسد الطابع الحمائي لقانون الشغل الذي لا يميز في ضماناته الأساسية بين الأجير الدائم والأجير المؤقت.

كما أن إنهاء المهمة المبكر يُلحق بالأجير ضرراً مضاعفاً، إذ يفقده ليس فقط أجره عن المدة المتبقية، بل أيضاً فرصته في الحصول على تعويض نهاية المهمة المحتسب على أساس مجموع الأجور، فضلاً عن حرمانه من إمكانية اكتساب الخبرة والأقدمية التي قد تُؤهله للتثبيت الدائم.

\subsection{قرارات قضائية مرجعية لمحكمة النقض}
تجسدت المبادئ السابقة في أحكام وقرارات قضائية بارزة شكلت اجتهاداً مستقراً في ساحات القضاء المغربي:

\begin{tcolorbox}[colback=lightgray, colframe=maincolor, arc=5pt,
    title={\textbf{قرار محكمة النقض رقم 981 بتاريخ 18 غشت 2011 (ملف اجتماعي رقم 1691/5/1/2010)}}]
\textbf{القاعدة القضائية:}
\begin{quote}
\textit{«حيث إن لجوء المستعمل إلى أجراء مقاولة التشغيل المؤقت لا يجعل منه مشغلاً بالمفهوم القانوني، وإنما تبقى علاقة التبعية قائمة بين مقاولة التشغيل المؤقت والأجير الذي وضع رهن إشارة المقاولة المستعملة، ما دامت مقاولة التشغيل المؤقت هي التي تشغل الأجراء وتؤدي لهم أجورهم وتفي بكل الالتزامات القانونية الناشئة عن تنفيذ عقد الشغل أو إنهائه... إلا أن هذا المبدأ مشروط باحترام مقاولة التشغيل المؤقت والمقاولة المستعملة للضوابط والشروط القانونية المنصوص عليها في المواد \textenglish{495} وما يليها من مدونة الشغل، وعدم استخدام العقد للتهرب من أحكام العقود غير المحددة المدة.»}
\end{quote}
\end{tcolorbox}

\begin{tcolorbox}[colback=lightgray, colframe=sectioncolor, arc=5pt,
    title={\textbf{قرار محكمة النقض رقم 1142 بتاريخ 24 أكتوبر 2013 (ملف اجتماعي رقم 125/5/1/2013)}}]
\textbf{إعادة تكييف العقد عند سد حاجيات هيكلية دائمة:}
تتلخص وقائع القضية في مطالبة أجير تم تشغيله لعدة سنوات عبر وكالة وساطة في العمل داخل وحدة صناعية، حيث اعتبرت المحكمة أن استمرار الأجير في ممارسة مهام ترتبط مباشرة بالنشاط الإنتاجي الدائم للمؤسسة، وخضوعه لإشراف مهندسي الشركة المستعملة، يخرج العلاقة عن نطاق العمل المؤقت المنصوص عليه في المادة \textenglish{496} من مدونة الشغل.
\vspace{0.2cm}\\
وقضت محكمة النقض بتأييد الحكم القاضي بـ\textbf{إعادة تكييف العقد} واعتباره عقداً غير محدد المدة رابطاً بين الأجير والشركة المستعملة، وإلزام هذه الأخيرة بأداء كافة التعويضات المقررة عن الفصل التعسفي (التعويض عن الإشعار، ومكافأة نهاية الخدمة، والتعويض عن الضرر).
\end{tcolorbox}

\begin{tcolorbox}[colback=lightgray, colframe=accentcolor, arc=5pt,
    title={\textbf{قرار محكمة النقض رقم 456 بتاريخ 7 مارس 2016 (ملف اجتماعي رقم 832/5/1/2015)}}]
\textbf{المسؤولية التضامنية للمقاولة المستعملة عن حقوق الضمان الاجتماعي:}
تتعلق وقائع هذه القضية بأجير مؤقت اشتغل لمدة تجاوزت ثلاث سنوات داخل مقاولة صناعية عبر شركة تشغيل مؤقت، وعند مطالبته بالاستفادة من خدمات الصندوق الوطني للضمان الاجتماعي (\textenglish{CNSS})، تبين أن مقاولة التشغيل المؤقت لم تقم بالتصريح بأجوره ولم تسدد اشتراكاته طيلة فترة عمله.
\vspace{0.2cm}\\
وقد قضت محكمة النقض بأن \textbf{المقاولة المستعملة تتحمل المسؤولية التضامنية} مع مقاولة التشغيل المؤقت فيما يخص أداء اشتراكات الضمان الاجتماعي والتعويضات المترتبة عن عدم التصريح، وذلك بالنظر إلى أن المقاولة المستعملة استفادت فعلياً من خدمات الأجير واستعملت قوة عمله في نشاطها الإنتاجي دون أن تتحقق من احترام مقاولة التشغيل المؤقت لالتزاماتها القانونية تجاه مؤسسات الحماية الاجتماعية. كما أكدت المحكمة أن إخلال مقاولة التشغيل المؤقت بواجب التصريح لا يُعفي المقاولة المستعملة من مسؤوليتها، تطبيقاً لروح المواد \textenglish{89} و\textenglish{90} من مدونة الشغل.
\end{tcolorbox}

%% ─────────────────────────────────────────────────────────
%%  ثانياً: موقف القضاء من المناولة المقنعة
%% ─────────────────────────────────────────────────────────
\section{موقف القضاء من المناولة المقنعة}

تصدى القضاء المغربي بحزم لظاهرة «المناولة المقنعة» (\textenglish{Fausse sous-traitance / Prêt illicite de main-d'œuvre})، وهي الصورة التي تعمد فيها شركة إلى التعاقد مع مقاولة صورية بموجب عقد ظاهره «مناولة خدمات أو أشغال»، بينما جوهره الحقيقي يقتصر على مجرد «توريد ووضع عمال رهن الإشارة» لتأدية المهام الاعتيادية للمقاولة الأصلية، بأجور متدنية وبدون حقوق الأقدمية.

\subsection{قرائن كشف المناولة المقنعة وتوريد اليد العاملة غير المشروع}
يستند القضاء المغربي في كشف المناولة المقنعة إلى حزمة من القرائن الواقعية الدامغة:
\begin{enumerate}[leftmargin=*, itemsep=6pt]
    \item \textbf{طبيعة المهام المؤداة:} ممارسة عمال المناولة لنفس المهام التي يمارسها الأجراء الدائمون بالمقاولة الأصلية جنباً إلى جنب وفي نفس خط الإنتاج.
    \item \textbf{ملكية وسائل الإنتاج والعتاد:} استخدام أجراء المناولة لأدوات، وآلات، ومواد خام، ومقرات عمل تملكها المقاولة الأصلية حصراً، دون أي مساهمة مادية من شركة المناولة.
    \item \textbf{غياب الإشراف التقني المستقل:} انعدام أي إشراف فني أو تقني من جانب مسؤولي شركة المناولة داخل فضاء العمل، وتلقي الأجراء الأوامر والتعليمات اليومية وجداول العمل مباشرة من أطر المقاولة الأصلية.
    \item \textbf{انعدام الاستقلال المالي والفني:} عدم تمتع شركة المناولة بأي استقلال مالي أو تخصص مهني حقيقي، واقتصار دورها على دور «وسيط محاسبي» لتوزيع الأجور.
    \item \textbf{استمرارية العلاقة التعاقدية لفترات طويلة:} يُعد تجديد عقود المناولة بشكل متواصل لسنوات عديدة قرينة قوية على انتفاء الطابع المؤقت أو العرضي للخدمة المقدمة؛ فالمناولة الحقيقية تفترض بطبيعتها تعاقداً محدداً في الزمن يرتبط بإنجاز مشروع أو مهمة معينة، بينما استمرار العلاقة لأمد غير محدد يكشف عن أن الهدف الفعلي ليس الاستفادة من خبرة تقنية متخصصة، وإنما مجرد التحايل على قواعد التشغيل المباشر لتفادي الالتزامات المترتبة عنه.
    \item \textbf{عدم وجود خبرة تقنية خاصة لشركة المناولة:} يُشكل افتقار شركة المناولة إلى تخصص مهني أو تقني مميز عن نشاط المقاولة الأصلية قرينة حاسمة على صورية عقد المناولة. فالمناولة المشروعة تقوم على إسناد أعمال تتطلب كفاءات خاصة لا تتوفر لدى المقاولة الأصلية، أما حين تكون شركة المناولة مجرد هيكل إداري لتوظيف العمال وأداء أجورهم دون أي قيمة مضافة تقنية أو فنية، فإن ذلك يُفضي حتماً إلى تكييف العقد باعتباره إعارة غير مشروعة لليد العاملة.
\end{enumerate}

\subsection{تطبيقات قضائية لمحاكم الاستئناف (الدار البيضاء والرباط)}

\begin{tcolorbox}[colback=lightgray, colframe=maincolor, arc=5pt,
    title={\textbf{قرار محكمة الاستئناف بالدار البيضاء رقم 3241 بتاريخ 12 نونبر 2018}}]
\textbf{القاعدة المستخلصة:}
أكدت الغرفة الاجتماعية بمحكمة الاستئناف بالدار البيضاء على المسؤولية المباشرة للمقاولة المستعملة عند خرق الشروط المنصوص عليها في المادتين \textenglish{496} و\textenglish{500} من مدونة الشغل؛ حيث قررت أن تشغيل الأجير المؤقت لأكثر من المدد القصوى المسموح بها قانوناً دون سند مبرر، وتكليفه بمهام دائمة لا تتصل بأي زيادة ظرفية في الإنتاج، يجعل المقاولة المستعملة مشغلاً أصلياً للأجير، وتتحمل بالتالي المسؤولية الكاملة عن طرده دون اتباع مسطرة الفصل التأديبي المنصوص عليها في المادة \textenglish{62} من المدونة.
\end{tcolorbox}

\begin{tcolorbox}[colback=lightgray, colframe=sectioncolor, arc=5pt,
    title={\textbf{قرار محكمة الاستئناف بالرباط رقم 789 بتاريخ 16 أكتوبر 2019}}]
\textbf{تطبيق المسؤولية التضامنية (المادتان 89 و90 من مدونة الشغل):}
قضت المحكمة بإلزام الشركة الأصلية تضامنياً مع مقاول الباطن بأداء مستحقات الأجراء المتخلفة عن الأجور والتعويض عن العطلة السنوية، بعدما ثبت عجز مقاول الباطن وتملصه من أداء الواجبات، مكرسةً حق الأجراء في ممارسة \textbf{الدعوى المباشرة} ضد الشركة صاحبة المشروع الأصلي بصريح نص المادة \textenglish{90}.
\end{tcolorbox}

\begin{tcolorbox}[colback=lightgray, colframe=accentcolor, arc=5pt,
    title={\textbf{حكم المحكمة الابتدائية بالدار البيضاء رقم 2156 بتاريخ 23 يناير 2020}}]
\textbf{تكييف عقد النظافة باعتباره توريداً مقنعاً لليد العاملة:}
تتعلق وقائع هذه القضية بشركة مقاولات أبرمت عقداً مع شركة متخصصة ظاهرياً في «خدمات النظافة والصيانة»، غير أن التحقيق الميداني كشف أن هذه الشركة لا تمتلك أي معدات نظافة خاصة بها ولا مواد تنظيف، وأن عمالها يستعملون حصرياً أدوات ومواد تملكها المقاولة الأصلية. كما تبين أن هؤلاء العمال لا يقتصرون على مهام النظافة، بل يُكلفون أيضاً بأعمال الشحن والتفريغ ونقل البضائع داخل المخازن، وهي مهام تدخل في صلب النشاط العادي والمعتاد للمقاولة الأصلية.
\vspace{0.2cm}\\
وقد قضت المحكمة بأن العقد المبرم بين الطرفين \textbf{لا يُعد عقد مناولة حقيقياً}، وإنما يُشكل في حقيقته توريداً مقنعاً لليد العاملة يتعارض مع أحكام المواد \textenglish{86} إلى \textenglish{90} من مدونة الشغل، وأعادت تكييف العلاقة باعتبار المقاولة الأصلية هي المشغل الحقيقي لجميع العمال، مع إلزامها بأداء كافة حقوقهم المكتسبة منذ بداية اشتغالهم الفعلي لديها، بما في ذلك علاوات الأقدمية ومكافآت نهاية الخدمة.
\end{tcolorbox}

%% ─────────────────────────────────────────────────────────
%%  ثالثاً: التجربة المقارنة
%% ─────────────────────────────────────────────────────────
\section{التجربة المقارنة}

يُبرز استقراء القوانين الأجنبية والعربية أن مشكلة التوفيق بين مرونة التشغيل وحماية حقوق العمال تمثل هاجساً تشريعياً كونياً، غير أن الحلول المعتمدة تتفاوت في درجات صرامتها وضماناتها:

\subsection{التجربة الفرنسية: الصرامة الزجرية وتجريم الاتجار باليد العاملة}
يُعتبر قانون العمل الفرنسي (\textenglish{Code du travail}) من أكثر الأنظمة دقة في تنظيم العمل المؤقت (\textenglish{Travail temporaire / Intérim})، حيث وضع حواجز صارمة لمنع استغلال العمال:
\begin{itemize}[leftmargin=*, itemsep=6pt]
    \item \textbf{قاعدة حظر شغل المناصب الدائمة (المادة \textenglish{L. 1251-5}):} تنص صراحة على أنه لا يجوز اللجوء إلى عقد الشغل المؤقت لشغل منصب يرتبط بالنشاط الطبيعي والدائم للمؤسسة المستعملة.
    \item \textbf{إعادة التكييف التلقائي بقوة القانون (المادة \textenglish{L. 1251-40}):} يحق للأجير المؤقت اللجوء إلى قضاء العمل (\textenglish{Conseil de prud'hommes}) للمطالبة بإعادة تكييف عقده إلى عقد غير محدد المدة (\textenglish{CDI}) مباشرة في مواجهة المقاولة المستعملة عند خرق قواعد اللجوء أو المدد، مع استحقاقه تعويضاً إجبارياً عن إعادة التكييف (\textenglish{Indemnité de requalification}) لا يقل عن أجر شهر كامل.
    \item \textbf{التجريم الجنائي للإعارة غير المشروعة للعمالة والاتجار بها:} يُعاقب القانون الفرنسي على «جريمة الاتجار باليد العاملة» (\textenglish{Délit de marchandage - Art. L. 8231-1}) و«الإعارة غير المشروعة لليد العاملة» (\textenglish{Prêt illicite de main-d'œuvre - Art. L. 8241-1}) بعقوبات حبسية تصل إلى \textenglish{3} سنوات وغرامات مالية ضخمة تصل إلى \textenglish{45.000} يورو للأشخاص الذاتيين وتتضاعف \textenglish{5} مرات للأشخاص الاعتباريين.
    \item \textbf{تعويض نهاية المهمة (\textenglish{Indemnité de fin de mission}):} إلزام المشغل بصرف تعويض نسبته \textenglish{10\%} من مجموع الأجور الإجمالية التي تقاضاها الأجير المؤقت كتعويض عن الهشاشة وعدم استقرار العمل.
    \item \textbf{مبدأ المساواة في الأجر:} إلزام المقاولة المستعملة بمنح الأجير المؤقت نفس الأجر والامتيازات التي يستفيد منها الأجير الدائم ذو المؤهلات المماثلة.
\end{itemize}

\subsection{التجربة الإسبانية: إصلاح 2021 ومحاصرة الهشاشة التعاقدية}
تعتبر التجربة الإسبانية نموذجاً حديثاً ملهماً في كبح جماح العمل المؤقت؛ فبعد عقود من انتشار العقود الهشة، أقرت إسبانيا تعديلاً جوهرياً بموجب المرسوم الملكي بقانون رقم \textenglish{32/2021} (\textenglish{Reforma Laboral 2021}):
\begin{itemize}[leftmargin=*, itemsep=5pt]
    \item افتراض أن جميع عقود الشغل مبرمة لأجل غير مسمى (\textenglish{Contrato indefinido}) كأصل مطلق، وتقييد العقود المؤقتة في أضيق الحدود (الظروف الإنتاجية الطارئة فقط وبسقف زمني صارم).
    \item فرض المساواة المطلقة في الأجور، وظروف الصحة والسلامة، والترقيات والتدريب بين أجراء شركات التشغيل المؤقت (\textenglish{Empresas de Trabajo Temporal - ETT}) والأجراء الدائمين في المقاولة المستعملة.
    \item تشديد العقوبات الإدارية والمالية على مقاولات المناولة التي يثبت تحايلها، وربط المسؤولية التضامنية بجميع الالتزامات المهنية.
\end{itemize}

\subsection{التجربة الألمانية: نموذج المشاركة والحوار الاجتماعي}
يُقدم النظام القانوني الألماني نموذجاً متميزاً في تنظيم التشغيل المؤقت يقوم على ركيزتين أساسيتين: المشاركة المؤسساتية للعمال والحوار الاجتماعي المنتظم. فقد أخضع المشرع الألماني نشاط مقاولات التشغيل المؤقت لقانون تنظيم إعارة العمال (\textenglish{Arbeitnehmerüberlassungsgesetz - AÜG}) الذي خضع لإصلاح جذري سنة \textenglish{2017} تضمن مقتضيات حمائية متقدمة:

\begin{itemize}[leftmargin=*, itemsep=6pt]
    \item \textbf{مبدأ المساواة في المعاملة من اليوم الأول (\textenglish{Equal Pay / Equal Treatment}):} أقر الإصلاح حق الأجير المؤقت في التمتع بنفس الأجر وظروف العمل المطبقة على الأجراء الدائمين في المقاولة المستعملة ابتداءً من أول يوم عمل، مع إمكانية الاستثناء بموجب اتفاقيات جماعية قطاعية لمدة أقصاها \textenglish{9} أشهر، على أن يتم التدرج في المساواة الأجرية بعد ذلك وصولاً إلى المساواة الكاملة بعد \textenglish{15} شهراً كحد أقصى.
    \item \textbf{سقف زمني أقصى للمهمة (\textenglish{18} شهراً):} حدد القانون مدة قصوى لوضع الأجير المؤقت رهن إشارة نفس المقاولة المستعملة بـ\textenglish{18} شهراً متواصلة، وعند تجاوز هذه المدة يُعاد تكييف العقد تلقائياً إلى عقد غير محدد المدة يربط الأجير بالمقاولة المستعملة مباشرة.
    \item \textbf{دور مجالس العمل (\textenglish{Betriebsrat}):} يتمتع مجلس العمل (مجلس المقاولة) بصلاحيات واسعة في مراقبة اللجوء إلى التشغيل المؤقت داخل المؤسسة، إذ يحق له الاطلاع على جميع عقود الوضع رهن الإشارة، والاعتراض على أي تشغيل مؤقت يعتبره مخالفاً للقانون أو ماساً بحقوق الأجراء الدائمين، مما يُشكل رقابة داخلية فعالة ومباشرة تُكمل الرقابة القضائية اللاحقة.
    \item \textbf{الاتفاقيات الجماعية القطاعية (\textenglish{Tarifverträge}):} يلعب التفاوض الجماعي على مستوى القطاعات الصناعية دوراً محورياً في تحسين أوضاع العمال المؤقتين بألمانيا؛ إذ تُبرم اتفاقيات جماعية بين نقابات العمال واتحادات أصحاب المقاولات تتضمن شروطاً أكثر ملاءمة من الحد الأدنى القانوني، لا سيما في مجالات العلاوات الخاصة بالقطاعات الصناعية والتكوين المهني المستمر وضمانات التثبيت بعد انتهاء المهمة.
\end{itemize}

ويُمثل النموذج الألماني مقاربة تشاركية تجمع بين المرونة التي يحتاجها الاقتصاد والحماية الفعلية للعمال، وهو ما يمكن أن يُشكل مصدر إلهام للمشرع المغربي في تطوير منظومته القانونية المتعلقة بالتشغيل المؤقت والمناولة.

\subsection{تجارب عربية مقارنة (تونس ومصر)}
\begin{itemize}[leftmargin=*, itemsep=6pt]
    \item \textbf{التجربة التونسية:} أفردت «مجلة الشغل التونسية» في المواد \textenglish{28} وما يليها تنظيماً حذراً لمقاولي اليد العاملة، حيث منعت المتاجرة بالأجراء واشترطت موافقة السلطة الإدارية. وتعيش تونس اليوم تحولاً تشريعياً ومجتمعياً نوعياً نحو المنع الشامل لظاهرة المناولة في الوظيفة العمومية والقطاع الخاص، واعتبارها شكلاً من أشكال العمل الهش الذي يتعارض مع مبدأ العدالة الاجتماعية وكرامة الشغالين.
    \item \textbf{التجربة المصرية:} نظّم قانون العمل المصري رقم \textenglish{12} لسنة \textenglish{2003} مسألة تشغيل العمال عبر متعهدي ووكالات التشغيل في الداخل والخارج؛ وحظر تشغيل العمال عن طريق وسطاء غير مرخص لهم، كما منع تحصيل أي مبالغ مالية من الأجراء مقابل إلحاقهم بالعمل، وألزم صاحب العمل الأصلي بالتضامن الكامل في الوفاء بالأجور والمستحقات العمالية عند الاستعانة بمقاولي الباطن.
\end{itemize}

%% ─────────────────────────────────────────────────────────
%%  رابعاً: المعايير الدولية للشغل
%% ─────────────────────────────────────────────────────────
\section{المعايير الدولية للشغل}

تُشكل معايير منظمة العمل الدولية (\textenglish{OIT}) الإطار المرجعي الأممي الذي يوجه السياسات الوطنية في تنظيم أسواق العمل:

\begin{enumerate}[leftmargin=*, itemsep=8pt]
    \item \textbf{الاتفاقية الدولية رقم 181 (لسنة 1997) بشأن وكالات الاستخدام الخاصة:}
          تعترف هذه الاتفاقية بدور وكالات التشغيل المؤقت في مرونة سوق الشغل، لكنها تُخضع نشاطها لمنظومة مشددة من الضمانات:
          \begin{itemize}[leftmargin=*, itemsep=3pt]
              \item منع وكالات التشغيل من اقتطاع أي أتعاب أو مصاريف مباشرة أو غير مباشرة من أجور العمال.
              \item ضمان الحق في المفاوضة الجماعية والحرية النقابية لعمال التشغيل المؤقت.
              \item إلزام الدول الأطراف بوضع آليات لتحديد مسؤوليات كل من الوكالة والمقاولة المستعملة في مجالات الأجور وحوادث الشغل.
          \end{itemize}
    \item \textbf{التوصية رقم 188 (لسنة 1997):} تدعو الدول إلى وضع آليات رقابية فعالة وعقوبات رادعة (بما فيها شطب وسحب تراخيص المقاولات المخالفة).
    \item \textbf{الاتفاقية رقم 94 بشأن بنود العمل في العقود العامة:} إلزام المقاولين وأصحاب الصفقات باحترام شروط الشغل اللائقة للأجراء المناولين.
    \item \textbf{مفهوم العمل اللائق (\textenglish{Decent Work}):} يؤكد إعلان منظمة العمل الدولية بشأن العدالة الاجتماعية من أجل عولمة منصفة (\textenglish{2008}) أن العمل ليس سلعة تجارية تباع وتشترى، وأن المرونة الاقتصادية تفقد مشروعيتها متى أفضت إلى حرمان الأجراء من الأمان الوظيفي والحماية الاجتماعية.
    \item \textbf{الاتفاقية رقم 122 (لسنة 1964) بشأن سياسة الاستخدام:}
          تُلزم هذه الاتفاقية الدول المصادقة عليها بانتهاج سياسة وطنية نشيطة ترمي إلى تحقيق \textbf{الاستخدام الكامل والمنتج والمختار بحرية}، بما يضمن لكل قادر على العمل وباحث عنه فرصة الحصول على شغل لائق ومستقر. وتُشكل هذه الاتفاقية أساساً معيارياً مهماً لمساءلة الدول التي تتوسع في اعتماد أشكال التشغيل الهش والمؤقت على حساب الاستخدام المستقر؛ إذ ينبغي ألا يُوظف التشغيل المؤقت كبديل منهجي عن خلق مناصب شغل دائمة ومنتجة، بل يجب أن يظل استثناءً محصوراً في الحاجيات الظرفية الحقيقية للمقاولات.
    \item \textbf{أجندة العمل اللائق \textenglish{2030} وأهداف التنمية المستدامة:}
          ربطت منظمة العمل الدولية بين مفهوم العمل اللائق وأهداف التنمية المستدامة (\textenglish{SDGs})، لا سيما \textbf{الهدف الثامن} (\textenglish{SDG 8}) المتعلق بـ«تعزيز النمو الاقتصادي المطرد والشامل والمستدام، والعمالة الكاملة والمنتجة، وتوفير العمل اللائق للجميع». ويستلزم هذا الهدف من الدول الأعضاء اتخاذ تدابير فعالة للقضاء على العمل الجبري والاتجار بالبشر وعمل الأطفال، وضمان بيئة عمل آمنة ومأمونة لجميع العمال بمن فيهم العمال المهاجرون والمؤقتون. كما تُلزم الغاية \textenglish{8.8} من هذا الهدف الدول بحماية حقوق العمل وتعزيز بيئات العمل المأمونة لجميع العمال، بما فيهم العمال الذين يعملون في ظروف هشة عبر عقود التشغيل المؤقت والمناولة.
\end{enumerate}

\subsection{تقييم مدى ملاءمة التشريع المغربي للمعايير الدولية}
يكشف استقراء المعايير الدولية المذكورة ومقارنتها بالمنظومة التشريعية المغربية عن وجود فجوات تستدعي التدخل التشريعي لسدها وتحقيق الملاءمة المطلوبة:

\begin{itemize}[leftmargin=*, itemsep=6pt]
    \item \textbf{على مستوى الضمانات الأجرية:} رغم نص المادة \textenglish{504} من مدونة الشغل على مبدأ المساواة في الأجر، إلا أن المشرع المغربي لم ينص على تعويض إلزامي عن الهشاشة (\textenglish{Précarité}) على غرار ما هو معمول به في فرنسا (نسبة \textenglish{10\%} من مجموع الأجور)، مما يُضعف الحماية المالية للأجير المؤقت ولا يُعوضه عن عدم استقرار وضعيته المهنية.
    \item \textbf{على مستوى الجزاءات:} تبقى الغرامات المنصوص عليها في مدونة الشغل المغربية ضعيفة وغير رادعة مقارنة بالعقوبات المقررة في التشريعات المتقدمة، كما يغيب التجريم الصريح لجريمة «الاتجار باليد العاملة» و«الإعارة غير المشروعة» بوصفهما جنحتين مستقلتين تستوجبان عقوبات حبسية.
    \item \textbf{على مستوى الرقابة المؤسساتية:} لم يُفرد المشرع المغربي دوراً مؤسساتياً لممثلي الأجراء أو المؤسسات النقابية في مراقبة اللجوء إلى التشغيل المؤقت والمناولة داخل المقاولات، على خلاف النموذج الألماني الذي يمنح مجالس العمل صلاحيات رقابية واسعة في هذا الشأن.
    \item \textbf{على مستوى المصادقة على الاتفاقيات:} لم يُصادق المغرب بعد على الاتفاقية الدولية رقم \textenglish{181} بشأن وكالات الاستخدام الخاصة، مما يُفقد الأجراء المؤقتين مرجعية دولية ملزمة يمكن الاحتجاج بها أمام القضاء الوطني لتعزيز حمايتهم.
\end{itemize}

%% ─── خلاصة الفصل ──────────────────────────────────────────
\section*{خلاصة الفصل}
\addcontentsline{toc}{section}{خلاصة الفصل}

يتبين من خلال استقراء الاجتهاد القضائي المغربي والتجارب المقارنة أن القضاء لم يقف مكتوف الأيدي أمام مخاطر استغلال الشغل المؤقت والمناولة كأدوات للتهرب من تكاليف قانون الشغل، بل سارع—مستعيناً بسلطته في تكييف العقود وبحث عناصر التبعية الواقعية—إلى ردع الممارسات الصورية وإعادة الأجراء إلى مظلة العقود الدائمة والمستقرة.

غير أن هذا المجهود القضائي، على عظمته، يظل علاجاً لاحقاً يقتضي من الأجير خوض غمار مساطر تقاضٍ طويلة ومكلفة في ظل تفاوت موازين القوى مع المقاولات الكبرى. وبالمقارنة مع التشريعات الدولية المتقدمة (كالنموذج الفرنسي والإسباني والألماني)، يتضح أن التشريع المغربي بحاجة ملحة إلى التطوير وسد الثغرات النصية، لا سيما بالنص الصريح على جزاء إعادة التكييف ورفع الغرامات وتجريم الإعارة غير المشروعة لليد العاملة كجنحة مستقلة.

ويستدعي الوضع الراهن تدخلاً تشريعياً عاجلاً يستلهم أفضل الممارسات الدولية، ولا سيما إقرار تعويض إلزامي عن الهشاشة لفائدة الأجراء المؤقتين، وتعزيز آليات الرقابة الإدارية والنقابية على مقاولات التشغيل المؤقت والمناولة، فضلاً عن المصادقة على الاتفاقيات الدولية ذات الصلة. إن الاجتهاد القضائي الفاعل، رغم أهميته البالغة في سد الفراغ التشريعي وحماية الحقوق الأساسية للعمال، لا يمكن أن يحل محل إرادة تشريعية واضحة وحازمة تُرسي قواعد آمرة لا تقبل التأويل ولا التحايل، وهو ما يشكل موضوع الدراسة والتحليل في الفصل العاشر الموالي.

\cleardoublepage

%% ─── الفصل العاشر: الاختلالات والتوصيات ──────────────────
%% ============================================================
%%  chapitre10.tex — الفصل العاشر: الاختلالات والتوصيات
%% ============================================================

\chapter{الاختلالات والتوصيات}

%% ─── مقدمة الفصل ──────────────────────────────────────────
\section*{مقدمة الفصل}
\addcontentsline{toc}{section}{مقدمة الفصل}

إذا كان المشرع المغربي قد سعى من خلال سن مقتضيات مدونة الشغل إلى مسايرة التحولات الاقتصادية عبر فتح الباب أمام أشكال التشغيل المرنة، ومنها الشغل المؤقت والمناولة، فإن الممارسة الميدانية خلال عقدين من تطبيق المدونة قد كشفت عن هوة سحيقة بين الفلسفة النظرية للنصوص القانونية والواقع الملموس داخل المقاولات.

فبدلاً من أن يشكل التشغيل المؤقت والمناولة آلية استثنائية وظرفية لمواكبة الطوارئ الإنتاجية أو الاستعانة بخبرات تخصصية دقيقة، تحولا في أحيان كثيرة إلى وسيلة ممنهجة للتحايل على استقرار العمل وتكريس الهشاشة الاجتماعية والمهنية. وتجسد الوضعية المهنية لمجموعتنا (المجموعة الثامنة)—التي يشتغل فيها عمال عبر شركة مناولة لسنوات طويلة داخل المصنع ذاته ويؤدون نفس مهام الأجراء الأصليين ولكن بأجور زهيدة وامتيازات منقوصة—النموذج الصارخ لهذه الاختلالات البنيوية.

بناءً على ذلك، يخصص هذا الفصل لتشخيص نقدي شامل لمكامن الخلل على المستويين التشريعي والعملي، تمهيداً لصياغة حزمة متكاملة من التوصيات الواقعية والموجهة لكافة الفاعلين المعنيين (المشرع، والإدارة، والمقاولات، والأجراء والنقابات)، للمساهمة في رد الاعتبار للأمن المهني وإعادة التوازن المفقود لعلاقات الشغل.

%% ─────────────────────────────────────────────────────────
\section{الاختلالات التشريعية في المنظومة القانونية المغربية}

أبان التحليل المعمق لنصوص مدونة الشغل (القانون رقم \textenglish{65.99}) عن جملة من الثغرات والقصور التشريعي التي تستغلها بعض المقاولات للالتفاف على القانون:

\subsection{غياب نص صريح يقرر إعادة التكييف التلقائي في مواجهة المقاولة المستعملة}
على خلاف التشريعات المقارنة (كالفقرة الأولى من المادة \textenglish{L. 1251-40} من قانون العمل الفرنسي)، خلت مدونة الشغل المغربية من نص تشريعي صريح وقاطع يقضي بـ«إعادة تكييف» عقد الشغل المؤقت إلى عقد شغل دائم وغير محدد المدة (\textenglish{CDI}) مباشرة في مواجهة المقاولة المستعملة عند خرق الشروط المنصوص عليها في المادة \textenglish{496} أو تجاوز السقوف الزمنية المحددة في المادة \textenglish{500}.

وقد ترتب عن هذا الفراغ التشريعي فتح الباب أمام تضارب الأحكام القضائية في محاكم الدرجة الأولى، وإلقاء عبء ثقيل ومكلف على عاتق الأجير المؤقت لإثبات التبعية القانونية وسوء نية المقاولة أمام القضاء لانتزاع حكم بإعادة التكييف.

\subsection{هزالة الجزاءات والغرامات المالية المقررة في المادة \textenglish{505}}
تنص المادة \textenglish{505} من مدونة الشغل على معاقبة كل مخالفة لأحكام مقاولات التشغيل المؤقت بغرامة تتراوح بين \textenglish{2000} و\textenglish{5000} درهم فقط. 

إن هذا المبلغ الهزيل لا يمثل أي قيمة رادعة أمام الشركات الصناعية والمقاولات الكبرى التي تحقق أرباحاً سنوية بملايين الدراهم؛ بل إن حسابات الربح والخسارة تدفع بعض المشغلين إلى تفضيل دفع هذه الغرامة الرمزية عند ضبط المخالفة، باعتبار ذلك أقل كلفة بكثير من تحمل تكاليف توظيف أجراء دائمين والوفاء بحقوق الأقدمية والتعويضات الاجتماعية.

\subsection{عدم تجريم «المناولة المقنعة» والاتجار باليد العاملة بنصوص جنائية خاصة}
يفتقر التشريع الجنائي والاجتماعي المغربي إلى تجريم صريح لما يُعرف في الفقه المقارن بـ«جريمة الاتجار باليد العاملة» (\textenglish{Délit de marchandage}) و«الإعارة غير المشروعة للعمالة» خارج إطار وكالات التشغيل المؤقت المعتمدة. هذا الفراغ يتيح لشركات المناولة غير المصنفة ومقاولات الواجهة تقديم خدمات «كراء الأيدي العاملة» دون أي مؤهلات فنية، متسترة وراء عقود مدنية وتجارية صورية.

\subsection{تعطيل آليات التتبع المؤسساتية (اللجنة الثلاثية للمادة \textenglish{496})}
نصت الفقرة الأخيرة من المادة \textenglish{496} من مدونة الشغل على إحداث «لجنة ثلاثية التركيب» تضم ممثلين عن الإدارة، ومنظمات المشغلين الأكثر تمثيلاً، والمنظمات النقابية للأجراء الأكثر تمثيلاً، تناط بها مهمة تتبع التطبيق السليم لمقتضيات التشغيل المؤقت. 

غير أن هذه اللجنة ظلت معطلة ومجمدة عملياً لسنوات طويلة، مما حرم المنظومة من آلية مؤسساتية استباقية للرصد والتقييم واقتراح التعديلات التشريعية الضرورية.

\subsection{قصور قواعد المسؤولية التضامنية في المقاولة من الباطن (المادة \textenglish{89})}
علقت المادة \textenglish{89} من مدونة الشغل قيام المسؤولية التضامنية للمقاول الأصلي عن أجور وتأمينات أجراء مقاول الباطن على شرطين مقيدين: أن يكون المقاول من الباطن \textbf{غير مقيد بالسجل التجاري} و\textbf{في حالة إعسار مثبت}. 

وقد أفرغ هذا التقييد المسؤولية التضامنية من فاعليتها؛ إذ يكفي أن يسجل مقاول الباطن نفسه في السجل التجاري كشركة ذات مسؤولية محدودة برأسمال رمزي، لكي يتحلل المقاول الأصلي تماماً من مسؤوليته التضامنية تجاه العمال، حتى لو اختفى مقاول الباطن أو أعلن إفلاسه لاحقاً!

%% ─────────────────────────────────────────────────────────
\section{الاختلالات الميدانية والعملية في واقع سوق الشغل}

على أرض الواقع، تفرز الممارسات غير القانونية آثاراً وخيمة تمس صلب الاستقرار الاجتماعي والكرامة الإنسانية للعمال، وهو ما تم رصده جلياً في وقائع نازلة الدراسة:

\subsection{التحايل الممنهج وسد الحاجيات الهيكلية الدائمة عبر المناولة والشغل المؤقت}
تبرز الوضعية المهنية للمجموعة الثامنة جوهر هذا الاختلال؛ حيث تلجأ الشركة الصناعية الكبرى إلى إبرام عقود مع شركة متخصصة لتوريد اليد العاملة، وتستخدم هؤلاء الأجراء لسنوات متواصلة داخل المصنع لأداء مهام الإنتاج والتجميع الأساسية التي تمثل صلب النشاط العادي والدائم للشركة. 

إن استمرار العمال لعدة سنوات في مناصبهم ذاتها ينسف كل طابع «مؤقت» أو «ظرفي» يدعيه المشغل، ويشكل تحايلاً موصوفاً على مبدأ استقرار العمل المنصوص عليه في المادة \textenglish{16} من المدونة.

\subsection{تكريس التمييز وازدواجية الأجور داخل الفضاء المهني الواحد}
يؤدي هذا الوضع إلى شرخ اجتماعي خطير داخل المصنع:
\begin{itemize}[leftmargin=*, itemsep=6pt]
    \item أجراء يمارسون نفس العمل المهني بنفس ساعات العمل، وبنفس درجة الجهد والإتقان.
    \item في المقابل، يتقاضى الأجراء الدائمون أجوراً مجزية ويستفيدون من منح الأقدمية، والتغطية الصحية التكميلية، والتعويضات الاجتماعية، والترقية.
    \item بينما يتقاضى عمال المناولة والتشغيل المؤقت الحد الأدنى للأجر فقط (أو دونه في بعض الحالات بعد اقتطاع عمولة الوكالة)، دون أي امتيازات تكميلية أو آفاق للتطور، مما يمثل خرقاً صارخاً للمبدأ الدستوري والدولي القاضي بـ\textbf{«الأجر المتساوي عن العمل المتساوي القيمة»}.
\end{itemize}

\subsection{هشاشة الوضع المهني وتكبيل الحقوق والحريات النقابية}
يعيش الأجير المؤقت في حالة دائمة من عدم الأمان النفسي والاجتماعي؛ فعقده المهدد دائماً بعدم التجديد يجعله عاجزاً عن التخطيط لحياته ومستقبل أسرته، ويحرمه من الولوج إلى القروض البنكية أو استئجار مسكن لائق. 

والأخطر من ذلك هو استخدام هاجس «إنهاء المهمة» وسيلة لترهيب الأجراء ومنعهم من ممارسة حقهم الدستوري في الانتماء النقابي أو المطالبة بتحسين ظروف العمل أو التصريح بحوادث الشغل.

\subsection{محدودية موارد جهاز تفتيش الشغل وضعف الرقابة الميدانية}
يصطدم مفتش الشغل—كما ظهر في تحليل الأدوار—بإكراهات بنيوية كبرى تحد من فاعلية تدخله:
\begin{itemize}[leftmargin=*, itemsep=5pt]
    \item قلة عدد مفتشي الشغل على الصعيد الوطني مقارنة بآلاف المنشآت والمصانع المتناثرة عبر التراب الوطني.
    \item غياب الوسائل اللوجستية والتقنية الكافية للتنقل وإجراء التحريات المعمقة في دفاتر الشركات.
    \item تعقيد المساطر القانونية وغياب سلطة توقيف المخالفات فورياً، مما يجعل دور المفتش في كثير من الأحيان مقتصراً على تحرير محاضر تطول إجراءات البت فيها أمام المحاكم.
\end{itemize}

\subsection{انتشار الشركات الوهمية (\textenglish{Sociétés écrans}) والتهرب الاجتماعي}
تلجأ بعض الأطراف إلى تأسيس شركات مناولة «صورية» عابرة لا تملك مقرات فعلية ولا أصولاً مادية، يتم من خلالها تشغيل مئات العمال لمدد معينة، ثم تتم تصفيتها أو التخلي عنها فجأة دون سداد واجبات الصندوق الوطني للضمان الاجتماعي (\textenglish{CNSS}) ومستحقات الأجراء، ليعاد تأسيس شركات جديدة بأسماء أخرى، في إفلات كامل من العقاب.

%% ─────────────────────────────────────────────────────────
\section{منظومة التوصيات المقترحة لتجاوز الاختلالات}

لمواجهة هذه الاختلالات واستعادة التوازن بين المرونة المشروعة والحماية الواجبة، نقترح منظومة متكاملة من التوصيات الإجرائية والعملية:

\subsection{توصيات موجهة للمشرع المغربي (إصلاحات تشريعية مستعجلة)}
\begin{tcolorbox}[colback=maincolor!5, colframe=maincolor, arc=5pt,
    title={\textbf{حزمة الإصلاحات التشريعية المقترحة}}]
\begin{enumerate}[leftmargin=*, itemsep=6pt]
    \item \textbf{التنصيص الصريح على إعادة التكييف بقوة القانون:} تعديل المادتين \textenglish{496} و\textenglish{500} بالنص القطعي على أن تشغيل أجير مؤقت في غير الحالات الاستثنائية الحصرية، أو تجاوز السقف الزمني المقرر، يجعل الأجير مرتبطاً تلقائياً بالمقاولة المستعملة بموجب عقد شغل غير محدد المدة (\textenglish{CDI})، وتسري أقدميته من تاريخ التحاقه الفعلي بالعمل.
    \item \textbf{تشديد العقوبات المالية والجنائية:} تعديل المادة \textenglish{505} برفع الغرامات لتتراوح بين \textenglish{30.000} و\textenglish{100.000} درهم، مع مضاعفتها بتعدد الأجراء المعنيين بالمخالفة، وإقرار عقوبة الحبس في حالة العود أو ثبوت التحايل المتعمد.
    \item \textbf{تجريم المناولة المقنعة والاتجار باليد العاملة:} إدراج نص قانوني يجرم «إعارة اليد العاملة غير المشروعة» بهدف الربح خارج إطار مقاولات التشغيل المؤقت المرخصة، ومصادرة الأرباح غير المشروعة الناتجة عنها.
    \item \textbf{تكريس مبدأ المساواة في الأجر والامتيازات:} إقرار نص تشريعي صريح يلزم المقاولة المستعملة بتطبيق مبدأ «الأجر المتساوي لنفس العمل»، وضمان استفادة الأجراء المؤقتين من كافة المكافآت والتعويضات العينية والمالية وخدمات النقل والمطعم الممنوحة للأجراء الدائمين.
    \item \textbf{إقرار تعويض نهاية المهمة (\textenglish{Prime de précarité}):} إلزام مقاولات التشغيل المؤقت بأداء تعويض مالي للأجير عند نهاية كل مهمة بنسبة لا تقل عن \textenglish{5\%} إلى \textenglish{10\%} من مجموع الأجر الخام المستحق، كتعويض عن طابع الهشاشة وعدم استقرار المنصب.
    \item \textbf{توسيع وتلقائية المسؤولية التضامنية في المناولة:} تعديل المادة \textenglish{89} بإلغاء اشتراط إعسار مقاول الباطن أو عدم تقييده بالسجل التجاري، وجعل المقاول الأصلي متضامناً بقوة القانون في أداء الأجور والواجبات الاجتماعية منذ اليوم الأول للتشغيل.
\end{enumerate}
\end{tcolorbox}

\subsection{توصيات موجهة للإدارة وأجهزة الرقابة والتفتيش}
\begin{itemize}[leftmargin=*, itemsep=6pt]
    \item \textbf{دعم جهاز تفتيش الشغل:} الرفع العاجل من أعداد مفتشي الشغل عبر مباريات توظيف مخصصة، وتزويدهم بالسيارات والمعدات الرقمية المتطورة لتمكينهم من القيام بحملات مراقبة فجائية واسعة النطاق.
    \item \textbf{إحداث منصة رقمية وطنية مركزية:} إنشاء منصة إلكترونية تلزم مقاولات التشغيل المؤقت والمقاولات المستعملة بإيداع عقود وضع الأجراء وعقود المهام بشكل آني، مع ربط المنصة بقواعد بيانات الصندوق الوطني للضمان الاجتماعي وإدارة الضرائب لرصد أي تكرار أو تجاوز غير قانوني للمدد.
    \item \textbf{تفعيل اللجنة الثلاثية الوطنية:} عقد اجتماعات دورية للجنة الثلاثية المنصوص عليها في المادة \textenglish{496}، وإصدار تقرير سنوي وطني يرصد خروقات التشغيل المؤقت والمناولة وينشر للعموم تكريساً لمبدأ الشفافية.
    \item \textbf{تشديد شروط منح وسحب التراخيص:} ربط منح ترخيص ممارسة نشاط التشغيل المؤقت بتقديم كفالة مالية بنكية كبرى تضمن سداد أجور الأجراء ومستحقاتهم في حالة الإفلاس، مع السحب الفوري للترخيص من كل شركة يثبت تورطها في مناولة مقنعة أو صورية.
\end{itemize}

\subsection{توصيات موجهة للمقاولات وأرباب العمل}
\begin{itemize}[leftmargin=*, itemsep=6pt]
    \item \textbf{تبني المسؤولية الاجتماعية للمقاولة (\textenglish{RSE}):} الكف عن اعتبار اليد العاملة مجرد كلفة مالية يجب تقليصها بأي ثمن، وإدراك أن الاستثمار في الرأسمال البشري المستقر هو السبيل الأوحد لرفع الجودة والإنتاجية وضمان السلم الاجتماعي المستدام.
    \item \textbf{قصر الشغل المؤقت على وظيفته النبيلة:} الالتزام الحرفي بعدم تشغيل أجراء مؤقتين في مناصب قارة تمس صلب نشاط المؤسسة، والاستجابة لاحتياجات الإنتاج الهيكلية عبر إبرام عقود غير محددة المدة.
    \item \textbf{اعتماد ميثاق أخلاقي للمناولة:} صياغة دفاتر تحملات صارمة للمناولة تلزم مقاول الباطن باحترام معايير العمل اللائق، وأداء الأجور القانونية في مواعيدها، وإخضاع ذلك لافتحاص دوري دقيق من طرف أطر الشركة الأصلية.
\end{itemize}

\subsection{توصيات موجهة للأجراء والتنظيمات النقابية والمجتمع المدني}
\begin{itemize}[leftmargin=*, itemsep=6pt]
    \item \textbf{توسيع التأطير النقابي للعمال الهشين:} تأسيس لجان ونقابات قطاعية متخصصة في تأطير عمال التشغيل المؤقت والمناولة، ومساندتهم في فتح قنوات الحوار والمفاوضة الجماعية مع المقاولات المستعملة.
    \item \textbf{نشر الثقافة والوعي القانوني:} إطلاق حملات تحسيسية لتبصير الأجراء بحقوقهم المكفولة قانوناً، وتوعيتهم بأهمية الاحتفاظ بالوثائق الثبوتية (نسخ العقود، وبطاقات الدخول للمصنع، وبيانات الحساب البنكي، وأوامر العمل المكتوبة) لإثبات العلاقة الشغلية عند النزاع.
    \item \textbf{إحداث مراصد قانونية للمساعدة القضائية:} تنسيق العمل بين النقابات والجمعيات الحقوقية وهيئات المحامين لتوفير المساعدة القانونية والمؤازرة القضائية المجانية للأجراء المؤقتين الذين يتعرضون للطرد التعسفي.
\end{itemize}

%% ─── خلاصة الفصل ──────────────────────────────────────────
\section*{خلاصة الفصل}
\addcontentsline{toc}{section}{خلاصة الفصل}

تؤكد دراسة الاختلالات التشريعية والميدانية أن ظاهرتي الشغل المؤقت والمناولة في المغرب قد بلغتا منعطفاً حاسماً لم يعد معه السكوت أو الاكتفاء بالترقيعات الجزئية ممكناً؛ فالواقع يكشف أن النمطين خرجا في حالات عديدة عن أهدافهما الاقتصادية المشروعة ليتحولا إلى أداة لإنتاج الهشاشة وتعميق الفوارق داخل فضاء العمل نفسه.

إن معالجة هذه المعضلة—كما جسدتها نازلة المصنع قيد التحليل—تتطلب شجاعة تشريعية تقضي على الفراغ القانوني وتسن جزاءات رادعة، مقترنة بيقظة إدارية ورقابة ميدانية حازمة، ووعي مجتمعي مسؤول من طرف أرباب العمل يؤمن بأن حماية كرامة العامل وحقه في عمل لائق ومستقر هما حجر الزاوية لكل إقلاع اقتصادي حقيقي.

\cleardoublepage

%% ─── الخاتمة والملاحق والمراجع ───────────────────────────
%% ============================================================
%%  conclusion.tex — الخاتمة وملاحق ولائحة المراجع
%% ============================================================

%% ════════════════════════════════════════════════════════════
%%  الخاتمة
%% ════════════════════════════════════════════════════════════
\chapter*{الخاتمة}
\addcontentsline{toc}{chapter}{الخاتمة}
\markboth{الخاتمة}{الخاتمة}

خلص هذا التقرير، بعد استعراض شامل لمختلف جوانب الشغل المؤقت والمناولة في ضوء
مدونة الشغل المغربية، إلى جملة من النتائج الجوهرية:

\begin{enumerate}[leftmargin=*, itemsep=10pt]
    \item يُتيح النظام القانوني المغربي اللجوءَ إلى الشغل المؤقت في حالات محددة
          وحصرية، غير أن التطبيق الفعلي يكشف عن فجوات بين النص والممارسة.
    \item تُكرّس مدونة الشغل مبدأ التضامن بين شركة التشغيل المؤقت والمقاولة
          المستعملة، بوصفه ضمانة أساسية لحماية الأجراء.
    \item يظل دور مفتش الشغل محورياً، غير أنه مقيَّد بمحدودية الموارد البشرية
          والمادية.
    \item تُسهم التجارب المقارنة في إثراء النموذج المغربي وتجويده.
    \item تبقى التوصيات الموجّهة رهينةَ الإرادة السياسية والوعي الاجتماعي
          لدى مختلف الأطراف المعنية.
\end{enumerate}

\vspace{1cm}

إن تحقيق التوازن بين متطلبات المرونة الاقتصادية التي تنشدها المقاولات وضرورة
الحماية الاجتماعية التي تستحقها الطبقة العاملة، لَيُعدّ تحدياً حقيقياً يستوجب
حواراً اجتماعياً جاداً وإصلاحات تشريعية مستمرة.

%% ════════════════════════════════════════════════════════════
%%  الملاحق
%% ════════════════════════════════════════════════════════════
\appendix

\chapter{الملحق الأول: نصوص قانونية مختارة}
\addcontentsline{toc}{chapter}{الملحق الأول: نصوص قانونية مختارة}

%% أدرج هنا نصوص المواد القانونية ذات الصلة

\section{المادة 495 من مدونة الشغل}
\textit{(أدرج نص المادة كاملاً)}

\section{المادة 496 من مدونة الشغل}
\textit{(أدرج نص المادة كاملاً)}

\section{المادة 86 من مدونة الشغل}
\textit{(أدرج نص المادة كاملاً)}

%% ─────────────────────────────────────────────────────────
\chapter{الملحق الثاني: نموذج عقد التشغيل المؤقت}
\addcontentsline{toc}{chapter}{الملحق الثاني: نموذج عقد التشغيل المؤقت}

\textit{(أدرج هنا نموذج العقد المرجعي المستخدم في الوضعية المدروسة)}

%% ════════════════════════════════════════════════════════════
%%  لائحة المراجع
%% ════════════════════════════════════════════════════════════
\chapter*{لائحة المراجع}
\addcontentsline{toc}{chapter}{لائحة المراجع}
\markboth{لائحة المراجع}{لائحة المراجع}

%% ─── أولاً: النصوص القانونية ──────────────────────────────
\section*{أولاً: النصوص القانونية}
\begin{itemize}[leftmargin=*, itemsep=6pt]
    \item القانون رقم \textenglish{65.99} المتعلق بمدونة الشغل، الجريدة الرسمية،
          \textenglish{2003}.
    \item الظهير الشريف رقم \textenglish{1.03.194} بتاريخ \textenglish{14} رجب
          \textenglish{1424} (\textenglish{11} غشت \textenglish{2003}).
    \item المراسيم التطبيقية ذات الصلة.
\end{itemize}

%% ─── ثانياً: الكتب والمؤلفات ─────────────────────────────
\section*{ثانياً: الكتب والمؤلفات}
\begin{itemize}[leftmargin=*, itemsep=6pt]
    \item المؤلف، \textit{عنوان الكتاب}، دار النشر، المدينة، السنة.
    \item المؤلف، \textit{عنوان الكتاب}، دار النشر، المدينة، السنة.
\end{itemize}

%% ─── ثالثاً: المقالات والدراسات ──────────────────────────
\section*{ثالثاً: المقالات والدراسات}
\begin{itemize}[leftmargin=*, itemsep=6pt]
    \item المؤلف، «عنوان المقال»، \textit{اسم المجلة}، العدد، السنة، الصفحات.
    \item المؤلف، «عنوان المقال»، \textit{اسم المجلة}، العدد، السنة، الصفحات.
\end{itemize}

%% ─── رابعاً: المراجع الإلكترونية ─────────────────────────
\section*{رابعاً: المراجع الإلكترونية}
\begin{itemize}[leftmargin=*, itemsep=6pt]
    \item الموقع الرسمي لوزارة الإدماج الاقتصادي والمقاولة الصغيرة والتشغيل
          والكفاءات: \url{www.emploi.gov.ma}
    \item موقع منظمة العمل الدولية: \url{www.ilo.org}
\end{itemize}



\end{document}

